\parte{Series y análisis periódico}{Los promedios de Cesàro preparan la reconstrucción de funciones periódicas. La medida normalizada del círculo mantiene visibles las constantes y permite pasar de coeficientes a convergencia puntual y en norma. Se utilizan integración de Lebesgue y geometría de Hilbert de los volúmenes anteriores.}
\chapter{Promedios y sumación de Cesàro}\label{cap:v4-01}
Una serie de Fourier puede conservar información sobre una función aunque sus sumas parciales no converjan. Antes de trabajar con funciones, estudiaremos una operación elemental que reduce las oscilaciones: tomar promedios. Conviene distinguir siempre los términos de una serie, sus sumas parciales y los promedios de esas sumas.

\section{Regularidad del promedio}\label{lectura:v4-01}
\index{Cesaro@Cesàro!sumacion@sumación}\index{serie!sumable por Cesaro@sumable por Cesàro}
\begin{proposicion}
Si $x_n\to x$ en un espacio normado $V$, entonces $n^{-1}\sum_{k=1}^nx_k\to x$.
\end{proposicion}
\begin{proof}
Fijemos $\eps>0$. Existe $N$ tal que $\norm{x_k-x}<\eps/2$ para $k>N$. Para $n>N$,
\[
 \left\|\frac1n\sum_{k=1}^nx_k-x\right\|
 \leq\frac1n\sum_{k=1}^N\norm{x_k-x}+\frac{n-N}{n}\frac\eps2.
\]
La primera suma tiene un número fijo de términos; su cociente entre $n$ es menor que $\eps/2$ cuando $n$ es suficientemente grande. No se ha usado completitud: el límite ya pertenece a $V$.
\end{proof}
\begin{definicion}
Para una serie $\sum_{n\geq1}v_n$ escribimos
\[
 s_n=\sum_{k=1}^nv_k,\qquad \sigma_n=\frac1n\sum_{k=1}^ns_k.
\]
La serie es \emph{sumable por Cesàro}, o $(C,1)$-sumable, con suma $v$, si $\sigma_n\to v$ en norma.
\end{definicion}
\begin{proposicion}
La sumación de Cesàro es lineal y consistente con la suma ordinaria. Además, si $\sigma_n$ converge, entonces $s_n/n\to0$.
\end{proposicion}
\begin{proof}
Para las series $\sum u_n$ y $\sum v_n$, las sumas parciales y sus promedios se combinan con los mismos escalares; la continuidad de suma y multiplicación escalar prueba linealidad. La proposición anterior, aplicada a $(s_n)$, prueba consistencia. Finalmente,
\[
 \frac{s_n}{n}=\sigma_n-\frac{n-1}{n}\sigma_{n-1}
 =(\sigma_n-\sigma_{n-1})+\frac{\sigma_{n-1}}n\quad(n\geq2).
\]
La primera diferencia tiende a cero y la segunda también, porque una sucesión convergente es acotada.
\end{proof}
\begin{ejemplo}
Si $z\in\C$, $|z|=1$ y $z\ne1$, entonces $s_n=z(1-z^n)/(1-z)$. En consecuencia,
\[
 \sigma_n=\frac z{1-z}-\frac{z^2(1-z^n)}{n(1-z)^2}\longrightarrow\frac z{1-z}.
\]
La serie $\sum z^n$ no converge ordinariamente, pues sus términos no tienden a cero. Esto incluye $z=e^{i}$, ejemplo de la lectura original. Los promedios no cambian la suma de una serie convergente, pero sí amplían la clase de series a las que puede asignarse una suma.
\end{ejemplo}

\section{Una condición tauberiana}
\index{Hardy--Landau, teorema de}\index{teorema!tauberiano}
\begin{teorema}[Hardy--Landau]
Si $\sigma_n\to v$ y $\norm{v_n}\leq A/n$ para todo $n$ suficientemente grande, entonces $s_n\to v$.
\end{teorema}
\begin{estrategia}
El promedio global $\sigma_n$ puede ocultar oscilaciones. La cota sobre $v_n$ asegura que las sumas parciales cambian poco dentro de un intervalo corto de índices. Compararemos $s_m$ con el promedio de $s_{n+1},\ldots,s_m$, tomando $m/n$ cercano a uno.
\end{estrategia}
\begin{proof}
Para $m>n$ después del índice donde vale la hipótesis,
\[
 \frac1{m-n}\sum_{k=n+1}^ms_k=\frac{m\sigma_m-n\sigma_n}{m-n}.
\]
Además, $\norm{s_m-s_k}\leq\sum_{j=k+1}^m A/j\leq A(m-k)/n$. Por tanto
\[
 \norm{s_m-v}\leq
 \frac{m\norm{\sigma_m-v}+n\norm{\sigma_n-v}}{m-n}
 +A\frac{m-n}{n}.
\]
Fijemos $0<\delta<1$ y, para cada $m$ grande, tomemos $n=\lfloor m/(1+\delta)\rfloor$. Entonces $n\to\infty$, $m/n\to1+\delta$ y los cocientes $m/(m-n)$ y $n/(m-n)$ permanecen acotados. El primer sumando tiende a cero, de modo que $\limsup_m\norm{s_m-v}\leq A\delta$. Como esto vale para todo $\delta>0$, el límite superior es cero.
\end{proof}
Esta prueba corrige las identidades de sumas intermedias de la lectura: no aparece ningún término adicional $s_{n+1}$ en $\sum_{k=n+1}^ms_k=m\sigma_m-n\sigma_n$.

\section{Ejercicios sobre sumabilidad y productos}
\begin{ejercicio}[Lectura 1, ejercicio 1]\label{ej:v4-l01-1}
Si $\sum v_n$ tiene suma de Cesàro $v$, demostrar que $\sum_{n\geq1}v_{n+1}$ tiene suma de Cesàro $v-v_1$.
\end{ejercicio}


\begin{ejercicio}[Lectura 1, ejercicio 2: rectificación del criterio de Abel]\label{ej:v4-l01-2}
El enunciado original afirma que la monotonía de $(\alpha_n)$ basta para que $\sum\alpha_nv_n$ sea sumable por Cesàro cuando lo es $\sum v_n$. Mostrar que es falso. Probar la siguiente versión suficiente en un espacio de Banach $V$: si $(\alpha_n)$ converge y
\[
 \sum_{n\geq1} n|\alpha_n-\alpha_{n+1}|<\infty,
\]
entonces el producto término a término conserva la sumabilidad de Cesàro.
\end{ejercicio}


\begin{ejercicio}[Lectura 1, ejercicio 3]\label{ej:v4-l01-3}
Sean $s_n=\sum_{k=1}^nv_k$, $t_n=\sum_{k=1}^nkv_k$ y $\sigma_n=n^{-1}\sum_{k=1}^ns_k$. Probar: (i) $t_n=(n+1)s_n-n\sigma_n$; (ii) si la serie es Cesàro-sumable, converge ordinariamente si y solo si $t_n/n\to0$; (iii) la serie es Cesàro-sumable si y solo si $\sum t_n/[n(n+1)]$ converge en $V$.
\end{ejercicio}


\begin{ejercicio}[Lectura 1, ejercicio 4]\label{ej:v4-l01-4}
En dimensión finita, demostrar que una serie es Cesàro-sumable si y solo si lo son sus series coordenadas respecto de una base fija.
\end{ejercicio}


\begin{ejercicio}[Lectura 1, ejercicio 5: uniformidad necesaria]\label{ej:v4-l01-5}
El original propone caracterizar la sumabilidad en $\ell^\infty$ por sumabilidad de cada coordenada y acotación de las sumas obtenidas. Dar un contraejemplo y formular el criterio correcto.
\end{ejercicio}


\begin{ejercicio}[Lectura 1, ejercicio 6]\label{ej:v4-l01-6}
En un álgebra de Banach $A$, sea $\sum u_n$ absolutamente convergente y sea $w_n=\sum_{k=1}^nu_kv_{n+1-k}$. Demostrar que $\sum w_n$ converge absolutamente, ordinariamente o por Cesàro si $\sum v_n$ tiene, respectivamente, esa propiedad. En cada caso su suma es el producto de las sumas en el orden indicado.
\end{ejercicio}


\begin{ejercicio}[Lectura 1, ejercicio 7]\label{ej:v4-l01-7}
Si $\sum u_n=a$ y $\sum v_n=b$ convergen en un álgebra normada, probar que su producto de Cauchy tiene suma de Cesàro $ab$. Determinar si necesariamente converge ordinariamente.
\end{ejercicio}


\begin{ejercicio}[Lectura 1, ejercicio 8]\label{ej:v4-l01-8}
Calcular la suma de Cesàro de $\sum_{n\geq1}(-1)^{n+1}$ y mostrar que su producto de Cauchy consigo misma no es Cesàro-sumable.
\end{ejercicio}


\paragraph{Enlace con las series de Fourier.}
Los promedios de Cesàro convertirán el núcleo oscilante de Dirichlet en el núcleo no negativo de Fejér. La positividad permitirá estimar el error separando una vecindad pequeña del origen de su complemento; esa idea reaparecerá en las identidades aproximadas de grupos.

\chapter{Integración y espacios periódicos}\label{cap:v4-02}
La elección de una medida fija las constantes de las fórmulas de Fourier. En la recta usaremos la medida de Lebesgue $dx$; en el círculo usaremos la medida de probabilidad $dm(t)=dt/(2\pi)$. Primero precisaremos los cambios de variable y los espacios de funciones en los que se interpretarán las series.

\section{Medida de Lebesgue y aproximación}\label{lectura:v4-02}
\index{Lebesgue!medida}\index{funcion@función!escalonada}
La construcción de Lebesgue parte de la longitud de intervalos semiabiertos y extiende ese volumen mediante medida exterior y el criterio de Carathéodory. Recordamos del volumen II la convergencia monótona, la convergencia dominada, Tonelli--Fubini, Hölder, Radon--Nikodym y la completitud de los espacios $L^p$. Aquí desarrollamos las consecuencias específicas que necesitaremos.
\begin{proposicion}
La medida de Lebesgue es invariante por traslaciones y regular. Toda medida de Borel regular, finita en compactos e invariante por traslaciones sobre $\R$ es $c\,dx$, con $c\geq0$.
\end{proposicion}
\begin{proof}
Trasladar una cubierta numerable por intervalos no altera la suma de sus longitudes. Aplicando esto en las dos direcciones se obtiene invariancia de la medida exterior; el criterio de Carathéodory muestra que también se preservan los conjuntos medibles. La regularidad por abiertos se obtiene tomando cubiertas cuya longitud total se aproxima a la medida exterior. Para conjuntos de medida finita, se obtiene regularidad interior por compactos aplicando esa aproximación al complemento dentro de un intervalo y truncando después la recta. Para conjuntos de medida infinita se trunca primero; esto da regularidad interior en todos los conjuntos de Lebesgue medibles.

Sea ahora $\nu$ una medida con las propiedades enunciadas. Cada punto tiene la misma masa $a$. Si $a>0$, un compacto con infinitos puntos tendría medida infinita; así $a=0$. Definamos $c=\nu([0,1])$. La aditividad y la invariancia dan $\nu([0,q])=cq$ para $q\geq0$ racional. Aproximando una longitud real por racionales desde abajo y arriba, la continuidad de la medida da $\nu([a,b])=c(b-a)$. Las medidas coinciden en los intervalos y, por unicidad de extensión en cada intervalo acotado, en todos los borelianos; una unión numerable de intervalos acotados cubre $\R$. La afirmación se refiere a medidas borelianas; sus completaciones coinciden cuando $c>0$.
\end{proof}
\begin{teorema}\label{v4:densidad-escalonadas}
Para $1\leq p<\infty$, las combinaciones lineales finitas de indicadoras de intervalos acotados son densas en $L^p(\R)$. Si $f$ es Riemann-integrable en $[a,b]$, su extensión por cero es Lebesgue-integrable y las dos integrales coinciden.
\end{teorema}
\begin{proof}
Las funciones simples con soporte de medida finita son densas en $L^p$: se trunca primero $f$ en espacio y en altura y se aproxima su rango mediante una malla finita. Basta aproximar $\mathbf1_A$, con $|A|<\infty$. Dado $\eta>0$, la regularidad proporciona un abierto $O\supset A$ con $|O\setminus A|<\eta/2$ y $|O|<\infty$. Todo abierto de $\R$ es unión numerable de intervalos abiertos disjuntos: cada componente contiene un racional, de modo que hay a lo sumo una cantidad numerable. Si $O=\bigcup_{j\geq1}I_j$, elegimos $r$ con $|O\setminus\bigcup_{j\leq r}I_j|<\eta/2$. El conjunto $B=\bigcup_{j\leq r}I_j$ satisface $|A\triangle B|<\eta$, y
\[
 \norm{\mathbf1_A-\mathbf1_B}_p^p=|A\triangle B|<\eta.
\]
La desigualdad triangular permite combinar estas aproximaciones para una función simple. No se exige que una unión finita de intervalos cerrados contenga $A$: esa afirmación, usada en la lectura original, es falsa para conjuntos densos de medida pequeña.

Para una función real Riemann-integrable existen funciones escalonadas inferiores y superiores $u_n\leq f\leq v_n$ cuyas integrales difieren en menos de $1/n$. Sus integrales de Lebesgue son las sumas de Darboux correspondientes. El conjunto donde $\limsup u_n<\liminf v_n$ tiene medida cero tras escoger una sucesión con errores sumables; equivalemente puede aplicarse el criterio de Lebesgue para integrabilidad de Riemann, probado en el volumen II. Por completitud de la medida, $f$ es medible y, por comparación, su integral coincide con el valor común de las integrales de Darboux. Para funciones complejas se aplica el argumento a las partes real e imaginaria.
\end{proof}

\section{Intervalos y cambios de variable}
En un intervalo no degenerado $I$ escribimos
\[
 \norm f_{p,I}=\left(\frac1{|I|}\int_I|f|^p\,dx\right)^{1/p}
 \quad(1\leq p<\infty),\qquad
 \norm f_{\infty,I}=\operatorname*{ess\,sup}_{I}|f|.
\]
Las funciones medibles se consideran iguales cuando coinciden casi en todas partes. Con estas normas, $L^p(I)$ es Banach; multiplicar la medida por la constante $|I|^{-1}$ no modifica las demostraciones de completitud. Para $1<p<\infty$, el dual se identifica con $L^{p'}(I)$ mediante $g\mapsto[f\mapsto |I|^{-1}\int_I f\overline g]$, donde $p'=p/(p-1)$.
\begin{proposicion}
Si $I=[a,b]$, $J=[c,d]$ y $T(x)=c+\alpha(x-a)$, con $\alpha=|J|/|I|>0$, entonces
\[
 \int_J f(y)\,dy=\alpha\int_I f(T(x))\,dx.
\]
La composición $f\mapsto f\circ T$ es una isometría sobreyectiva entre los espacios $L^p$ normalizados, también para $p=\infty$.
\end{proposicion}
\begin{proof}
Extendamos la fórmula afín a toda la recta. Las preimágenes de intervalos tienen longitud dividida entre $\alpha$. Por aproximación mediante cubiertas, $|T^{-1}(E)|=\alpha^{-1}|E|$ para todo conjunto medible $E$; en particular las preimágenes de nulos son nulas. Un conjunto de Lebesgue medible difiere de un boreliano en un nulo, por lo que $f\circ T$ es medible. La identidad integral vale primero para indicadoras, después para funciones simples, luego para funciones no negativas por convergencia monótona y finalmente para funciones integrables complejas por descomposición en cuatro partes no negativas. Aplicándola a $|f|^p$ se obtiene la igualdad de normas; para $p=\infty$, las desigualdades esenciales se preservan porque los nulos se preservan en ambas direcciones. La composición con $T^{-1}$ es la inversa.
\end{proof}

\section{Ejercicios sobre el cambio de variable}
\begin{observacion}[Hipótesis y derivada correctas]
En los cuatro ejercicios de la lectura 2, $\varphi'\ne0$ asegura que $\varphi$ es estrictamente monótona, pero no que sea sobreyectiva sobre $\R$ ni que $\varphi'$ sea creciente. Escribiremos $J=\varphi(\R)$, un intervalo abierto, y demostraremos
\[
 \int_J f(y)\,dy=\int_\R f(\varphi(x))|\varphi'(x)|\,dx.
\]
Si se desea integrar sobre toda $\R$ en ambos lados, se añade $J=\R$; si se omite el valor absoluto, se exige $\varphi'>0$. Por ejemplo, $\varphi(x)=e^x$ muestra que la sobreyectividad no se deduce de $\varphi'\ne0$.
\end{observacion}
\begin{ejercicio}[Lectura 2, ejercicio 1]\label{ej:v4-l02-1}
Sea $\varphi\in C^1(\R)$ con $\varphi'(x)\ne0$ para todo $x$, y sea $f$ Lebesgue-medible en $J$. Probar que $f\circ\varphi$ es Lebesgue-medible.
\end{ejercicio}

\begin{ejercicio}[Lectura 2, ejercicio 2]\label{ej:v4-l02-2}
Definir $\mu(A)=\int_\R\mathbf1_A(\varphi(x))\,dx$ para conjuntos medibles $A\subseteq\R$. Probar que $\mu$ es una medida absolutamente continua respecto de Lebesgue.
\end{ejercicio}

\begin{ejercicio}[Lectura 2, ejercicio 3: derivada de la medida imagen]\label{ej:v4-l02-3}
Hallar la derivada de Radon--Nikodym de $\mu$. Corregir la afirmación original $d\mu/dx=\varphi'$.
\end{ejercicio}

\begin{ejercicio}[Lectura 2, ejercicio 4]\label{ej:v4-l02-4}
Completar la demostración del cambio de variable, con las hipótesis corregidas anteriores.
\end{ejercicio}


\section{El círculo como espacio de medida}\label{lectura:v4-03}
\index{circulo@círculo}\index{espacio!Lp periodico@Lp periódico}
Identificamos $\T$ con $\R/(2\pi\Z)$. Una función medible de $\T$ se representa por una función $2\pi$-periódica en $\R$, y
\[
 \int_\T f\,dm=\frac1{2\pi}\int_0^{2\pi}f(t)\,dt.
\]
La integral es independiente del intervalo de longitud $2\pi$ elegido: al partir un intervalo trasladado en el extremo del período y trasladar una de las partes, se obtiene la misma integral. Cualquier intervalo se reduce a ese caso quitando un múltiplo entero del período. La restricción al intervalo semiabierto $[0,2\pi)$ es una biyección entre funciones periódicas y funciones sobre ese intervalo; la extensión periódica de una función medible es medible porque se define por traslaciones sobre una partición numerable.
\begin{proposicion}[Transporte de seminormas]
Si $T:V\to W$ es lineal y sobreyectivo y $p$ es una seminorma en $W$, entonces $p_T(v)=p(Tv)$ es la única seminorma que hace a $T$ isométrica. Si $N_W=\{p=0\}$ y $N_V=T^{-1}(N_W)$, la aplicación inducida $V/N_V\to W/N_W$ es un isomorfismo isométrico. Una forma sesquilineal no negativa en $W$ se transporta mediante $(u,v)\mapsto\ip{Tu}{Tv}$.
\end{proposicion}
\begin{proof}
La homogeneidad y la desigualdad triangular se conservan al componer con una aplicación lineal. La condición de isometría exige exactamente $p_T(v)=p(Tv)$, lo que prueba unicidad. El mapa cociente está bien definido porque dos representantes difieren en $N_V$ exactamente cuando sus imágenes difieren en $N_W$; es inyectivo, sobreyectivo e isométrico por esas definiciones. La forma transportada tiene cuadrado de seminorma $p_T^2$ y es única por polarización. En el original debe leerse $T^{-1}(N_W)$, no $T^{-1}(W)$.
\end{proof}
En consecuencia, $L^p(\T)$ es Banach y es separable para $p<\infty$: las funciones escalonadas con extremos racionales y valores racionales complejos forman un conjunto denso numerable. Para $p=2$ es Hilbert con
\[
 \ip fg=\int_\T f\overline g\,dm.
\]
El producto interno es lineal en su primer argumento. La periodicidad permite pasar de igualdad casi en todas partes en un período a igualdad casi en todas partes en $\R$, tomando una unión numerable de conjuntos nulos trasladados.

\section{Ortogonalidad y coeficientes que desaparecen}
\index{Riemann--Lebesgue, lema de}\index{familia!total}
Una familia $A$ de un espacio normado es \emph{completa} si $\overline{\operatorname{span}A}=V$. Una familia de funcionales es \emph{total} si el único vector que todos anulan es cero. En Hilbert identificamos un vector $a$ con el funcional $x\mapsto\ip xa$.
\begin{proposicion}
En Hilbert, una familia de vectores es completa si y solo si es total.
\end{proposicion}
\begin{proof}
Si su espacio generado es denso y $x$ es ortogonal a sus elementos, por linealidad y continuidad es ortogonal a todo el espacio, en particular a sí mismo. Recíprocamente, sea $M=\overline{\operatorname{span}A}$. Si $M\ne H$, la proyección ortogonal de un vector $x\notin M$ deja un residuo no nulo $x-P_Mx\in M^\perp$, que anula a toda la familia. Esto contradice totalidad.
\end{proof}
\begin{teorema}[Riemann--Lebesgue en el círculo]\label{v4:rl-circulo}
Si $f\in L^1(\T)$, entonces $\widehat f(n)=\int_\T f(t)e^{-int}\,dm(t)\to0$ cuando $|n|\to\infty$. Además $|\widehat f(n)|\leq\norm f_1$.
\end{teorema}
\begin{proof}
Para un intervalo $I=(a,b)$ y $n\ne0$,
\[
 \int_I e^{-int}\,dt=\frac{e^{-inb}-e^{-ina}}{-in},
\]
cuyo módulo es a lo sumo $2/|n|$. Por linealidad, los coeficientes de toda función escalonada tienden a cero. Dada $f$, aproximémosla por una escalonada $g$ con $\norm{f-g}_1<\eps$. Como $|e^{-int}|=1$,
\[
 |\widehat f(n)|\leq\norm{f-g}_1+|\widehat g(n)|<\eps+|\widehat g(n)|.
\]
Tomando límite superior y luego $\eps\downarrow0$ se obtiene el resultado. La cota de la norma se deduce directamente de la desigualdad triangular integral. Esto también prueba el caso $L^2$ de la lectura, pues $\norm f_1\leq\norm f_2$ en un espacio de probabilidad.
\end{proof}
\begin{proposicion}
Las funciones $e_n(t)=e^{int}$, $n\in\Z$, forman una familia ortonormal en $L^2(\T)$.
\end{proposicion}
\begin{proof}
Su producto interno es $(2\pi)^{-1}\int_0^{2\pi}e^{i(n-k)t}\,dt$. Vale uno si $n=k$; si $n\ne k$, la primitiva $e^{i(n-k)t}/(i(n-k))$ toma el mismo valor en los extremos, así que la integral es cero. La completitud de esta familia se probará después de estudiar Fejér.
\end{proof}

\section{Ejercicios sobre familias totales}
\begin{ejercicio}[Lectura 3, ejercicio 1]\label{ej:v4-l03-1}
Sea $1<p<\infty$ y sea $A\subset L^p(X,\mu)$. Demostrar que $A$ es completa en $L^p$ si y solo si los funcionales $g\mapsto\int_X a\overline g\,d\mu$, $a\in A$, son totales en $L^{p'}$.
\end{ejercicio}

\begin{ejercicio}[Lectura 3, ejercicio 2: normalización real]\label{ej:v4-l03-2}
Establecer Riemann--Lebesgue y ortogonalidad para los senos y cosenos. Precisar su normalización ortonormal.
\end{ejercicio}


\paragraph{Hacia la reconstrucción.}
La ortogonalidad proporciona coeficientes y cotas; por sí sola no asegura que los coeficientes reconstruyan todas las funciones. El siguiente capítulo desarrolla los núcleos que harán posible esa reconstrucción y distinguirá cuidadosamente convergencia puntual, uniforme y en norma.

\chapter{Series de Fourier y núcleos de sumación}\label{cap:v4-03}
Una expansión ortogonal es primero una afirmación sobre vectores de un espacio de Hilbert. Evaluarla en un punto requiere argumentos adicionales. En este capítulo esos argumentos se expresan mediante núcleos: Dirichlet para las sumas parciales y Fejér para sus promedios.

\section{Coeficientes y núcleo de Dirichlet}\label{lectura:v4-04}
\index{Fourier!serie}\index{Dirichlet!nucleo@núcleo}
Para $f\in L^1(\T)$ definimos
\[
 c_n=\widehat f(n)=\int_\T f(t)e^{-int}\,dm(t),\quad
 S_Nf(x)=\sum_{|n|\leq N}c_ne^{inx}.
\]
La notación $\widehat f$ denota aquí la sucesión de coeficientes, no el valor de la serie: no se presupone convergencia. La mera medibilidad no garantiza que los coeficientes existan; la hipótesis $L^1$ sí lo hace. Para $f\in L^2$, Bessel da $\sum_n|c_n|^2\leq\norm f_2^2$, y $S_N$ es la proyección ortogonal sobre el espacio generado por $e_{-N},\ldots,e_N$.
\begin{proposicion}\label{v4:dirichlet}
Sea $D_N(t)=\sum_{|n|\leq N}e^{int}$. Entonces
\[
 S_Nf(x)=\int_\T f(x-t)D_N(t)\,dm(t),\qquad
 D_N(t)=\frac{\sin((N+\tfrac12)t)}{\sin(t/2)}\quad(t\notin2\pi\Z).
\]
El núcleo es real y par, $\int_\T D_N\,dm=1$ y $\int_0^\pi D_N(t)\,dm(t)=1/2$. Para $0<\delta<\pi$,
\[
 \int_{-\delta}^{\delta}D_N(t)\,dm(t)\longrightarrow1.
\]
\end{proposicion}
\begin{proof}
Sustituimos $c_n$ en la suma finita, intercambiamos suma e integral y cambiamos $t$ por $x-t$; la invariancia de la integral periódica justifica recuperar un período fijo. La suma geométrica proporciona
\[
 D_N(t)=e^{-iNt}\frac{1-e^{i(2N+1)t}}{1-e^{it}}
 =\frac{\sin((N+1/2)t)}{\sin(t/2)}.
\]
En los múltiplos de $2\pi$, su valor es $2N+1$, por la definición finita. Los términos con índices opuestos son conjugados y su suma es par. Al integrar los exponenciales solo sobrevive el término constante, y la paridad divide la integral en dos mitades iguales. Finalmente, sobre $[\delta,\pi]$ el cociente $1/\sin(t/2)$ es acotado; Riemann--Lebesgue, aplicado a esa función multiplicada por $e^{it/2}$, muestra que la integral de $D_N$ sobre ese intervalo tiende a cero. Por paridad ocurre lo mismo en $[-\pi,-\delta]$. Restar esas colas de la integral total prueba la última afirmación.
\end{proof}
Usamos núcleos sin el factor $1/(2\pi)$ porque integramos contra $dm$. Los núcleos de las lecturas 4--5, cuando se integran contra $dt$, deben contener ese factor; mezclar ambas convenciones alteraría todas las constantes.

\section{Convergencia en puntos regulares}
\begin{teorema}[Dirichlet para funciones suaves a trozos]\label{v4:dirichlet-puntual}
Sea $f$ periódica, acotada y $C^1$ a trozos: en un período hay una partición finita en cuyos intervalos abiertos es $C^1$, con límites laterales finitos de $f$ y $f'$. Entonces
\[
 S_Nf(x)\longrightarrow\frac{f(x+)+f(x-)}2
\]
para todo $x$. En los puntos de continuidad, el límite es $f(x)$.
\end{teorema}
\begin{proof}
La paridad del núcleo permite escribir
\[
 S_Nf(x)-\frac{f(x+)+f(x-)}2
 =\frac1{2\pi}\int_0^\pi
 \frac{f(x+t)-f(x+)+f(x-t)-f(x-)}{\sin(t/2)}
 \sin((N+\tfrac12)t)\,dt.
\]
Cada diferencia del numerador es $O(t)$ cerca de cero: en cada lado, el valor medio o la integral de una derivada acotada da esa cota. Como $\sin(t/2)/t\to1/2$, el cociente es acotado cerca de cero. En el resto del intervalo es integrable por acotación de $f$ y porque el denominador no se anula. Riemann--Lebesgue, absorbiendo el factor $e^{it/2}$ como antes, hace tender la integral a cero. Los valores asignados a $f$ en sus saltos no intervienen en los coeficientes.
\end{proof}
\begin{ejemplo}[La función $|x|$]
Extendamos periódicamente $f(x)=|x|$ desde $[-\pi,\pi]$. Por paridad,
\[
 c_0=\frac\pi2,\quad
 c_n=\frac1\pi\int_0^\pi x\cos(nx)\,dx
 =\frac{(-1)^n-1}{\pi n^2}\quad(n\ne0).
\]
La integración por partes usa la primitiva $\sin(nx)/n$ y después $\int\sin(nx)\,dx=-\cos(nx)/n$. Por el teorema anterior,
\[
 |x|=\frac\pi2-\frac4\pi\sum_{j\geq1}\frac{\cos((2j-1)x)}{(2j-1)^2}.
\]
Evaluar en $x=0$ da $\sum_{j\geq1}(2j-1)^{-2}=\pi^2/8$.
\end{ejemplo}

\section{Ejercicios sobre coeficientes trigonométricos}
\begin{ejercicio}[Lectura 4, ejercicio 1]\label{ej:v4-l04-1}
Para $f\in L^1(\T)$, expresar la serie compleja como $a_0/2+\sum_{n\geq1}(a_n\cos nx+b_n\sin nx)$.
\end{ejercicio}

\begin{ejercicio}[Lectura 4, ejercicio 2]\label{ej:v4-l04-2}
Si $f\in L^1(\T)$ es real, probar que sus coeficientes $a_n,b_n$ son reales.
\end{ejercicio}

\begin{ejercicio}[Lectura 4, ejercicio 3]\label{ej:v4-l04-3}
Hallar la serie de $x^2$ en $[-\pi,\pi]$, extendida periódicamente, y deducir $\sum_{n\geq1}n^{-2}=\pi^2/6$.
\end{ejercicio}

\begin{ejercicio}[Lectura 4, ejercicio 4]\label{ej:v4-l04-4}
Sea $f$ periódica, continua y $C^1$ a trozos. Si $\sum_{n\in\Z}|c_n|<\infty$, probar que su serie converge a $f$ absoluta y uniformemente.
\end{ejercicio}


\section{Derivación, integración y convergencia uniforme}\label{lectura:v4-05}
\index{Fourier!derivacion@derivación}\index{Fourier!integracion@integración}
\begin{proposicion}
Si $f$ es periódica, continua y $C^1$ a trozos, entonces $\widehat{f'}(n)=in\widehat f(n)$, y su serie de Fourier converge absoluta y uniformemente a $f$.
\end{proposicion}
\begin{proof}
Integrando por partes en cada intervalo de suavidad, los términos de frontera interiores se cancelan porque $f$ es continua, y los extremos del período se cancelan por periodicidad. Queda $\int_\T f'e^{-int}\,dm=in\int_\T fe^{-int}\,dm$. Como $f'$ es acotada a trozos, pertenece a $L^2$. Bessel y Cauchy--Schwarz implican
\[
 \sum_{n\ne0}|c_n|
 \leq\left(\sum_{n\ne0}n^{-2}\right)^{1/2}
 \left(\sum_{n\ne0}|\widehat{f'}(n)|^2\right)^{1/2}<\infty.
\]
El ejercicio~\ref{ej:v4-l04-4} termina la prueba.
\end{proof}
\begin{proposicion}[Interpretación de la derivada de la serie]
Si $f$ es continua y $C^2$ a trozos, la serie formalmente derivada tiene coeficientes $inc_n$ y converge en cada punto al promedio de los límites laterales de $f'$. Donde $f'$ es continua, ese límite es $f'(x)$. Si además $f\in C^1(\T)$ y $f'$ es $C^1$ a trozos, la serie puede derivarse término a término con convergencia uniforme de las derivadas.
\end{proposicion}
\begin{proof}
La proposición anterior identifica $inc_n$ con los coeficientes de $f'$. Aplicar Dirichlet a $f'$ prueba la primera afirmación. Bajo la hipótesis adicional de continuidad de $f'$, su serie converge absoluta y uniformemente por la proposición anterior. El teorema de derivación de series de funciones, usando además convergencia de la serie original en un punto, permite derivar término a término. Sin esa continuidad no se afirma convergencia uniforme de las derivadas. El original escribe $f(x)$ en la conclusión para la serie de $f'$; debe decir $f'(x)$.
\end{proof}
\begin{proposicion}
Si $f$ es periódica y continua a trozos, $F(x)=\int_0^xf(t)\,dt$ es periódica si y solo si $c_0=0$. En ese caso sus coeficientes son $C_n=c_n/(in)$ para $n\ne0$ y $C_0=\int_\T F\,dm$. En general $G(x)=F(x)-c_0x$ es periódica y tiene esos coeficientes no constantes.
\end{proposicion}
\begin{proof}
La diferencia $F(x+2\pi)-F(x)=\int_x^{x+2\pi}f=2\pi c_0$ es independiente de $x$. Además, $F$ es continua y $C^1$ a trozos con $F'=f$ fuera de los puntos de partición. Si $c_0=0$, la fórmula de derivación aplicada a $F$ da $c_n=inC_n$. Si $c_0\ne0$, aplicamos el mismo razonamiento a $G'=f-c_0$, cuya media es cero. La constante se obtiene integrando $F$ o $G$ en un período; no se divide entre $n=0$.
\end{proof}

\section{Fejér: positividad y aproximación}\label{v4:fejer-seccion}
\index{Fejer@Fejér!nucleo@núcleo}\index{identidad!aproximada}
\begin{definicion}
Para $N\geq1$ definimos
\[
 K_N=\frac1N\sum_{j=0}^{N-1}D_j,\qquad
 \sigma_Nf=\frac1N\sum_{j=0}^{N-1}S_jf=f*K_N,
\]
donde en el círculo $f*g(x)=\int_\T f(x-t)g(t)\,dm(t)$.
\end{definicion}
\begin{proposicion}\label{v4:fejer-kernel}
Se tiene
\[
 K_N(t)=\sum_{|n|<N}\left(1-\frac{|n|}N\right)e^{int}
 =\frac1N\left|\sum_{j=0}^{N-1}e^{ijt}\right|^2
 =\frac{\sin^2(Nt/2)}{N\sin^2(t/2)}.
\]
En particular $K_N\geq0$, $\int_\T K_N\,dm=1$ y, para $\delta\leq|t|\leq\pi$,
\[
 K_N(t)\leq\frac1{N\sin^2(\delta/2)}.
\]
En $t\in2\pi\Z$, $K_N(t)=N$.
\end{proposicion}
\begin{proof}
El término $e^{int}$ aparece en $D_j$ exactamente cuando $j\geq|n|$, esto es, $N-|n|$ veces. Al multiplicar la suma geométrica por su conjugada, la diferencia $j-k=n$ también aparece $N-|n|$ veces. Esto prueba las dos primeras igualdades; la última se obtiene de la fórmula geométrica. La positividad es un cuadrado; la integral es el coeficiente constante; la cota fuera del origen usa $|\sin(Nt/2)|\leq1$ y $\sin(|t|/2)\geq\sin(\delta/2)$.
\end{proof}
Esta elección corrige a la vez el extremo de la suma y los factores del núcleo de la lectura 5. Promediar $S_1,\ldots,S_N$ en lugar de $S_0,\ldots,S_{N-1}$ produce el mismo límite: la diferencia es $(S_N-S_0)/N$, cuyo supremo tiende a cero porque $c_n\to0$ y $N^{-1}\sum_{|n|\leq N}|c_n|\to0$.
\begin{teorema}[Fejér]\label{v4:fejer-uniforme}
Si $f$ es continua y periódica, $\sigma_Nf\to f$ uniformemente. Si $f$ es acotada y tiene límites laterales en $x$, entonces $\sigma_Nf(x)\to\tfrac12(f(x+)+f(x-))$. Por tanto los polinomios trigonométricos son uniformemente densos en $C(\T)$.
\end{teorema}
\begin{proof}
Para $f$ continua,
\[
 |\sigma_Nf(x)-f(x)|\leq\int_\T|f(x-t)-f(x)|K_N(t)\,dm(t).
\]
Dado $\eps>0$, la continuidad uniforme permite elegir $\delta$ para que la diferencia sea $<\eps$ cuando la distancia circular de $t$ al origen sea $<\delta$. Esa parte de la integral es a lo sumo $\eps$, pues el núcleo tiene masa uno. La parte restante es a lo sumo $2\norm f_\infty/[N\sin^2(\delta/2)]$, independientemente de $x$. Esto prueba convergencia uniforme.

En un punto con límites laterales separamos $(-\pi,0)$ y $(0,\pi)$. En cada mitad restamos el límite lateral correspondiente. La masa de cada mitad es $1/2$ por paridad. Cerca de cero la diferencia es pequeña por la definición de límite lateral; lejos de cero es acotada y la masa del núcleo tiende a cero. Así se obtiene el promedio de ambos límites. Finalmente, $\sigma_Nf$ es un polinomio trigonométrico por la fórmula finita del núcleo, y aproxima uniformemente a $f$.
\end{proof}
El argumento de identidades aproximadas periódicas sigue también la exposición de Maximenko~\cite{maximenkoPeriodicos}, con normalización e intervalos de integración explicitados. La lectura 5 termina con el título «Ejercicios», sin una lista numerada; sus demostraciones dejadas como ejercicio se han desarrollado en las proposiciones anteriores.

\paragraph{Dos límites distintos.}
La masa total uno no basta por sí sola para reconstruir una función: Dirichlet cambia de signo, mientras que Fejér es no negativo y concentra su masa cerca del origen. Esta diferencia explica por qué los promedios convergen uniformemente para toda función continua, aunque las sumas parciales ordinarias no tengan esa propiedad.

\chapter{Fourier en los espacios de Lebesgue}\label{cap:v4-04}
Ahora pasamos de funciones regulares a clases de funciones integrables. La convergencia en norma no determina el valor en un punto, pero permite reconstruir toda la clase de equivalencia. Los puntos de Lebesgue proporcionarán después la conexión precisa con la convergencia puntual de los promedios.

\section{Densidad suave y completitud ortogonal}\label{lectura:v4-06}
\index{funcion@función!suave de soporte compacto}\index{Fourier!base ortonormal}
\begin{lema}\label{v4:densidad-suave}
Para $1\leq p<\infty$, $C_c^\infty(\R)$ es denso en $L^p(\R)$ y las funciones periódicas $C^\infty$ son densas en $L^p(\T)$.
\end{lema}
\begin{proof}
Definamos $r(t)=e^{-1/t}$ si $t>0$ y $r(t)=0$ si $t\leq0$. Cada derivada en $t>0$ es $e^{-1/t}$ multiplicada por un polinomio en $1/t$. Para todo entero $m$, $s^me^{-s}\to0$ cuando $s\to\infty$: basta comparar la serie de $e^s$ con su término de grado $m+1$. Todas las derivadas se prolongan entonces con valor cero en $t=0$. Por inducción en el orden, $r\in C^\infty(\R)$.

La función $\eta(t)=r(t)/(r(t)+r(1-t))$ es suave, vale cero si $t\leq0$, vale uno si $t\geq1$ y está entre cero y uno. Para $a<b$ y $\delta>0$,
\[
 h_\delta(x)=\eta\left(\frac{x-a+\delta}{\delta}\right)
 \eta\left(\frac{b+\delta-x}{\delta}\right)
\]
vale uno en $[a,b]$, se anula fuera de $[a-\delta,b+\delta]$ y satisface
$\norm{h_\delta-\mathbf1_{[a,b]}}_p^p\leq2\delta$.
La densidad de funciones escalonadas del teorema~\ref{v4:densidad-escalonadas}, seguida de una aproximación para cada uno de sus finitos intervalos, prueba el resultado en $\R$.

En $[0,2\pi]$ se puede primero cortar una función integrable en las franjas $(0,\delta)$ y $(2\pi-\delta,2\pi)$: la integral de $|f|^p$ en esas franjas tiende a cero. Aproximamos el resto por funciones escalonadas apoyadas en un subintervalo interior y después por las funciones suaves anteriores, escogiendo sus soportes todavía dentro de $(0,2\pi)$. Sus extensiones periódicas son suaves, pues se anulan en una vecindad de los extremos. Esto prueba la densidad periódica y justifica la suavidad que la lectura original dejaba sin verificar.
\end{proof}
\begin{teorema}\label{v4:base-fourier}
La familia $(e_n)_{n\in\Z}$ es una base ortonormal de $L^2(\T)$. Para toda $f\in L^2(\T)$,
\[
 S_Nf\longrightarrow f\text{ en }L^2,
 \qquad\norm f_2^2=\sum_{n\in\Z}|\widehat f(n)|^2.
\]
La aplicación $f\mapsto(\widehat f(n))$ es un isomorfismo unitario sobre $\ell^2(\Z)$.
\end{teorema}
\begin{proof}
Fejér aproxima uniformemente cada función continua mediante polinomios trigonométricos, y convergencia uniforme implica convergencia en $L^2$ porque $m(\T)=1$ y $\norm h_2\leq\norm h_\infty$. Por el lema anterior, las funciones continuas son densas en $L^2$. Así, el espacio generado por los exponenciales es denso.

Para precisar la convergencia de las sumas parciales ordinarias, dado $\eps>0$ escojamos un polinomio $P$ con $\norm{f-P}_2<\eps$. Si $N$ contiene todas sus frecuencias, $S_NP=P$. Como una proyección ortogonal tiene norma a lo sumo uno,
\[
 \norm{f-S_Nf}_2\leq\norm{f-P}_2+\norm{S_N(P-f)}_2<2\eps.
\]
La ortogonalidad da $\norm{S_Nf}_2^2=\sum_{|n|\leq N}|\widehat f(n)|^2$; tomar límites prueba Parseval. Si $(a_n)\in\ell^2$, las sumas $\sum_{|n|\leq N}a_ne_n$ son de Cauchy por Pitágoras. La completitud de $L^2$ da un límite $f$, y la continuidad de cada producto interno implica $\widehat f(n)=a_n$. Así la isometría es sobreyectiva.
\end{proof}
\begin{corolario}
Los coeficientes de Fourier determinan una función de $L^2$ salvo igualdad casi en todas partes. Además, alguna subsucesión de sus sumas parciales converge a ella casi en todas partes.
\end{corolario}
\begin{proof}
Si todos los coeficientes son cero, Parseval da $\norm f_2=0$. Para la segunda afirmación elegimos $N_j$ de modo que $\norm{S_{N_j}f-f}_2^2\leq2^{-j}$. Tonelli da $\int\sum_j|S_{N_j}f-f|^2<\infty$; por tanto la serie numérica es finita fuera de un nulo y sus términos tienden a cero.
\end{proof}
\begin{observacion}[Panorama: el teorema de Carleson]
El resultado mucho más profundo de Carleson afirma que la sucesión completa $S_Nf$ converge casi en todas partes para $f\in L^2(\T)$. Como en la lectura 6, se menciona como un teorema externo y no se utiliza en las demostraciones de este volumen. La prueba de convergencia en norma y la extracción de una subsucesión anteriores no prueban ese teorema. La referencia original se identifica al final~\cite{carleson1966}.
\end{observacion}
\begin{ejemplo}
Para $f(x)=\pi-x$ en $(0,2\pi)$, $a_n=0$ y $b_n=2/n$: integrar $\int(\pi-x)\sin nx$ por partes deja el término de frontera $2\pi/n$. Así $f=2\sum_{n\geq1}\sin(nx)/n$ en $L^2$ y en cada punto interior por Dirichlet. Parseval da
\[
 \frac1{2\pi}\int_0^{2\pi}(\pi-x)^2dx=\frac{\pi^2}3
 =2\sum_{n\geq1}\frac1{n^2}.
\]
No es necesario invocar Carleson para este ejemplo suave a trozos.
\end{ejemplo}

\section{Ejercicios de Parseval}
\begin{ejercicio}[Lectura 6, ejercicio 1]\label{ej:v4-l06-1}
Probar que senos, cosenos y la constante son totales en $L^2(\T)$ y obtener la identidad real de Parseval.
\end{ejercicio}

\begin{ejercicio}[Lectura 6, ejercicio 2]\label{ej:v4-l06-2}
Si $c_n,d_n$ son los coeficientes de $f,g\in L^2(\T)$, demostrar $\ip fg=\sum_n c_n\overline{d_n}$.
\end{ejercicio}

\begin{ejercicio}[Lectura 6, ejercicio 3]\label{ej:v4-l06-3}
Para $f$ real, periódica y $C^1$, demostrar mediante los coeficientes que $f'\perp f$.
\end{ejercicio}

\begin{ejercicio}[Lectura 6, ejercicio 4]\label{ej:v4-l06-4}
Usar la serie de $x^2$ para calcular $\sum_{n\geq1}n^{-4}$.
\end{ejercicio}

\begin{ejercicio}[Lectura 6, ejercicio 5]\label{ej:v4-l06-5}
Calcular la serie de $f(x)=x(\pi-|x|)$ en $[-\pi,\pi]$ y la suma $\sum_{j\geq1}(2j-1)^{-6}$.
\end{ejercicio}

\begin{ejercicio}[Lectura 6, ejercicio 6]\label{ej:v4-l06-6}
Sea $\alpha\in\R\setminus\Z$ y extiéndase periódicamente $f(x)=\pi e^{i(\pi-x)\alpha}/\sin(\pi\alpha)$ desde $[0,2\pi)$. Calcular sus coeficientes y demostrar $\sum_{n\in\Z}(n+\alpha)^{-2}=\pi^2/\sin^2(\pi\alpha)$.
\end{ejercicio}

\begin{ejercicio}[Lectura 6, ejercicio 7: Wirtinger]\label{ej:v4-l06-7}
Si $f$ es periódica, continua, $C^1$ a trozos y tiene media cero, demostrar $\int|f|^2\leq\int|f'|^2$ y caracterizar la igualdad.
\end{ejercicio}


\section{Localización y criterio de Dini}\label{lectura:v4-07}
\index{Dini!criterio}\index{localizacion@localización!de Fourier}
Por Riemann--Lebesgue, la transformación $L^1(\T)\to c_0(\Z)$ que asigna coeficientes es lineal y continua, de norma uno: la cota ya se probó y la constante $f=1$ alcanza la norma. La desaparición de coeficientes por sí sola no caracteriza cuáles sucesiones son transformadas; se demostrará su falta de sobreyectividad en el capítulo siguiente.
\begin{teorema}[Localización y Dini]\label{v4:dini}
Sea $f\in L^1(\T)$, $a\in\C$, $x\in\R$ y $0<\delta<\pi$. Entonces $S_Nf(x)\to a$ si y solo si
\[
 \int_0^\delta\frac{f(x+t)+f(x-t)-2a}{t}\sin((N+\tfrac12)t)\,dt\longrightarrow0.
\]
En particular, basta que $[f(x+t)+f(x-t)-2a]/t$ sea integrable cerca de cero.
\end{teorema}
\begin{proof}
Pongamos $g(t)=f(x+t)+f(x-t)-2a$. La fórmula de Dirichlet y la paridad dan
\[
 S_Nf(x)-a=\frac1\pi\int_0^\pi\frac{g(t)}{2\sin(t/2)}\sin((N+\tfrac12)t)\,dt.
\]
En $[\delta,\pi]$, $g/(2\sin(t/2))$ es integrable, así que su integral oscilante tiende a cero. En $(0,\delta)$, la diferencia
\[
 \frac1{2\sin(t/2)}-\frac1t
\]
se prolonga continuamente por cero en el origen: el desarrollo $2\sin(t/2)=t+O(t^3)$ da un error $O(t)$. Multiplicada por $g\in L^1(0,\delta)$ produce otra función integrable, cuya integral oscilante tiende a cero. Queda exactamente la equivalencia anunciada. Bajo la condición de Dini se aplica Riemann--Lebesgue directamente al último cociente. El signo entre $f(x+t)$ y $f(x-t)$ es suma; la resta que aparece en el enunciado original es una errata.
\end{proof}

\section{Puntos de Lebesgue y promedios de Fourier}
\index{Lebesgue!punto de}\index{Fejer@Fejér!teorema de Fejer--Lebesgue@teorema de Fejér--Lebesgue}
\begin{definicion}
Un punto $x$ es punto de Lebesgue de una función localmente integrable $f$ si $f(x)$ es finito y
\[
 \frac1{2h}\int_{-h}^h|f(x-t)-f(x)|\,dt\longrightarrow0\quad(h\downarrow0).
\]
\end{definicion}
Las dos integrales unilaterales son no negativas y están acotadas por la integral bilateral; por ello la condición equivale a exigir ambas versiones unilaterales. Todo punto de continuidad es de Lebesgue. El teorema de diferenciación del volumen II asegura que casi todo punto de una función localmente integrable es de Lebesgue para un representante apropiado.
\begin{teorema}[Fejér--Lebesgue]\label{v4:fejer-lebesgue}
Si $f\in L^1(\T)$ y $x$ es un punto de Lebesgue, entonces $\sigma_Nf(x)\to f(x)$.
\end{teorema}
\begin{proof}
Sea $g(t)=|f(x+t)+f(x-t)-2f(x)|$ para $0<t<\pi$ y $A(h)=\int_0^hg(t)\,dt$. La condición de Lebesgue implica $A(h)=o(h)$. La paridad de $K_N$ da
\[
 |\sigma_Nf(x)-f(x)|\leq\frac1{2\pi}\int_0^\pi g(t)K_N(t)\,dt.
\]
La fórmula del núcleo y $\sin(t/2)\geq t/\pi$ proporcionan
$K_N(t)\leq\min\{N,\pi^2/(Nt^2)\}$.
Fijemos $\eps>0$ y elijamos $\delta$ con $A(h)\leq\eps h$ para $0<h\leq\delta$. Para $N>1/\delta$ dividimos la integral en $(0,1/N)$, $(1/N,\delta)$ y $(\delta,\pi)$. La primera parte, incluyendo el factor $1/(2\pi)$, es a lo sumo $NA(1/N)/(2\pi)\leq\eps/(2\pi)$. Para la segunda, integración por partes de una función absolutamente continua da
\[
 \int_{1/N}^\delta\frac{g(t)}{t^2}\,dt
 =\frac{A(\delta)}{\delta^2}-N^2A(1/N)+2\int_{1/N}^\delta\frac{A(t)}{t^3}\,dt
 \leq2\eps N.
\]
Su contribución es a lo sumo $\pi\eps$. En la tercera, $K_N\leq1/[N\sin^2(\delta/2)]$, así que la contribución tiende a cero porque $g\in L^1$. El límite superior del error es a lo sumo $(\pi+1/(2\pi))\eps$. Hacer $\eps\downarrow0$ termina la prueba.
\end{proof}
La lectura 7 no contiene ejercicios numerados después de su encabezado final. Sus pasos de localización, integrabilidad del error y estimaciones de Fejér se han desarrollado aquí. La convergencia casi en todas partes de los \emph{promedios} para $L^1$ procede de Fejér--Lebesgue; no debe confundirse con la convergencia de las sumas ordinarias de Carleson.

\paragraph{Qué reconstruyen los coeficientes.}
En $L^2$, Fourier conserva exactamente la norma y convierte el espacio en $\ell^2$. En $L^1$ sigue habiendo unicidad, pero ni las normas ni las condiciones de convergencia se trasladan de la misma forma. La convolución permitirá formular esa diferencia y extender la aproximación a todos los exponentes finitos.

\parte{Convolución y frecuencias discretas}{La convolución reúne traslaciones de una función. En el círculo y en grupos finitos se puede comparar directamente con multiplicación de coeficientes; en la recta intervienen integrabilidad, continuidad y aproximaciones de la identidad.}
\chapter{Convolución periódica y Fourier finita}\label{cap:v4-05}
La convolución combina una función con todas las traslaciones de otra. En el círculo sirve para suavizar y aproximar; en un grupo finito es una suma que se puede calcular sin cuestiones de convergencia. Ambos modelos muestran la misma identidad: Fourier convierte convolución en producto.

\section{El álgebra de convolución del círculo}\label{lectura:v4-08}
\index{convolucion@convolución!periodica@periódica}
Para $f,g\in L^1(\T)$ definimos, donde la integral sea absolutamente convergente,
\[
 (f*g)(x)=\int_\T f(t)g(x-t)\,dm(t).
\]
La lectura 8 utiliza $dt$ y cambia su norma $L^1$ por la no normalizada. Aquí mantenemos $dm=dt/(2\pi)$ en todas las normas y convoluciones periódicas; las dos convenciones se relacionan multiplicando la convolución por $2\pi$.
\begin{teorema}\label{v4:conv-T}
La convolución está definida casi en todas partes, pertenece a $L^1(\T)$ y satisface
\[
 \norm{f*g}_1\leq\norm f_1\norm g_1,\qquad
 \int_\T f*g\,dm=\left(\int_\T f\,dm\right)\left(\int_\T g\,dm\right).
\]
Además $\widehat{f*g}(n)=\widehat f(n)\widehat g(n)$.
\end{teorema}
\begin{proof}
Elijamos representantes borelianos de $f,g$, modificándolos en conjuntos nulos si hace falta. El integrando $(x,t)\mapsto f(t)g(x-t)$ es medible. Las modificaciones nulas no afectan al resultado porque las traslaciones preservan nulos y el cálculo de Tonelli siguiente controla las secciones excepcionales. Por invariancia de $dm$,
\[
 \int_\T\int_\T |f(t)g(x-t)|\,dm(x)dm(t)
 =\int_\T|f(t)|\norm g_1\,dm(t)=\norm f_1\norm g_1.
\]
Fubini da existencia casi en todas partes y medibilidad de la integral. La desigualdad triangular, seguida de integración en $x$, da la cota. Aplicar Fubini sin valores absolutos da la fórmula para la integral. Finalmente, el factor $e^{-inx}$ tiene módulo uno y no altera la integrabilidad absoluta; con $u=x-t$,
\[
 \widehat{f*g}(n)=\int_\T\int_\T f(t)g(u)e^{-in(t+u)}\,dm(u)dm(t)
 =\widehat f(n)\widehat g(n).
\]
No se identifica la $\sigma$-álgebra producto de dos completaciones con toda la completación del producto: el uso de representantes borelianos evita esa afirmación incorrecta de la lectura.
\end{proof}

\section{Identidades aproximadas y unicidad}
\begin{lema}\label{v4:traslacion-T}
Si $1\leq p<\infty$ y $f\in L^p(\T)$, entonces $\norm{f(\cdot-t)-f}_p\to0$ cuando $t\to0$. La aplicación de traslación es uniformemente continua como aplicación de $\T$ en $L^p(\T)$.
\end{lema}
\begin{proof}
Para una función continua la conclusión se deduce de su continuidad uniforme y $\norm h_p\leq\norm h_\infty$. Para $f$ general escojamos continua $g$ con $\norm{f-g}_p<\eps/3$. La invariancia de la norma por traslación da
\[
 \norm{f(\cdot-t)-f}_p\leq2\norm{f-g}_p+\norm{g(\cdot-t)-g}_p.
\]
El último término es pequeño para $t$ pequeño. La diferencia entre dos traslaciones solo depende de la diferencia de sus parámetros, lo que da uniformidad.
\end{proof}
\begin{definicion}
Una sucesión de Dirac en $\T$ es una sucesión $k_j\geq0$ de $L^1$ con integral uno y $\int_{|t|\geq\delta}k_j\,dm\to0$ para cada $0<\delta<\pi$. Es fuerte si cada $k_j$ es esencialmente acotada y $\sup_{\delta\leq|t|\leq\pi}k_j(t)\to0$.
\end{definicion}
\begin{ejemplo}
Los núcleos $K_N$ de Fejér forman una sucesión fuerte. Otro ejemplo es
\[
 k_j(t)=\frac{\cos^{2j}(t/2)}{\int_\T\cos^{2j}(u/2)\,dm(u)}.
\]
Para $|t|\geq\delta$, el numerador es a lo sumo $\cos^{2j}(\delta/2)$. Al integrar el denominador sólo sobre $[-\delta/2,\delta/2]$, obtenemos
\[
 k_j(t)\leq\frac1{m([-\delta/2,\delta/2])}
 \left(\frac{\cos(\delta/2)}{\cos(\delta/4)}\right)^{2j}.
\]
El cociente de cosenos es menor que uno, así que la cota tiende a cero. Positividad e integral uno son inmediatas; la cota uniforme de la cola implica concentración de masa.
\end{ejemplo}
\begin{teorema}\label{v4:dirac-T}
Si $(k_j)$ es una sucesión de Dirac, entonces $f*k_j\to f$ en $L^p(\T)$ para $1\leq p<\infty$. Para $f$ continua, la convergencia es uniforme.
\end{teorema}
\begin{proof}
Para $p=1$, Tonelli y la desigualdad triangular dan
\[
 \norm{f*k_j-f}_1\leq\int_\T k_j(t)\norm{f(\cdot-t)-f}_1\,dm(t).
\]
Para $1<p<\infty$, Jensen aplicado a la probabilidad $k_jdm$ da
\[
 \norm{f*k_j-f}_p^p\leq\int_\T k_j(t)\norm{f(\cdot-t)-f}_p^p\,dm(t).
\]
En $|t|<\delta$ la diferencia de normas es uniformemente pequeña por el lema anterior; fuera, su potencia está acotada por $(2\norm f_p)^p$ y la masa tiende a cero. Esto prueba convergencia en norma. Para $f$ continua se aplica el mismo argumento al supremo de $|f(x-t)-f(x)|$, pequeño por continuidad uniforme.
\end{proof}
\begin{corolario}\label{v4:unicidad-T}
Si $f\in L^1(\T)$ tiene todos sus coeficientes nulos, entonces $f=0$ casi en todas partes. El álgebra de convolución $L^1(\T)$ no tiene unidad.
\end{corolario}
\begin{proof}
La fórmula finita de Fejér muestra que $f*K_N=0$ cuando todos los coeficientes de $f$ son cero. El teorema anterior da $\norm f_1=0$. Si hubiera unidad $u$, convolucionar con $e_n$ daría $\widehat u(n)e_n=e_n$ para todo $n$, por lo que $\widehat u(n)=1$. Esto contradice Riemann--Lebesgue.
\end{proof}
\begin{proposicion}
La aplicación de coeficientes $L^1(\T)\to c_0(\Z)$ no es sobreyectiva.
\end{proposicion}
\begin{proof}
Si fuera sobreyectiva, sería un isomorfismo lineal continuo entre espacios de Banach por el teorema de la inversa acotada, pues acabamos de probar inyectividad. Existiría $C$ tal que $\norm f_1\leq C\sup_n|\widehat f(n)|$. Para $f=D_N$, el supremo de los coeficientes es uno, pero $\norm{D_N}_1\to\infty$. Verifiquemos esto último. Con $M=N+1/2$ y $\sin(t/2)\leq t/2$,
\[
 \norm{D_N}_1\geq\frac2\pi\int_0^1\frac{|\sin(Mt)|}{t}\,dt.
\]
En cada intervalo $[k\pi/M,(k+1)\pi/M]\subset[0,1]$, $k\geq1$, la integral es al menos
\[
 \frac{M}{(k+1)\pi}\int_{k\pi/M}^{(k+1)\pi/M}|\sin(Mt)|\,dt
 =\frac2{(k+1)\pi}.
\]
El número de esos intervalos crece linealmente con $N$ y la suma armónica diverge. Esto contradice la supuesta cota de la inversa y completa la afirmación anunciada en la lectura 7.
\end{proof}

\section{Ejercicios de convolución periódica}
\begin{ejercicio}[Lectura 8, ejercicio 1]\label{ej:v4-l08-1}
Demostrar que $L^1(\T)$, con convolución, es un álgebra de Banach conmutativa.
\end{ejercicio}
\begin{solucion}
La bilinealidad sigue de la linealidad integral y la cota submultiplicativa del teorema~\ref{v4:conv-T}. La completitud es la de $L^1$. Cambiar $t$ por $x-t$ prueba conmutatividad. Para asociatividad, la integral triple del valor absoluto está acotada por $\norm f_1\norm g_1\norm h_1$. Fubini permite entonces calcular casi en todas partes
\[
 ((f*g)*h)(x)=\int_\T\int_\T f(s)g(t-s)h(x-t)\,dm(s)dm(t)
 =\int_\T f(s)\int_\T g(u)h(x-s-u)\,dm(u)dm(s),
\]
con $u=t-s$. El último miembro es $f*(g*h)(x)$.
\end{solucion}
\begin{ejercicio}[Lectura 8, ejercicio 2: cota del operador]\label{ej:v4-l08-2}
Si $f\in L^1(\T)$, $g\in L^p(\T)$ y $1<p<\infty$, demostrar que $f*g\in L^p$ y $\norm{f*g}_p\leq\norm f_1\norm g_p$. Precisar la norma de $T_f:g\mapsto f*g$.
\end{ejercicio}
\begin{solucion}
Si $\norm f_1=0$ la conclusión es inmediata. En otro caso, Hölder para la medida finita $|f(t)|dm(t)$ da
\[
 |f*g(x)|^p\leq\norm f_1^{p-1}\int_\T |f(t)|\,|g(x-t)|^pdm(t).
\]
Tonelli e invariancia de la medida hacen integrable el lado derecho y, tras integrar, dan $\norm{f*g}_p^p\leq\norm f_1^p\norm g_p^p$. En particular la integral de convolución existe absolutamente casi en todas partes. Por tanto $\norm{T_f}\leq\norm f_1$.

La igualdad afirmada en el original es falsa en general. Para $p=2$ y $f(t)=1+e^{it}$, Fourier muestra que $T_f$ es la proyección ortogonal sobre $\operatorname{span}\{1,e^{it}\}$, de norma uno, mientras que
\[
 \norm f_1=\frac1{2\pi}\int_0^{2\pi}2|\cos(t/2)|\,dt=\frac4\pi>1.
\]
La cota es la conclusión válida para todo $p$.
\end{solucion}
\begin{ejercicio}[Lectura 8, ejercicio 3]\label{ej:v4-l08-3}
Probar que $f\mapsto T_f$ es un homomorfismo continuo de álgebras de $L^1(\T)$ en $\Bcal(L^p(\T))$.
\end{ejercicio}
\begin{solucion}
La linealidad es inmediata. Si $f,h\in L^1$ y $g\in L^p$, aproximamos $g$ en $L^p$ por funciones acotadas $g_j$ (truncación en altura). Para cada $g_j$, las integrales de valor absoluto en la identidad de asociatividad son finitas porque el círculo tiene medida uno y $g_j$ es acotada. Así $T_{f*h}g_j=T_fT_hg_j$. La cota del ejercicio anterior permite pasar al límite en $L^p$ y obtener $T_{f*h}=T_fT_h$. También $\norm{T_f-T_h}\leq\norm{f-h}_1$, lo que prueba continuidad.
\end{solucion}
\begin{ejercicio}[Lectura 8, ejercicio 4]\label{ej:v4-l08-4}
Si $1<p<\infty$, $(k_j)$ es una sucesión de Dirac y $g\in L^p(\T)$, probar $\norm{g-g*k_j}_p\to0$.
\end{ejercicio}
\begin{solucion}
Jensen para la medida de probabilidad $k_jdm$ produce
\[
\norm{g-g*k_j}_p^p\leq\int k_j(t)\norm{g-g(\cdot-t)}_p^pdm(t).
\]
Para $\eps>0$ escojamos $\delta$ con $\norm{g-g(\cdot-t)}_p<\eps$ si $|t|<\delta$. La integral es a lo sumo
\[
 \eps^p+(2\norm g_p)^p\int_{|t|\geq\delta}k_j(t)\,dm(t).
\]
El límite superior es $\eps^p$; hacer $\eps\downarrow0$ da la conclusión. Esta es la versión $L^p$ completa de la sugerencia original, con una misma letra $g$ en hipótesis y conclusión.
\end{solucion}

\section{Un laboratorio finito: el grupo cíclico}
\index{Fourier!transformada discreta}\index{grupo!ciclico finito@cíclico finito}\index{matriz!circulante}
El complemento siguiente desarrolla el modelo finito presentado por Maximenko~\cite{maximenkoConvolucionFinita}. Sea $G=\Z/q\Z$, representado por $0,\ldots,q-1$, con medida de conteo. Los índices de una función $a:G\to\C$ se entienden módulo $q$. Definimos
\[
 (a*b)(r)=\sum_{k=0}^{q-1}a(k)b(r-k),\qquad
 \widehat a(j)=\sum_{k=0}^{q-1}a(k)e^{2\pi ijk/q}.
\]
\begin{teorema}\label{v4:dft}
Se cumplen inversión, Plancherel y el teorema de convolución:
\[
 a(k)=\frac1q\sum_{j=0}^{q-1}\widehat a(j)e^{-2\pi ijk/q},\qquad
 \sum_k|a(k)|^2=\frac1q\sum_j|\widehat a(j)|^2,
 \qquad\widehat{a*b}=\widehat a\,\widehat b.
\]
\end{teorema}
\begin{proof}
La suma geométrica da $\sum_{j=0}^{q-1}e^{2\pi ijr/q}=q$ si $r=0$ en $G$, y cero en otro caso. Sustituir la definición de $\widehat a$ en la fórmula inversa deja únicamente el término con índice $k$. Expandiendo $|\widehat a(j)|^2$ y sumando en $j$, la misma ortogonalidad elimina los términos cruzados y da $q\sum_k|a(k)|^2$.

Para convolución cambiamos $r=k+s$; para cada $k$, la aplicación $s\mapsto k+s$ permuta exactamente todos los elementos de $G$. Por tanto
\[
 \widehat{a*b}(j)=\sum_k\sum_s a(k)b(s)e^{2\pi ij(k+s)/q}
 =\widehat a(j)\widehat b(j).
\]
Esta biyección de índices es toda la justificación del cambio de suma; no hay ningún paso analítico oculto. El operador unitario $\mathcal F_q=q^{-1/2}\widehat{\phantom a}$ satisface $\mathcal F_q(a*b)=\sqrt q(\mathcal F_qa)(\mathcal F_qb)$: la normalización unitaria introduce el factor $\sqrt q$.
\end{proof}
La matriz de $T_a:b\mapsto a*b$ tiene entrada $a(r-k)$; cada columna es una traslación de la anterior. Se llama circulante. Los vectores $k\mapsto q^{-1/2}e^{-2\pi ijk/q}$ forman una base ortonormal de eigenvectores de $T_a$, con eigenvalores $\widehat a(j)$: sustituirlos en la definición de convolución lo verifica directamente. Así $\norm{T_a}_{\ell^2\to\ell^2}=\max_j|\widehat a(j)|$.
\begin{ejercicio}[Complemento: cálculo en cuatro puntos]\label{ej:v4-ampl-dft4}
En $\Z/4\Z$, calcular la convolución de $a=(1,2,0,0)$ y $b=(1,0,1,0)$ mediante sumas y mediante Fourier.
\end{ejercicio}
\begin{solucion}
Como $b$ vale uno en $0$ y $2$, $(a*b)(r)=a(r)+a(r-2)$, por lo que $a*b=(1,2,1,2)$. Las transformadas son
\[
 \widehat a=(3,1+2i,-1,1-2i),\qquad \widehat b=(2,0,2,0).
\]
El producto es $(6,0,-2,0)$. Invirtiendo, $(a*b)(k)=\tfrac14(6-2e^{-\pi ik})$, que vale uno si $k$ es par y dos si es impar. El cálculo exhibe la diferencia entre convolución cíclica y producto de polinomios sin reducción: aquí se trabaja módulo $z^4-1$.
\end{solucion}

\paragraph{La misma operación en tres modelos.}
Las sumas finitas, las series en grupos discretos y las integrales periódicas comparten la misma identidad de convolución. Al pasar a la recta desaparece la medida total finita; las hipótesis de integrabilidad y las estimaciones de Young serán las que aseguren que las operaciones sigan teniendo sentido.

\chapter{Convolución en el espacio euclidiano}\label{cap:v4-06}
Las estimaciones de convolución no dependen esencialmente de trabajar en una sola variable. Desarrollaremos la lectura 9 en $\R^d$, de modo que el caso original $d=1$ y la ampliación multidimensional compartan una misma prueba. La medida es Lebesgue sin normalización.

\section{Existencia, asociatividad e integrales}\label{lectura:v4-09}
\index{convolucion@convolución!euclidiana}
Para funciones medibles definimos $(f*g)(x)=\int_{\R^d}f(y)g(x-y)\,dy$ cuando la integral converge absolutamente. Si converge, el cambio $y=x-z$ prueba $f*g(x)=g*f(x)$. En general la función se define casi en todas partes y representa una clase de equivalencia.
\begin{teorema}\label{v4:conv-Rn}
Si $f,g\in L^1(\R^d)$, la convolución pertenece a $L^1$ y
\[
 \norm{f*g}_1\leq\norm f_1\norm g_1,\qquad
 \int_{\R^d}f*g=\left(\int_{\R^d}f\right)\left(\int_{\R^d}g\right).
\]
La convolución es bilineal, asociativa y conmutativa; hace de $L^1(\R^d)$ un álgebra de Banach.
\end{teorema}
\begin{proof}
Podemos elegir representantes borelianos de las clases. La aplicación $(x,y)\mapsto(y,x-y)$ es un isomorfismo lineal de determinante de valor absoluto uno, por lo que preserva la medida y los nulos. El integrando es medible también para representantes de Lebesgue. Tonelli y la traslación $z=x-y$ dan
\[
 \int_{\R^d}\int_{\R^d}|f(y)g(x-y)|\,dy\,dx
 =\int_{\R^d}|f(y)|\int_{\R^d}|g(z)|\,dz\,dy
 =\norm f_1\norm g_1<\infty.
\]
Fubini asegura existencia casi en todas partes y medibilidad. La desigualdad triangular integral produce la cota de norma; Fubini sin valores absolutos da la identidad de integrales. Cambiar un representante en un nulo no modifica la clase de la convolución, por el mismo cálculo.

Para tres funciones de $L^1$, la integral triple de $|f(z)g(y-z)h(x-y)|$ sobre $(x,y,z)$ es $\norm f_1\norm g_1\norm h_1$. Casi todo $x$ permite, por tanto, intercambiar las integrales en $y,z$. El cambio $y=z+t$ da
\[
 ((f*g)*h)(x)=\int f(z)\int g(t)h(x-z-t)\,dt\,dz=(f*(g*h))(x).
\]
Más generalmente, la misma igualdad vale en un punto fijo $x$ si $(|f|*|g|)*|h|(x)<\infty$, con las integrales internas definidas en las secciones pertinentes. La bilinealidad es inmediata y la completitud ya es la de $L^1$. En la fórmula de integrales de la lectura 9 deben figurar integrales sobre $\R$, no sobre $[0,2\pi]$.
\end{proof}

\section{La desigualdad de Young}
\index{Young!desigualdad de convolucion@desigualdad de convolución}
\begin{lema}[Hölder para varios factores]
Si $p_j\in[1,\infty]$ y $\sum_{j=1}^m1/p_j=1$, entonces $\norm{\prod_j f_j}_1\leq\prod_j\norm{f_j}_{p_j}$.
\end{lema}
\begin{proof}
Los factores de exponente infinito se acotan por su supremo esencial. Para los restantes, después de dividir cada factor por su norma, la desigualdad de Young ponderada da $\prod_j|f_j|\leq\sum_j|f_j|^{p_j}/p_j$. Integrando, el lado derecho vale uno. La desigualdad ponderada se obtiene de la convexidad de la exponencial: $\exp(\sum\theta_j a_j)\leq\sum\theta_j e^{a_j}$, con $\theta_j=1/p_j$ y $a_j=p_j\log|f_j|$; los ceros se tratan por límite. Si alguna norma es cero, el producto se anula casi en todas partes.
\end{proof}
\begin{teorema}[Young]\label{v4:young-Rn}
Sean $p,q,r\in[1,\infty]$ con $1/p+1/q=1+1/r$. Si $f\in L^p(\R^d)$ y $g\in L^q(\R^d)$, entonces
\[
 \norm{f*g}_r\leq\norm f_p\norm g_q.
\]
La convolución existe absolutamente casi en todas partes; si $r=\infty$, puede definirse en todo punto usando representantes adecuados, mediante Hölder.
\end{teorema}
\begin{proof}
Si $r=\infty$, $q=p'$ y Hölder aplicado a $f(y)$ y $g(x-y)$ da la cota para cada $x$. Si $p=1$ y $r=q<\infty$, Hölder para la medida $|f(y)|dy$ da
\[
 |f*g(x)|^q\leq\norm f_1^{q-1}\int|f(y)||g(x-y)|^qdy.
\]
Integrar en $x$ y aplicar Tonelli da la conclusión, incluida existencia casi en todas partes. Para $q=1$ se intercambian los factores. El caso $p=q=r=1$ ya está probado.

Queda $1<p,q\leq r<\infty$. Escribimos
\[
 |f(y)g(x-y)|=
 (|f(y)|^p|g(x-y)|^q)^{1/r}
 |f(y)|^{1-p/r}|g(x-y)|^{1-q/r}.
\]
Los exponentes de Hölder son $r$, $(1/p-1/r)^{-1}$ y $(1/q-1/r)^{-1}$; sus recíprocos suman uno. Un factor con exponente de potencia cero se omite. Como $|f|^p,|g|^q\in L^1$, el teorema anterior asegura que su convolución es finita casi en todas partes. Hölder da en esos puntos
\[
 |f*g(x)|^r\leq\norm f_p^{r-p}\norm g_q^{r-q}(|f|^p*|g|^q)(x).
\]
Integrar y usar la fórmula de la integral de convolución produce $\norm{f*g}_r^r\leq\norm f_p^r\norm g_q^r$. Todos los cambios se han hecho primero sobre integrandos no negativos, por lo que Tonelli justifica también la integrabilidad que se concluye.
\end{proof}
En particular, $T_f:g\mapsto f*g$ es lineal y acotado en $L^p$, con $\norm{T_f}\leq\norm f_1$. La igualdad de norma afirmada en la lectura 9 no vale para todos los exponentes. Por ejemplo, para $p=2$, $f=\mathbf1_{[0,1]}-\mathbf1_{[1,2]}$ tiene norma $L^1$ igual a dos, pero Plancherel dará $\norm{T_f}=\sup|\widehat f|<2$: en cada frecuencia la desigualdad triangular es estricta, pues las fases no pueden ser constantes en ambos intervalos con signos opuestos; la transformada es continua y se anula en infinito, por lo que alcanza su máximo y este es estrictamente menor que dos.

\section{Continuidad y comportamiento en infinito}
\begin{teorema}\label{v4:conv-continuidad}
Si $1\leq p<\infty$, $f\in L^p(\R^d)$ y $g\in L^{p'}(\R^d)$, entonces $f*g$ es acotada y uniformemente continua. Si $1<p<\infty$, además pertenece a $C_0(\R^d)$.
\end{teorema}
\begin{proof}
Hölder da $|f*g(x)|\leq\norm f_p\norm g_{p'}$ para todo $x$. La continuidad de traslaciones en $L^p$, probada en el ejercicio~\ref{ej:v4-l09-1}, implica
\[
 |f*g(x+h)-f*g(x)|\leq\norm g_{p'}\norm{f(\cdot+h)-f}_p\longrightarrow0,
\]
uniformemente en $x$. Si $1<p<\infty$, aproximamos $f,g$ en sus respectivas normas por funciones continuas de soporte compacto $u,v$. Hölder da
\[
 \norm{f*g-u*v}_\infty\leq\norm{f-u}_p\norm g_{p'}+\norm u_p\norm{g-v}_{p'}.
\]
La convolución $u*v$ es continua y se anula fuera de la suma compacta de los soportes. Por ello $f*g$ es límite uniforme de funciones de $C_c$, y pertenece a $C_0$. Para $p=1$, la aproximación de $g\in L^\infty$ por funciones de soporte compacto no está garantizada; el ejemplo $g=1$ explica por qué la conclusión adicional puede fallar.
\end{proof}
\begin{ejemplo}
En $\R$, $\mathbf1_{[0,1]}*\mathbf1_{[0,1]}(x)$ es la longitud de $[0,1]\cap[x-1,x]$. Vale $x$ para $0\leq x\leq1$, $2-x$ para $1\leq x\leq2$, y cero fuera de $[0,2]$. La convolución de dos funciones discontinuas puede así ser continua. En $\R^d$, la convolución de las indicadoras de $[0,1]^d$ es el producto de estas funciones triangulares, por Tonelli.
\end{ejemplo}

\section{Ejercicios: traslaciones, soportes y aproximación}
\begin{ejercicio}[Lectura 9, ejercicio 1]\label{ej:v4-l09-1}
Para $1\leq p<\infty$ y $f\in L^p(\R^d)$, demostrar que $t\mapsto f(\cdot+t)$ y $t\mapsto f(\cdot-t)$ son uniformemente continuas con valores en $L^p$.
\end{ejercicio}

\begin{ejercicio}[Lectura 9, ejercicio 2]\label{ej:v4-l09-2}
Sean $f\in L^1(\R^d)$ y $g$ acotada, con $g(x)\to0$ cuando $|x|\to\infty$. Probar que $f*g(x)\to0$ y mostrar que no puede omitirse una condición de este tipo en un enunciado general.
\end{ejercicio}

\begin{ejercicio}[Lectura 9, ejercicio 3]\label{ej:v4-l09-3}
Si $f,g$ tienen soporte compacto y su convolución existe casi en todas partes, demostrar $\supp(f*g)\subseteq\supp f+\supp g$, entendiendo soporte esencial cuando corresponda.
\end{ejercicio}

\begin{ejercicio}[Lectura 9, ejercicio 4]\label{ej:v4-l09-4}
Sea $(k_j)\subset L^1(\R^d)$, con $k_j\geq0$, $\int k_j=1$ y $\int_{|t|\geq\delta}k_j(t)dt\to0$ para cada $\delta>0$. Probar $f*k_j\to f$ en $L^p$ para $1\leq p<\infty$.
\end{ejercicio}

\begin{ejercicio}[Lectura 9, ejercicio 5: continuidad en los puntos del compacto]\label{ej:v4-l09-5}
Sea $f$ un representante acotado de $L^\infty(\R^d)$ que es continuo, como función definida en $\R^d$, en cada punto de un compacto $K$. Probar que $f*k_j\to f$ uniformemente en $K$. Explicar por qué no basta continuidad de la restricción $f|_K$.
\end{ejercicio}


Las traslaciones y las identidades aproximadas han convertido integrales en límites en norma. Una identidad gaussiana permitirá ahora invertir Fourier mediante cálculos explícitos.

\parte{La transformada en el espacio euclidiano}{El núcleo gaussiano permite justificar inversión y Plancherel sin intercambios formales de integrales. La descripción de subespacios invariantes desemboca en funciones de banda limitada, núcleos reproductores y reconstrucción por muestras.}
\chapter{Transformada euclidiana e inversión}\label{cap:v4-07}
La transformada en la recta reemplaza las frecuencias enteras por un parámetro continuo. Conservaremos el signo positivo del curso original y mantendremos una notación distinta para la transformada unitaria que aparecerá en el capítulo siguiente. Las pruebas se desarrollan en $\R^d$ cuando no hay una dificultad adicional; los argumentos específicos de la recta se identifican expresamente.

\section{La transformada de funciones integrables}\label{lectura:v4-10}
\index{Fourier!transformada euclidiana}\index{Riemann--Lebesgue, lema de}
Para $f\in L^1(\R^d)$ definimos
\[
 \widehat f(\xi)=\int_{\R^d}f(x)e^{ix\cdot\xi}\,dx.
\]
La integral existe para todo $\xi$, $|\widehat f(\xi)|\leq\norm f_1$ y no depende del representante de $f$. El signo difiere del empleado para los coeficientes periódicos; cambiarlo únicamente refleja la variable de frecuencia.
\begin{teorema}\label{v4:fourier-L1}
La transformada es una aplicación lineal continua de norma uno de $L^1(\R^d)$ en $C_0(\R^d)$. Además
\[
 \widehat{f*g}=\widehat f\,\widehat g,\quad
 \widehat{f(\cdot+a)}(\xi)=e^{-ia\cdot\xi}\widehat f(\xi),\quad
 \widehat{e^{ia\cdot x}f(x)}(\xi)=\widehat f(\xi+a).
\]
\end{teorema}
\begin{proof}
La linealidad y la cota de norma proceden de la integral. Si $f\geq0$, $\widehat f(0)=\norm f_1$, de modo que la norma del operador es uno. Para continuidad uniforme,
\[
 \sup_\xi|\widehat f(\xi+h)-\widehat f(\xi)|
 \leq\int|f(x)|\,|e^{ix\cdot h}-1|\,dx\longrightarrow0
\]
por convergencia dominada, con mayorante $2|f|$.

Para la indicadora de un rectángulo acotado $Q=\prod_j[a_j,b_j]$, Fubini expresa la transformada como el producto de las integrales unidimensionales $\int_{a_j}^{b_j}e^{ix_j\xi_j}dx_j$. Cada factor está acotado por $b_j-a_j$ y, si $\xi_j\ne0$, por $2/|\xi_j|$. Si $|\xi|\to\infty$, alguna coordenada tiene módulo al menos $|\xi|/\sqrt d$, y el producto tiende a cero. Las funciones escalonadas en rectángulos son densas en $L^1(\R^d)$ por regularidad de Lebesgue; la cota $\norm{\widehat f-\widehat g}_\infty\leq\norm{f-g}_1$ extiende la desaparición en infinito a todo $f$.

Para convolución, la integral doble de $|f(y)g(x-y)|$ es $\norm f_1\norm g_1$. Fubini y $z=x-y$ dan
\[
 \widehat{f*g}(\xi)=\int\int f(y)g(z)e^{i(y+z)\cdot\xi}\,dzdy
 =\widehat f(\xi)\widehat g(\xi).
\]
Los cambios de variable $y=x+a$ y la identidad $e^{ia\cdot x}e^{i\xi\cdot x}=e^{i(\xi+a)\cdot x}$ prueban las dos últimas fórmulas.
\end{proof}

\section{Caracteres y derivación}\label{lectura:v4-11}
\index{caracter@carácter!de la recta}
\begin{proposicion}\label{v4:caracteres-R}
Todo homomorfismo continuo $\chi:(\R,+)\to\C^\times$ tiene la forma $\chi(t)=e^{zt}$, $z\in\C$. Si toma valores en $\T$, entonces $z=i\xi$ para un único $\xi\in\R$. Los caracteres de $\R^d$ son exactamente $x\mapsto e^{ix\cdot\xi}$.
\end{proposicion}
\begin{proof}
Sea $F(t)=\int_0^t\chi(s)ds$. Como $F'(0)=\chi(0)=1$, existe $a$ con $F(a)\ne0$. La propiedad de homomorfismo da
\[
 F(a)\chi(t)=\int_t^{t+a}\chi(s)ds=F(t+a)-F(t).
\]
Así $\chi$ es diferenciable. Dividir $\chi(t+h)-\chi(t)=\chi(t)(\chi(h)-1)$ entre $h$ y pasar al límite da $\chi'(t)=\chi'(0)\chi(t)$. La derivada de $e^{-zt}\chi(t)$ es cero para $z=\chi'(0)$; su valor en cero es uno, por lo que $\chi(t)=e^{zt}$. Si $|\chi(t)|=1$, entonces $e^{t\Re z}=1$ para todo $t$, de donde $\Re z=0$. La derivada en cero prueba unicidad. En $\R^d$, restringimos a cada eje y multiplicamos las imágenes de $x_1e_1+\cdots+x_de_d$.
\end{proof}
\begin{proposicion}\label{v4:derivacion-fourier}
Si $x_jf\in L^1(\R^d)$, entonces $\partial_{\xi_j}\widehat f=\widehat{ix_jf}$. En la recta, si $f$ es localmente absolutamente continua y $f,f'\in L^1$, entonces
\[
 \widehat{f'}(\xi)=-i\xi\widehat f(\xi),\qquad
 f(x)=-\int_x^\infty f'(t)dt=\int_{-\infty}^xf'(t)dt.
\]
\end{proposicion}
\begin{proof}
El cociente incremental del exponencial está acotado por $|x_j|$, pues $|e^{iu}-1|\leq|u|$. Convergencia dominada permite derivar bajo la integral; además el error del cociente incremental está dominado independientemente de $\xi$, de modo que la convergencia es uniforme en la frecuencia.

Para la segunda afirmación, la continuidad absoluta local da $f(b)-f(a)=\int_a^bf'$. Como $f'$ es integrable, los límites de $f$ en ambos infinitos existen. Deben ser cero: un límite no nulo contradice $f\in L^1$. Pasar al límite en la identidad anterior da las dos fórmulas. La integración por partes en $[-R,R]$ produce
\[
 \int_{-R}^Rf'(x)e^{ix\xi}dx
 =[f(x)e^{ix\xi}]_{-R}^R-i\xi\int_{-R}^Rf(x)e^{ix\xi}dx.
\]
El término de frontera tiende a cero y las integrales convergen absolutamente. La hipótesis de continuidad absoluta local se hace explícita porque la mera existencia de una derivada integrable no justifica automáticamente el teorema fundamental de Lebesgue.
\end{proof}

\section{La gaussiana como puente de inversión}
\index{gaussiana}\index{nucleo@núcleo!del calor}
\begin{lema}\label{v4:gaussiana}
Se cumple $\int_\R e^{-x^2}dx=\sqrt\pi$ y
\[
 \widehat{e^{-a|x|^2}}(\xi)=\left(\frac\pi a\right)^{d/2}e^{-|\xi|^2/(4a)}\quad(a>0).
\]
Por tanto, para $t>0$,
\[
 H_t(x)=(4\pi t)^{-d/2}e^{-|x|^2/(4t)},\quad
 \widehat H_t(\xi)=e^{-t|\xi|^2},\quad
 H_t(x)=(2\pi)^{-d}\int e^{-t|\xi|^2}e^{-ix\cdot\xi}d\xi.
\]
\end{lema}
\begin{proof}
La integral $I=\int e^{-x^2}dx$ es positiva y finita por comparación en las colas con $e^{-|x|}$. Tonelli y coordenadas polares dan
\[
 I^2=\int_{\R^2}e^{-(x^2+y^2)}dxdy
 =2\pi\int_0^\infty re^{-r^2}dr=\pi.
\]
Si $F(\xi)=\int e^{-x^2}e^{ix\xi}dx$, la proposición anterior y $(e^{-x^2})'=-2xe^{-x^2}$ dan $F'(\xi)=-\xi F(\xi)/2$. Puesto que $F(0)=\sqrt\pi$, se obtiene $F(\xi)=\sqrt\pi e^{-\xi^2/4}$. El cambio $u=\sqrt a\,x$ y la factorización en coordenadas por Fubini prueban la fórmula general. Aplicar esa fórmula directamente a la gaussiana en la variable $\xi$ da la última identidad; todavía no se ha usado ningún teorema de inversión.
\end{proof}
Las funciones $H_t$ son no negativas y tienen integral uno. Para cada $\delta>0$, el cambio $x=\sqrt t\,u$ muestra $\int_{|x|\geq\delta}H_t(x)dx\to0$ cuando $t\downarrow0$. Forman, por tanto, una identidad aproximada y $f*H_t\to f$ en $L^p$ para $p<\infty$.
\begin{teorema}[Inversión]\label{v4:inversion-Rn}
Si $f\in L^1(\R^d)$ y $\widehat f\in L^1(\R^d)$, entonces
\[
 f(x)=(2\pi)^{-d}\int_{\R^d}\widehat f(\xi)e^{-ix\cdot\xi}d\xi
\]
casi en todas partes. La igualdad vale en todo punto de continuidad de $f$ y, más generalmente, en todo punto de Lebesgue.
\end{teorema}
\begin{proof}
La fórmula explícita de $H_t$ y Fubini, cuya hipótesis se verifica con $\norm f_1\int e^{-t|\xi|^2}d\xi<\infty$, dan
\[
 f*H_t(x)=(2\pi)^{-d}\int\widehat f(\xi)e^{-t|\xi|^2}e^{-ix\cdot\xi}d\xi.
\]
Si $\widehat f\in L^1$, convergencia dominada hace tender el lado derecho uniformemente en $x$ a la integral de inversión: la cota del error es $(2\pi)^{-d}\int|\widehat f||e^{-t|\xi|^2}-1|$. Por otro lado, $f*H_t\to f$ en $L^1$. Una sucesión $t_j\downarrow0$ posee una subsucesión convergente casi en todas partes, pues se puede escoger con errores $L^1$ sumables y aplicar Tonelli. La unicidad del límite prueba la fórmula casi en todas partes.

Precisemos la afirmación puntual. Sea $x$ punto de Lebesgue y $A(r)=\int_{|y|<r}|f(x-y)-f(x)|dy=o(r^d)$. Escribamos $h_t(r)=(4\pi t)^{-d/2}e^{-r^2/(4t)}$. Como $h_t(r)=\int_r^\infty-h_t'(s)ds$, Tonelli da
\[
 \int |f(x-y)-f(x)|H_t(y)dy
 =\int_0^\infty[-h_t'(s)]A(s)ds.
\]
Dado $\eps>0$, elegimos $\delta$ con $A(s)\leq\eps v_ds^d$ para $s<\delta$, donde $v_d$ es el volumen de la bola unidad. La integral en ese intervalo es a lo sumo $\eps$, porque $\int_0^\infty[-h_t'(s)]v_ds^d ds=\int H_t=1$. Para $s\geq\delta$ usamos $A(s)\leq\norm f_1+|f(x)|v_ds^d$. La contribución queda acotada por
\[
 \norm f_1h_t(\delta)+|f(x)|\left(v_d\delta^dh_t(\delta)+\int_{|y|\geq\delta}H_t(y)dy\right),
\]
integrando por partes el segundo término. Esta cota tiende a cero: la exponencial $e^{-\delta^2/(4t)}$ domina cualquier potencia de $1/t$ y la cola gaussiana tiende a cero. Así $f*H_t(x)\to f(x)$ y la fórmula puntual queda probada. Todo punto de continuidad es de Lebesgue.
\end{proof}
La prueba mediante calor sigue el enfoque de Maximenko~\cite{maximenkoInversion}, desarrollando también las constantes y el argumento en $\R^d$. En particular se aplica al error $f(x-y)-f(x)$, indispensable en el cálculo del punto de Lebesgue.

\section{Inversión truncada y localización en la recta}
\begin{proposicion}[Criterio de Dini en $\R$]\label{v4:dini-R}
Sea $f\in L^1(\R)$ y $S_Rf(x)=(2\pi)^{-1}\int_{-R}^R\widehat f(\xi)e^{-ix\xi}d\xi$. Si
\[
 \int_0^\delta\frac{|f(x+t)+f(x-t)-2f(x)|}{t}dt<\infty
\]
para algún $\delta>0$, entonces $S_Rf(x)\to f(x)$.
\end{proposicion}
\begin{proof}
Fubini se aplica al dominio finito de frecuencias, pues la integral absoluta es a lo sumo $2R\norm f_1$. Integrando el exponencial se obtiene
\[
 S_Rf(x)=\frac1\pi\int_0^\infty[f(x+t)+f(x-t)]\frac{\sin Rt}{t}dt.
\]
Usamos la integral impropia $\int_0^\infty\sin(Rt)/t\,dt=\pi/2$ para $R>0$. Para verificarla, la integral amortiguada $J_a(R)=\int_0^\infty e^{-at}\sin(Rt)/t\,dt$ satisface $J_a'(R)=a/(a^2+R^2)$ y $J_a(0)=0$, luego $J_a(R)=\arctan(R/a)$. Integración por partes da la cota uniforme $|\int_M^Ne^{-at}\sin(Rt)/t\,dt|\leq3/(RM)$ para $a\geq0$. Se pasa primero $a\downarrow0$ en un intervalo finito y después $M\to\infty$, obteniendo $\pi/2$.

Restamos $f(x)$ usando esa integral y separamos $(0,\delta)$ del resto. En el primer intervalo, la condición de Dini y Riemann--Lebesgue hacen tender a cero la integral del error. En $(\delta,\infty)$, $[f(x+t)+f(x-t)]/t$ es integrable, por lo que su integral oscilante tiende a cero. El término constante restante es un múltiplo de $\int_\delta^\infty\sin(Rt)/t\,dt$, acotado por $3/(R\delta)$ y por ello convergente a cero. Es necesario tratar este término aparte: la función constante dividida entre $t$ no es integrable en infinito, contrario al uso directo de Riemann--Lebesgue sugerido en la lectura 10.
\end{proof}

\section{Ejercicios de la transformada integrable}
\begin{ejercicio}[Lectura 10, ejercicio 1: transferencia]\label{ej:v4-l10-1}
Si $f,g\in L^1(\R^d)$, probar $\int f\widehat g=\int\widehat f g$.
\end{ejercicio}

\begin{ejercicio}[Lectura 10, ejercicio 2: involución]\label{ej:v4-l10-2}
Definir $f^*(x)=\overline{f(-x)}$. Probar que es una involución isométrica de $L^1(\R^d)$ y que $\widehat{f^*}=\overline{\widehat f}$. Explicar por qué $L^1(\R)$ no es un álgebra $C^*$ con esa norma.
\end{ejercicio}

\begin{ejercicio}[Lectura 10, ejercicio 3]\label{ej:v4-l10-3}
Si $(k_j)$ es una sucesión de Dirac en $L^1(\R^d)$, demostrar $\widehat k_j(\xi)\to1$ y que, para todo $h\in C_0(\R^d)$, $\norm{h\widehat k_j-h}_\infty\to0$.
\end{ejercicio}


\section{Ejercicios de inversión y funciones de Schwartz}
\index{Schwartz!espacio de}
\begin{ejercicio}[Lectura 11, ejercicio 1]\label{ej:v4-l11-1}
Si $f\in L^1(\R^d)$ y $\widehat f=0$, demostrar que $f=0$ casi en todas partes.
\end{ejercicio}

\begin{ejercicio}[Lectura 11, ejercicio 2: factor de normalización]\label{ej:v4-l11-2}
Calcular la transformada de $\mathbf1_{[-1,1]}$ y demostrar que no es integrable.
\end{ejercicio}

\begin{ejercicio}[Lectura 11, ejercicio 3: Parseval integrable]\label{ej:v4-l11-3}
Si $f,g,\widehat f,\widehat g\in L^1(\R^d)$, demostrar
$\int\widehat f\,\overline{\widehat g}=(2\pi)^d\int f\overline g$.
\end{ejercicio}

\begin{ejercicio}[Lectura 11, ejercicio 4]\label{ej:v4-l11-4}
Para $\varphi(x)=e^{-x^2}$, calcular su integral y transformada. Probar que $k_n(x)=n\pi^{-1/2}e^{-n^2x^2}$ forma una sucesión de Dirac que satisface la fórmula de inversión.
\end{ejercicio}

\begin{ejercicio}[Lectura 11, ejercicio 5]\label{ej:v4-l11-5}
Una función $f\in C^\infty(\R^d)$ es de Schwartz si $\sup_x|x^\alpha\partial^\beta f(x)|<\infty$ para todos los multiíndices $\alpha,\beta$. Probar: (i) $C_c^\infty\subset\Sscr$ y las gaussianas pertenecen a $\Sscr$ sin tener soporte compacto; $\Sscr$ es denso en $L^p$, $p<\infty$; (ii) Fourier preserva $\Sscr$; (iii) $L^1(\R)$ tiene divisores de cero.
\end{ejercicio}


\section{Los caracteres del álgebra de convolución}\label{lectura:v4-12}
\index{caracter@carácter!de L1}\index{Gelfand!transformada de}
Un carácter de un álgebra de Banach compleja es un funcional lineal multiplicativo, continuo y no nulo. Cada evaluación $\phi_\xi(f)=\widehat f(\xi)$ es un carácter de $L^1(\R^d)$ por el teorema de convolución; es no nulo, por ejemplo, sobre una gaussiana modulada.
\begin{lema}
Todo carácter $\phi$ de un álgebra de Banach satisface $\norm\phi\leq1$.
\end{lema}
\begin{proof}
Si $|\phi(a)|>\norm a$, pongamos $b=a/\phi(a)$. Entonces $\phi(b)=1$ y $\norm b<1$. La serie $x=-\sum_{n\geq1}b^n$ converge absolutamente y satisface $b+x=bx$. Aplicar $\phi$ daría $1+\phi(x)=\phi(x)$, imposible. El argumento no requiere una unidad en el álgebra.
\end{proof}
\begin{teorema}\label{v4:caracteres-L1R}
Todo carácter de $L^1(\R^d)$ es $f\mapsto\widehat f(\xi)$ para una única frecuencia $\xi$. Esta identificación es un homeomorfismo entre $\R^d$ y el espacio de caracteres con su topología débil estrella.
\end{teorema}
\begin{proof}
Por la representación del dual de $L^1(\R^d)$ existe $h\in L^\infty$ con $\phi(f)=\int fh$. Para $R_yf(x)=f(x-y)$, Fubini y multiplicatividad dan, para cada $f,g\in L^1$,
\[
 \int g(y)\phi(R_yf)dy=\phi(f*g)=\phi(f)\phi(g)=\int g(y)\phi(f)h(y)dy.
\]
La integral doble es absolutamente convergente por $\norm h_\infty\norm f_1\norm g_1$. Como $g$ es arbitraria, $\phi(R_yf)=\phi(f)h(y)$ para casi todo $y$.

Elijamos $f_0$ con $\phi(f_0)\ne0$ y definamos $H(y)=\phi(R_yf_0)/\phi(f_0)$. La continuidad de traslaciones hace a $H$ continua y acotada, y $H=h$ casi en todas partes. Para cualquier $f$, las funciones continuas $y\mapsto\phi(R_yf)$ y $y\mapsto\phi(f)H(y)$ coinciden casi en todas partes, luego en todo punto: una diferencia continua no nula sería no nula en un abierto de medida positiva. Así se obtiene una identidad válida para cada $f$ y todo $y$, sin elegir ilegítimamente un único nulo para todos los $f$.

Aplicándola a $R_zf_0$, se deduce $H(y+z)=H(y)H(z)$ y $H(0)=1$. En particular $H$ nunca se anula. La acotación y las potencias positivas y negativas implican $|H(y)|=1$; si fuese mayor o menor que uno, alguna sucesión de potencias sería no acotada. Por la proposición~\ref{v4:caracteres-R}, $H(y)=e^{iy\cdot\xi}$, y por tanto $\phi=\phi_\xi$. Si dos frecuencias dan el mismo funcional, sus exponenciales son iguales casi en todas partes por dualidad y luego en todo punto por continuidad; restringir a los ejes y derivar en cero prueba que las frecuencias coinciden.

La continuidad de $\xi\mapsto\phi_\xi$ es la continuidad de cada $\widehat f$. Para la inversa, si una red $\phi_{\xi_\alpha}\to\phi_\xi$ débil estrella, evaluamos en $f(x)=e^{-|x|^2}e^{-ix\cdot\xi}$. Su transformada es $\pi^{d/2}e^{-|\eta-\xi|^2/4}$. La convergencia de esas evaluaciones implica $e^{-|\xi_\alpha-\xi|^2/4}\to1$, y de ahí $\xi_\alpha\to\xi$. Esto prueba continuidad inversa incluso para redes.
\end{proof}
La transformada de Gelfand asigna a $f$ la función $\phi\mapsto\phi(f)$ sobre el espacio de caracteres. Bajo el homeomorfismo anterior es exactamente la transformada de Fourier. La lectura 12 no contiene una lista de ejercicios; sus resultados de norma, continuidad del representante y topología del espectro se han demostrado aquí.

\paragraph{Del álgebra a la geometría.}
La transformada ya distingue funciones integrables y respeta su producto de convolución. El siguiente paso cambia de norma: Plancherel extenderá Fourier a todo $L^2$, incluso cuando la integral que lo define no converja absolutamente.

\chapter{Plancherel, invariancia y muestreo}\label{cap:v4-08}
La integral de Fourier no está definida absolutamente para toda función de $L^2$. La identidad de Plancherel permite extenderla por continuidad, mientras que la inversión determina su rango. Esta extensión transforma problemas de traslación en problemas de multiplicación y conduce a una reconstrucción concreta de funciones de banda limitada.

\section{La extensión unitaria}\label{lectura:v4-13}
\index{Plancherel, teorema de}\index{Fourier!transformada unitaria}
\begin{teorema}[Plancherel en $\R^d$]\label{v4:plancherel}
Si $f\in L^1(\R^d)\cap L^2(\R^d)$, entonces $\widehat f\in L^2$ y $\norm{\widehat f}_2=(2\pi)^{d/2}\norm f_2$. Por tanto
\[
 \mathcal Ff=(2\pi)^{-d/2}\widehat f
\]
se extiende de manera única a un operador unitario de $L^2(\R^d)$ sobre sí mismo. Su inversa y adjunto es la extensión de
\[
 \mathcal F^{-1}g(x)=(2\pi)^{-d/2}\int g(\xi)e^{-ix\cdot\xi}d\xi.
\]
\end{teorema}
\begin{proof}
Sea $f^*(x)=\overline{f(-x)}$ y $g=f*f^*$. Por Young, $g\in L^1$; por el teorema~\ref{v4:conv-continuidad}, $g\in C_0$, y $g(0)=\int|f|^2$. Además $\widehat g=|\widehat f|^2$, aun sin saber que esta última función sea integrable. La fórmula gaussiana ya demostrada da
\[
 (g*H_t)(0)=(2\pi)^{-d}\int |\widehat f(\xi)|^2e^{-t|\xi|^2}d\xi.
\]
La izquierda tiende a $g(0)$ porque $g$ es continua y acotada. Al tomar $t=1/n$, el integrando de la derecha crece hacia $|\widehat f|^2$; convergencia monótona prueba simultáneamente su integrabilidad y la identidad de normas. No se aplica inversión a $g$ antes de probar que su transformada es integrable.

El subespacio $L^1\cap L^2$ es denso en $L^2$: truncar una función en espacio y altura produce funciones de ese subespacio que convergen por dominada. Para $f_n\to f$ en $L^2$, la igualdad de normas hace a $(\mathcal Ff_n)$ de Cauchy. Su límite define $\mathcal Ff$, es independiente de la aproximación y conserva la norma. Repetir la construcción con signo negativo define otra isometría $V$.

Si $f,g\in L^1\cap L^2$, Fubini da $\ip{\mathcal Ff}g=\ip f{Vg}$: la integral doble absoluta se controla por $\norm f_1\norm g_1$. Por densidad, $V=\mathcal F^*$ en todo $L^2$. Como ambas aplicaciones son isometrías, $\mathcal F^*\mathcal F=I$ y $V^*V=I$. La segunda identidad es $\mathcal F\mathcal F^*=I$, luego $\mathcal F$ es sobreyectiva y $V$ su inversa.
\end{proof}
La identidad de producto interno es $\ip{\mathcal Ff}{\mathcal Fg}=\ip fg$. Si se conserva la notación no normalizada $\widehat f=(2\pi)^{d/2}\mathcal Ff$ para clases de $L^2$, se escribe
\[
 \int|f|^2=(2\pi)^{-d}\int|\widehat f|^2.
\]
No alternaremos esas dos normalizaciones bajo un mismo símbolo, como ocurre en algunos pasos de la lectura 13. Las fórmulas de traslación y modulación pasan a $L^2$ aproximando por $L^1\cap L^2$ y usando continuidad de los operadores implicados.
\begin{ejemplo}
Para $f=\mathbf1_{[-1,1]}$, $\widehat f(\xi)=2\sin\xi/\xi$. Plancherel da $4\int_\R\sin^2\xi/\xi^2\,d\xi=2\pi\cdot2$, es decir,
\[
 \int_\R\frac{\sin^2\xi}{\xi^2}d\xi=\pi.
\]
La transformada no está en $L^1$, pero sí en $L^2$; este es precisamente el tipo de situación que justifica la extensión.
\end{ejemplo}

\section{Ejercicios de Plancherel}
\begin{ejercicio}[Lectura 13, ejercicio 1]\label{ej:v4-l13-1}
Para $t\in\R$, demostrar
\[
 \int_{-1}^1\frac{\sin(t-y)}{t-y}dy
 =\int_{-1}^1\frac{\sin y}{y}e^{iyt}dy,
\]
prolongando los cocientes por continuidad en cero.
\end{ejercicio}

\begin{ejercicio}[Lectura 13, ejercicio 2]\label{ej:v4-l13-2}
Calcular la transformada de $f(x)=e^{-|x|}$ y la integral $\int_\R t^2/(1+t^2)^2\,dt$.
\end{ejercicio}

\begin{ejercicio}[Lectura 13, ejercicio 3]\label{ej:v4-l13-3}
Para $a>0$ sean $p_a(x)=a/[\pi(x^2+a^2)]$ y $s_a(x)=\sin(ax)/(\pi x)$. Probar $p_a*p_b=p_{a+b}$ y $s_a*s_b=s_{\min(a,b)}$.
\end{ejercicio}

\begin{ejercicio}[Lectura 13, ejercicio 4]\label{ej:v4-l13-4}
Demostrar, para $a,b>0$,
\[
 \int_\R\frac{\sin(at)\sin(bt)}{t^2}dt=\pi\min(a,b),\qquad
 \int_\R\frac{t^2}{(t^2+a^2)(t^2+b^2)}dt=\frac\pi{a+b}.
\]
\end{ejercicio}


\section{Subespacios invariantes por traslación}\label{lectura:v4-14}
\index{subespacio!invariante por traslacion@invariante por traslación}\index{multiplicador!de Fourier}
\begin{teorema}\label{v4:subespacios-invariantes}
Para un conjunto medible $A\subseteq\R^d$, sea
\[
 H_A=\{f\in L^2(\R^d):\mathcal Ff=0\text{ casi en todas partes en }A\}.
\]
Entonces $H_A$ es cerrado e invariante bajo todas las traslaciones; $H_A=H_B$ si y solo si $A\triangle B$ es nulo; y todo subespacio \emph{cerrado} invariante bajo traslaciones es de esa forma.
\end{teorema}
\begin{proof}
El espacio de funciones que se anulan en $A$ es el núcleo del operador acotado $g\mapsto\mathbf1_Ag$, por lo que es cerrado. Su preimagen bajo la aplicación unitaria $\mathcal F$ es $H_A$. Las traslaciones se convierten en multiplicación por $e^{-ia\cdot\xi}$; multiplicar por ese factor no cambia el conjunto de ceros, lo que prueba invariancia.

Si $A\triangle B$ es nulo, las condiciones de anulación coinciden. Recíprocamente, si los espacios coinciden y $|B\setminus A|>0$, alguna intersección $(B\setminus A)\cap[-n,n]^d$ tiene medida positiva y finita. Su indicadora pertenece al espacio de funciones que se anulan en $A$, pero no al que se anula en $B$, contradicción. Se intercambian $A,B$ para concluir igualdad módulo nulos.

Sea ahora $H$ cerrado e invariante por todas las traslaciones, y $W=\mathcal F(H)$. Entonces $W$ es cerrado e invariante por multiplicación por todos los caracteres, incluidas sus inversas. Su ortogonal $W^\perp$ también es invariante: si $u\in W$, $v\in W^\perp$, entonces $\ip{e^{ia\cdot\xi}v}u=\ip v{e^{-ia\cdot\xi}u}=0$. Para esos $u,v$, la función $u\overline v$ pertenece a $L^1$ por Hölder y
\[
 \int e^{ia\cdot\xi}u(\xi)\overline{v(\xi)}d\xi=0\quad(a\in\R^d).
\]
La unicidad de Fourier en $L^1$ implica $u\overline v=0$ casi en todas partes.

Por separabilidad de $L^2(\R^d)$, el subespacio $W$ tiene una sucesión densa $(u_j)$. Sea $E=\bigcup_j\{u_j\ne0\}$. Todo elemento de $W$ se anula fuera de $E$: esto es cierto para $u_j$ y pasa al límite en $L^2$ por continuidad de multiplicación por $\mathbf1_{E^c}$. Para cada $v\in W^\perp$, las identidades $u_j\overline v=0$, reunidas fuera de una unión numerable de nulos, implican que $v=0$ en $E$ casi en todas partes. Si $g$ está apoyada en $E$, descomponemos $g=w+v$ con $w\in W$, $v\in W^\perp$. Tanto $g$ como $w$ se anulan fuera de $E$, así que $v$ también; pero $v$ se anula en $E$, luego $v=0$. Por tanto $W=L^2(E)$, y $H=H_{E^c}$.
\end{proof}
La proyección ortogonal sobre $H_A$ es $\mathcal F^{-1}M_{\mathbf1_{A^c}}\mathcal F$. Así la descripción mediante proyecciones de la lectura original se concreta en un multiplicador indicador. La palabra «cerrado» no puede omitirse: $\Sscr(\R)$ es un subespacio propio, denso e invariante por traslaciones de $L^2$, y no puede ser ninguno de los espacios cerrados $H_A$. La lectura 14 no contiene ejercicios numerados.

\section{Banda limitada y núcleo reproductor}
\index{banda limitada}\index{nucleo@núcleo!reproductor}\index{Paley--Wiener!espacio de}
Este complemento sigue el desarrollo de Lozano Vite y Maximenko~\cite{lozanoMaximenkoPaleyWiener}. Sea $a>0$ y
\[
 PW_a=\{f\in L^2(\R):\mathcal Ff=0\text{ casi en todas partes fuera de }[-a,a]\}.
\]
Toda clase de este espacio tiene un representante continuo canónico: si $g=\mathcal Ff$ está apoyada en $[-a,a]$, Hölder da $\norm g_1\leq\sqrt{2a}\norm g_2$, y se define
\[
 f(x)=\frac1{\sqrt{2\pi}}\int_{-a}^ag(\xi)e^{-ix\xi}d\xi.
\]
La integral representa la clase inversa en $L^2$, es continua y pertenece a $C_0$. En adelante $PW_a$ consta de esos representantes, con norma $L^2$.
\begin{proposicion}\label{v4:rkhs-pw}
La evaluación en $y$ es un funcional continuo de $PW_a$, representado por
\[
 K_y(x)=\frac{\sin(a(x-y))}{\pi(x-y)},\qquad K_y(y)=\frac a\pi.
\]
Se tiene $f(y)=\ip f{K_y}$ y $\norm f_\infty\leq\sqrt{a/\pi}\norm f_2$.
\end{proposicion}
\begin{proof}
Por Cauchy--Schwarz en la integral inversa,
\[
|f(y)|\leq(2\pi)^{-1/2}\sqrt{2a}\norm g_2=\sqrt{a/\pi}\norm f_2.
\]
La función $K_y$ es la inversa unitaria de $(2\pi)^{-1/2}\mathbf1_{[-a,a]}(\xi)e^{iy\xi}$, pues integrar el exponencial da la fórmula seno. Por Parseval,
\[
 \ip f{K_y}=\frac1{\sqrt{2\pi}}\int_{-a}^ag(\xi)e^{-iy\xi}d\xi=f(y).
\]
Su norma al cuadrado es $\norm{\mathcal FK_y}_2^2=(2\pi)^{-1}2a=a/\pi$. Este espacio es, por definición, un espacio de Hilbert con núcleo reproductor. La norma es la raíz cuadrada de la integral de $|f|^2$; no debe omitirse esa raíz.
\end{proof}

\section{Reconstrucción por muestras}
\index{muestreo!teorema de Whittaker--Shannon}\index{sinc}
Definimos $\operatorname{sinc}(t)=\sin(\pi t)/(\pi t)$ para $t\ne0$, con valor uno en cero.
\begin{teorema}[Whittaker--Shannon en $L^2$]\label{v4:muestreo}
Para $f\in PW_a$,
\[
 f(x)=\sum_{k\in\Z}f\left(\frac{k\pi}a\right)
 \operatorname{sinc}\left(\frac{ax}\pi-k\right),
 \qquad
 \norm f_2^2=\frac\pi a\sum_{k\in\Z}\left|f\left(\frac{k\pi}a\right)\right|^2.
\]
Las sumas simétricas convergen en $L^2$ y uniformemente en toda $\R$.
\end{teorema}
\begin{proof}
El cambio lineal de variable desde el círculo muestra que
\[
 \phi_k(\xi)=(2a)^{-1/2}e^{ik\pi\xi/a},\quad k\in\Z,
\]
es una base ortonormal de $L^2([-a,a],d\xi)$. Aplicando la inversa unitaria de Fourier obtenemos la base ortonormal
\[
 e_k(x)=\sqrt{\frac a\pi}\operatorname{sinc}\left(\frac{ax}\pi-k\right)
 =\sqrt{\frac\pi a}\,K_{k\pi/a}(x)
\]
de $PW_a$. La propiedad reproductora da $\ip f{e_k}=\sqrt{\pi/a}\,f(k\pi/a)$. La expansión en esa base y Parseval proporcionan las dos fórmulas del enunciado, con convergencia en norma. Si $r_N$ es el residuo de la suma parcial, todavía pertenece a $PW_a$; por la proposición anterior,
\[
 \norm{r_N}_\infty\leq\sqrt{a/\pi}\norm{r_N}_2\longrightarrow0.
\]
Por tanto la convergencia es uniforme en toda la recta, no únicamente en intervalos compactos. La demostración utiliza Plancherel y una base ortonormal; no requiere el teorema de Carleson.
\end{proof}
\begin{ejercicio}[Complemento: cambio de ancho de banda]\label{ej:v4-ampl-muestreo}
En la convención $\widehat f(\nu)=\int f(x)e^{2\pi ix\nu}dx$, resolver el ejercicio de Lozano Vite y Maximenko que sustituye la banda $[-1/2,1/2]$ por $[-b,b]$: hallar la separación entre muestras, el núcleo y Parseval.
\end{ejercicio}


\paragraph{Del espacio euclidiano al grupo.}
Traslación, medida invariante, caracteres y convolución han aparecido en cada construcción. Los capítulos siguientes abstraen esas cuatro estructuras. El trabajo previo en $\R^d$, $\T$ y los grupos finitos servirá para comprobar las constantes y reconocer los casos concretos de la teoría general.

\parte{Grupos topológicos y medida invariante}{Las operaciones algebraicas se hacen compatibles con la topología. Sobre grupos localmente compactos construimos medida de Haar, distinguimos invariancia izquierda y derecha y estudiamos la integración en cocientes.}
\chapter{La topología de un grupo}\label{cap:v4-09}
En un grupo topológico, conocer lo que ocurre cerca de la identidad permite trasladar esa información a todos los puntos. Esta homogeneidad explica los resultados de separación y la construcción de cocientes. No se presupone Hausdorff hasta que se indique; los capítulos de medida de Haar sí trabajarán con grupos localmente compactos y de Hausdorff.

\section{Operaciones continuas y vecindades}\label{lectura:v4-15}
\index{grupo!topologico@topológico}\index{vecindad!de la identidad}
\begin{definicion}
Un grupo topológico es un grupo $G$ con una topología para la cual $(x,y)\mapsto xy$ y $x\mapsto x^{-1}$ son continuas. Su identidad se denota $e$. Para subconjuntos, $AB=\{ab:a\in A,b\in B\}$ y $A^{-1}=\{a^{-1}:a\in A\}$. El conjunto $AA$ se distingue de $\{a^2:a\in A\}$.
\end{definicion}
Todo grupo discreto es topológico porque todas las aplicaciones desde un espacio discreto son continuas y el producto finito de espacios discretos es discreto. Los grupos aditivos $\R^d$, $\C$ y el grupo multiplicativo $\T$ con sus topologías usuales son ejemplos, por continuidad de sus operaciones. El producto cartesiano de grupos topológicos, con operaciones coordenadas y topología producto, es otro grupo topológico: una aplicación al producto es continua si lo son todas sus coordenadas.
\begin{proposicion}
Las traslaciones $x\mapsto ax$, $x\mapsto xa$ y la inversión son homeomorfismos. Si $\mathcal V$ es una base de vecindades abiertas de $e$, los conjuntos $aV$, $a\in G$, $V\in\mathcal V$, forman una base de la topología; lo mismo ocurre con $Va$.
\end{proposicion}
\begin{proof}
Cada traslación es continua por continuidad del producto con un factor fijo, y su inversa es la traslación por $a^{-1}$. La inversión es su propia inversa. Si $O$ es abierto y $a\in O$, $a^{-1}O$ es una vecindad abierta de $e$; contiene algún $V\in\mathcal V$, y $a\in aV\subset O$. El argumento derecho es idéntico usando $Oa^{-1}$.
\end{proof}
Una traslación individual no es, en general, un homomorfismo de grupos: no manda $e$ en $e$. Lo que sí forma una representación es la acción sobre funciones $L_af(x)=f(a^{-1}x)$, para la cual $L_aL_b=L_{ab}$.
\begin{proposicion}
Si $A$ es abierto, $AB$ y $BA$ son abiertos para todo $B$. Si $A,B$ son compactos, $AB$ es compacto. Si $A$ es compacto y $B$ cerrado, $AB$ y $BA$ son cerrados.
\end{proposicion}
\begin{proof}
$AB=\bigcup_{b\in B}Ab$ es unión de abiertos trasladados, y de modo similar para $BA$. El producto de dos compactos es compacto, y $AB$ es su imagen continua bajo multiplicación.

Sea $x\notin AB$. Para cada $a\in A$, $a^{-1}x\notin B$. Como $B$ es cerrado y $(u,v)\mapsto u^{-1}v$ es continua, existen abiertos $U_a\ni a$, $V_a\ni x$ con $U_a^{-1}V_a\subseteq G\setminus B$. Una cantidad finita $U_{a_1},\ldots,U_{a_r}$ cubre $A$. Entonces $V=\bigcap_jV_{a_j}$ es una vecindad de $x$ disjunta de $AB$: si $v=ab$, escogemos $j$ con $a\in U_{a_j}$, lo que daría $b=a^{-1}v\notin B$. Así el complemento es abierto. Para $BA$ se usa la aplicación $(v,u)\mapsto vu^{-1}$.
\end{proof}
\begin{proposicion}\label{v4:vecindades-grupo}
Toda vecindad $U$ de $e$ contiene una vecindad abierta simétrica $V=V^{-1}$. Puede escogerse además con $VV\subset U$, e incluso con $\overline V\subset U$. Para cada compacto $K\subset G$ y vecindad $U$ de $e$, existe una vecindad $V$ de $e$ tal que $xVx^{-1}\subset U$ para todo $x\in K$.
\end{proposicion}
\begin{proof}
Si $W$ es abierta y $e\in W\subset U$, entonces $W\cap W^{-1}$ es abierta y simétrica. Por continuidad de multiplicación en $(e,e)$, existen vecindades $A,B$ con $AB\subset U$; tomar una vecindad simétrica dentro de $A\cap B$ da $VV\subset U$. Si $x\in\overline V$, la vecindad $xV$ de $x$ corta $V$, así que $xu=v$ para algunos $u,v\in V$. Entonces $x=vu^{-1}\in VV\subset U$.

Para la uniformidad de conjugación, elijamos una vecindad simétrica $W$ con $WWW\subset U$. Los abiertos $Wx$, $x\in K$, cubren $K$; bastan $Wx_1,\ldots,Wx_r$. Sea $V=\bigcap_jx_j^{-1}Wx_j$. Si $x=wx_j\in K$, con $w\in W$, entonces $xVx^{-1}\subset wWw^{-1}\subset WWW\subset U$.
\end{proof}
\begin{teorema}[Construcción por una base en la identidad]\label{v4:base-identidad}
Sea $G$ un grupo y $\mathcal V$ una familia no vacía de subconjuntos que contienen $e$, dirigida por refinamiento de intersecciones finitas. Supongamos:
\begin{enumerate}
\item para cada $U\in\mathcal V$ existe $V\in\mathcal V$ con $VV\subset U$;
\item para cada $U\in\mathcal V$ existe $V\in\mathcal V$ con $V^{-1}\subset U$;
\item para cada $U\in\mathcal V$ y $x\in U$ existe $V\in\mathcal V$ con $xV\subset U$;
\item para cada $U\in\mathcal V$ y $x\in G$ existe $V\in\mathcal V$ con $xVx^{-1}\subset U$.
\end{enumerate}
Entonces los conjuntos $xU$ forman una base de una topología de grupo y $\mathcal V$ es una base de vecindades abiertas de $e$.
\end{teorema}
\begin{proof}
Todos los puntos están cubiertos porque $x\in xU$. Si $z\in xU\cap yV$, (iii) proporciona $A,B\in\mathcal V$ con $(x^{-1}z)A\subset U$ y $(y^{-1}z)B\subset V$. Elegimos $W\subset A\cap B$ en la familia, y entonces $zW\subset xU\cap yV$. Esto prueba el axioma de base.

Para continuidad del producto en $(x,y)$, consideremos una vecindad básica $xyU$. Escojamos $W$ con $WW\subset U$, $A$ con $y^{-1}Ay\subset W$ y $B\subset W$. Entonces
\[
 (xA)(yB)=xy(y^{-1}Ay)B\subset xyWW\subset xyU.
\]
Para inversión, dada $x^{-1}U$, combinamos (ii) y (iv) para escoger $A$ con $xA^{-1}x^{-1}\subset U$. Se obtiene $(xA)^{-1}=x^{-1}(xA^{-1}x^{-1})\subset x^{-1}U$. Las operaciones son continuas. Las propiedades (i)--(iv) son necesarias para una base abierta real, por continuidad y por definición de abierto. La condición (iii) exige $x\in U$: pedirla para todo $x\in G$, como en la lectura 15, forzaría $U=G$.
\end{proof}

\section{Separación, subgrupos y cocientes}\label{lectura:v4-16}
\index{grupo!cociente}\index{subgrupo!abierto}
\begin{proposicion}
Para un grupo topológico son equivalentes: ser $T_1$, ser Hausdorff y que $\{e\}$ sea cerrado. La clausura de un subgrupo es un subgrupo, y todo subgrupo abierto es cerrado.
\end{proposicion}
\begin{proof}
Si $\{e\}$ es cerrado, todos los puntos lo son por traslación, y el espacio es $T_1$. Si es $T_1$ y $x\ne y$, elegimos una vecindad $U$ de $e$ que no contenga $xy^{-1}$, y una vecindad abierta simétrica $V$ con $VV\subset U$. Los abiertos $Vx,Vy$ son disjuntos: una igualdad $vx=wy$ daría $xy^{-1}=v^{-1}w\in VV$. Así es Hausdorff. La implicación Hausdorff a puntos cerrados es general.

Si $H$ es subgrupo y $x,y\in\overline H$, la continuidad de $(u,v)\mapsto uv^{-1}$ implica que toda vecindad de $xy^{-1}$ contiene algún $hk^{-1}$ con $h,k\in H$: se eligen vecindades de $x,y$ cuyo producto con inversión esté en ella. Por tanto $xy^{-1}\in\overline H$. Si $H$ es abierto, todas sus clases izquierdas son abiertas; su complemento es unión de las clases distintas de $H$, luego también es abierto.
\end{proof}
Para un subgrupo $H\leq G$, el conjunto $G/H$ de clases izquierdas lleva la topología cociente: $O\subset G/H$ es abierto si y solo si $q^{-1}(O)$ es abierto en $G$. Es la topología \emph{más fina} que hace continua a $q$; la descripción como «más gruesa» en el original está invertida.
\begin{teorema}\label{v4:cociente-topologico}
La proyección $q:G\to G/H$ es abierta. La acción $G\times G/H\to G/H$, $(x,yH)\mapsto xyH$, es continua y abierta. El cociente es $T_1$, y de hecho Hausdorff, si y solo si $H$ es cerrado; es discreto si y solo si $H$ es abierto. Si $H$ es normal, $G/H$ es un grupo topológico.
\end{teorema}
\begin{proof}
Si $U$ es abierto en $G$, $q^{-1}(q(U))=UH=\bigcup_{h\in H}Uh$ es abierto, luego $q(U)$ es abierto. La aplicación $\mathrm{id}_G\times q$ es continua, abierta y sobreyectiva, por lo que es una aplicación cociente. La composición de la acción con ella es $q\circ m$, donde $m$ es multiplicación en $G$; es continua, y la propiedad cociente prueba continuidad de la acción. Para un abierto básico $U\times V$, su imagen es $\bigcup_{x\in U}xV$, abierta, porque cada traslación del cociente es un homeomorfismo.

La preimagen del punto $xH$ es $xH$, de modo que todos los puntos son cerrados exactamente cuando $H$ lo es; el mismo argumento con abiertos caracteriza discreción. Si $H$ es cerrado y $xH\ne yH$, la continuidad de $(u,v)\mapsto v^{-1}u$ y $y^{-1}x\notin H$ dan abiertos $U\ni x$, $V\ni y$ con $V^{-1}U\cap H=\varnothing$. Sus imágenes abiertas por $q$ son disjuntas, pues compartir una clase produciría $v^{-1}u\in H$. Esto prueba Hausdorff sin exigir normalidad.

Si $H$ es normal, las operaciones algebraicas del cociente están bien definidas. La aplicación $q\times q$ es abierta, sobreyectiva y continua, por tanto cociente. La identidad $m_{G/H}\circ(q\times q)=q\circ m_G$ prueba continuidad del producto; la identidad correspondiente para inversión y el carácter cociente de $q$ prueban continuidad de la inversa.
\end{proof}

\section{La componente de la identidad}
\index{componente!de la identidad}\index{grupo!totalmente disconexo}
\begin{teorema}
La componente conexa $G^\circ$ de $e$ es un subgrupo normal cerrado, y $G/G^\circ$ es Hausdorff y totalmente disconexo.
\end{teorema}
\begin{proof}
La imagen de $G^\circ$ por inversión es conexa y contiene $e$, luego está contenida en $G^\circ$. Si $x\in G^\circ$, $xG^\circ$ es conexo y contiene $e$ porque $x^{-1}\in G^\circ$; así $xG^\circ\subset G^\circ$. Esto prueba que es subgrupo. Para todo $a\in G$, la conjugación manda $G^\circ$ a un conexo que contiene $e$, por lo que $aG^\circ a^{-1}\subset G^\circ$. La conjugación inversa da igualdad. Las componentes conexas son cerradas, así que el teorema anterior hace Hausdorff al cociente.

Para la última afirmación usaremos el siguiente hecho: una aplicación abierta y sobreyectiva con fibras conexas tiene preimagen conexa de todo conjunto conexo. En efecto, si la preimagen de $S$ se separara en dos abiertos relativos $A,B$, cada fibra, por conexidad, estaría entera en uno de ellos. Sus imágenes serían disjuntas, no vacías y abiertas en $S$: si $A=O\cap q^{-1}(S)$, entonces $q(A)=q(O)\cap S$. Eso separaría $S$.

Las fibras de $q:G\to G/G^\circ$ son las clases conexas $xG^\circ$. Si $S$ es un subconjunto conexo del cociente que contiene la identidad, su preimagen es conexa y contiene $G^\circ$, por lo que está contenida en la componente $G^\circ$. Entonces $S$ es un punto. Las traslaciones llevan este resultado a cualquier punto del cociente.
\end{proof}

\section{Ejercicios de estructuras topológicas}
\begin{ejercicio}[Lectura 15, ejercicio 1]\label{ej:v4-l15-1}
Demostrar que basta exigir continuidad de $(x,y)\mapsto xy^{-1}$ para obtener un grupo topológico.
\end{ejercicio}

\begin{ejercicio}[Lectura 15, ejercicio 2]\label{ej:v4-l15-2}
Probar que un grupo infinito con la topología cofinita no es un grupo topológico.
\end{ejercicio}

\begin{ejercicio}[Lectura 15, ejercicio 3]\label{ej:v4-l15-3}
Sea $P=\prod_{\chi\in\operatorname{Hom}(\Z,\T)}\T$ y $j(n)=(\chi(n))_\chi$. Probar: (i) $P$ es un grupo topológico compacto; (ii) $j$ es un monomorfismo; (iii) la topología que induce sobre $\Z$ no es discreta.
\end{ejercicio}

\begin{ejercicio}[Lectura 15, ejercicio 4]\label{ej:v4-l15-4}
Construir la topología profinita de un grupo $G$ a partir de sus subgrupos normales de índice finito.
\end{ejercicio}


\section{Ejercicios de subgrupos y compacidad}
\begin{ejercicio}[Lectura 16, ejercicio 1]\label{ej:v4-l16-1}
Probar que un subgrupo $H$ es abierto si y solo si tiene interior no vacío.
\end{ejercicio}

\begin{ejercicio}[Lectura 16, ejercicio 2]\label{ej:v4-l16-2}
Sea $\mathcal A$ una familia no vacía de vecindades de $e$, dirigida por refinamiento, tal que para cada $U$ hay miembros $V,W$ con $VV\subset U$ y $W^{-1}\subset U$. Probar que $H=\bigcap_{U\in\mathcal A}U$ es un subgrupo cerrado.
\end{ejercicio}

\begin{ejercicio}[Lectura 16, ejercicio 3]\label{ej:v4-l16-3}
Bajo las hipótesis anteriores, si para cada $U\in\mathcal A$ y $x\in G$ existe $V\in\mathcal A$ con $x^{-1}Vx\subset U$, probar que $H$ es normal.
\end{ejercicio}

\begin{ejercicio}[Lectura 16, ejercicio 4]\label{ej:v4-l16-4}
Probar que un subgrupo es discreto en su topología de subespacio si y solo si tiene un punto aislado.
\end{ejercicio}

\begin{ejercicio}[Lectura 16, ejercicio 5]\label{ej:v4-l16-5}
Si $G$ es Hausdorff y $H$ un subgrupo compacto, demostrar que $q:G\to G/H$ es cerrada.
\end{ejercicio}

\begin{ejercicio}[Lectura 16, ejercicio 6]\label{ej:v4-l16-6}
Si $H\leq G$ y $U\cap H$ es cerrado en $G$ para alguna vecindad $U$ de $e$, probar que $H$ es cerrado.
\end{ejercicio}

\begin{ejercicio}[Lectura 16, ejercicio 7]\label{ej:v4-l16-7}
Probar que $\overline{\{e\}}$ es un subgrupo normal.
\end{ejercicio}

\begin{ejercicio}[Lectura 16, ejercicio 8]\label{ej:v4-l16-8}
Si $F$ es compacto y $O$ es abierto con $F\subset O$, hallar una vecindad $V$ de $e$ con $FV\cup VF\subset O$.
\end{ejercicio}

\begin{ejercicio}[Lectura 16, ejercicio 9]\label{ej:v4-l16-9}
Si $F$ es una vecindad compacta y abierta de $e$, probar que contiene un subgrupo compacto y abierto, que también es cerrado en $G$.
\end{ejercicio}


\paragraph{Entrada a la integración invariante.}
Las vecindades pequeñas, los soportes compactos y los cocientes abiertos son las herramientas topológicas de la construcción de Haar. La local compacidad añadirá funciones continuas que localizan la integración sin perder invariancia por traslación.

\chapter{La medida de Haar}\label{v4:haar}\label{cap:v4-10}
\index{medida!de Haar}
La longitud de un intervalo no cambia al trasladarlo. La medida de Haar expresa esta idea en un grupo: mide conjuntos de modo compatible con la multiplicación. El problema no es escribir la condición de invariancia, sino construir una medida que la satisfaga sin coordenadas, derivadas ni una fórmula explícita para el grupo.

Este capítulo desarrolla las lecturas 17--19 de Valencia Bucio \cite{valencia4}; la lectura 19 lleva por error el encabezado «Lectura 18». Se completan la construcción por recubrimientos, su paso al límite y la prueba de unicidad. El contexto general y la convención de regularidad se contrastan con Folland \cite{folland2016}.

\section{Preparación topológica y regularidad}\label{lectura:v4-17}
En todo el capítulo, $G$ es un grupo localmente compacto y Hausdorff. Escribimos $\Cc(G)$ para las funciones complejas continuas de soporte compacto y
\[
 L_xf(t)=f(x^{-1}t),\qquad R_xf(t)=f(tx).
\]
El conjunto $\Cc^+(G)$ contendrá las funciones de $\Cc(G)$ no negativas; se dirá explícitamente cuando se excluye la función cero.

\begin{proposicion}\label{v4:haar-subgrupo}
Un subgrupo localmente compacto $H$ de un grupo Hausdorff $G$, con la topología inducida, es cerrado.
\end{proposicion}
\begin{proof}
Elijamos una vecindad compacta $K$ de la identidad en $H$. Existe un abierto $U$ de $G$ tal que $e\in U$ y $U\cap H\subset K$. La regularidad de los grupos Hausdorff permite elegir un abierto $V$ con $e\in V$ y $\overline V\subset U$. Si $x\in V\cap\overline H$, toda vecindad de $x$ contenida en $U$ intersecta $H\cap U$. Por tanto $x\in\overline{H\cap U}\subset K$: la última inclusión usa que $K$ es compacto en $G$ y, por ser $G$ Hausdorff, cerrado. Así $V\cap\overline H\subset H$. El subgrupo $H$ contiene una vecindad de la identidad del grupo $\overline H$; sus traslaciones muestran que es abierto en ese grupo. Un subgrupo abierto es también cerrado, porque su complemento es unión de clases laterales abiertas. Como $H$ es denso en $\overline H$, necesariamente $H=\overline H$.
\end{proof}

\begin{lema}[Funciones de corte]\label{v4:haar-corte}\index{funcion@función!de corte}
Si $X$ es localmente compacto y Hausdorff, $K\subset X$ es compacto y $U\supset K$ es abierto, existe $\phi\in\Cc(X)$ tal que $0\leq\phi\leq1$, $\phi=1$ en $K$ y $\supp\phi\subset U$.
\end{lema}
\begin{proof}
Para cada punto de $K$ elijamos una vecindad cuyo cierre sea compacto y esté contenido en $U$. Una subcubierta finita produce un abierto $V$ con $K\subset V\subset\overline V\subset U$ y $\overline V$ compacto. Si $X$ es compacto, apliquemos el lema de Urysohn a los cerrados disjuntos $K$ y $X\setminus V$. Si no lo es, hagámoslo en la compactificación por un punto $X\cup\{\infty\}$, que es compacta y Hausdorff y, por tanto, normal. La función obtenida vale cero fuera de $V$; su restricción a $X$ tiene soporte contenido en $\overline V$. Este argumento utiliza normalidad de un espacio compacto, sin presuponer que todo espacio localmente compacto sea normal.
\end{proof}

\begin{lema}[Continuidad uniforme de las traslaciones]\label{v4:haar-uniforme}
Para $f\in\Cc(G)$ y $\eps>0$ existe una vecindad $V$ de $e$ tal que
\[
 \norm{L_vf-f}_\infty<\eps,\qquad
 \norm{R_vf-f}_\infty<\eps\quad(v\in V).
\]
Además, si $v$ recorre una vecindad relativamente compacta fija de $e$, los soportes de estas funciones están contenidos en un mismo compacto.
\end{lema}
\begin{proof}
Sea $S=\supp f$ y sea $W$ una vecindad compacta y simétrica de $e$. Los soportes de $L_vf$ y $R_vf$, para $v\in W$, están contenidos en $WS$ y $SW$, respectivamente. Fuera del compacto $C=S\cup WS\cup SW$ todas las diferencias consideradas son cero. La función $(v,t)\mapsto f(v^{-1}t)-f(t)$ es continua y se anula en $\{e\}\times C$. Para cada $t\in C$, la continuidad proporciona vecindades $V_t$ de $e$ y $O_t$ de $t$ en las cuales su módulo es menor que $\eps$. Cubrimos $C$ con un número finito de los $O_t$ e intersectamos los correspondientes $V_t$. Esto da la primera desigualdad uniformemente en $t$. El mismo argumento aplicado a $(v,t)\mapsto f(tv)-f(t)$ da la segunda. Intersectamos las dos vecindades obtenidas con el interior de $W$.
\end{proof}

\begin{definicion}[Convención de medida de Haar]\label{v4:haar-def}
Una medida de Haar izquierda es una medida de Borel $\mu$ no nula, finita en compactos, exteriormente regular en los conjuntos de Borel e interiormente regular en los abiertos, tal que
\[
 \mu(xA)=\mu(A)\qquad(x\in G,\ A\in\mathcal B(G)).
\]
La regularidad interior también vale para todo Borel de medida finita. Para las integrales podemos usar la completación de esta medida. Una medida de Haar derecha satisface $\mu(Ax)=\mu(A)$.
\end{definicion}
\begin{observacion}[Una precisión necesaria sobre «Radon»]\label{v4:haar-radon}\index{regularidad!de una medida}
La terminología varía: algunos autores reservan «Radon» para medidas interiormente regulares en todos los Borel. En grupos $\sigma$-compactos las condiciones usuales coinciden; fuera de ellos hay que especificarlas. La definición de las lecturas exige simultáneamente regularidad interior y exterior en todos los Borel. Esa exigencia es demasiado fuerte en la generalidad anunciada.

Por ejemplo, sea $D$ un grupo discreto no numerable y $G=\R\times D$. Normalicemos la medida de Haar exteriormente regular para que cada intervalo de longitud uno en una componente tenga medida uno. El cerrado $E=\{0\}\times D$ tiene intersección de medida cero con cada compacto, porque la proyección de un compacto sobre $D$ es finita. Sin embargo, todo abierto que contiene $E$ tiene medida infinita. En efecto, contiene un intervalo de longitud positiva en cada componente; entre una familia no numerable de longitudes positivas, alguna subfamilia infinita tiene longitudes mayores que $1/n$. Las uniones finitas de esos intervalos tienen medidas arbitrariamente grandes. Por regularidad exterior, $\mu(E)=\infty$. No hay regularidad interior en $E$.

Si $G$ es $\sigma$-compacto, escribimos $G=\bigcup K_n$ con $K_n$ compactos crecientes. Para un Borel $A$, cada $A\cap K_n$ tiene medida finita e interiormente regular; la continuidad creciente de la medida prueba la regularidad interior de $A$. Esta es la situación de $\R^n$, los grupos compactos y los grupos discretos numerables.
\end{observacion}

\begin{proposicion}\label{v4:haar-integral}
Para una medida con la regularidad de la definición anterior, la invariancia izquierda equivale a
\[
 \int_G L_xf\,d\mu=\int_G f\,d\mu
 \qquad(x\in G,\ f\in\Cc^+(G)).
\]
La medida $\nu(A)=\mu(A^{-1})$ es de Haar derecha si y sólo si $\mu$ es de Haar izquierda.
\end{proposicion}
\begin{proof}
La inversión es un homeomorfismo; conserva compactos y abiertos y transforma $(Ax)^{-1}$ en $x^{-1}A^{-1}$. Esto prueba la última afirmación, incluida la regularidad. Si $\mu$ es invariante, la identidad integral vale para indicadores, pues $L_x1_A=1_{xA}$; por linealidad vale para funciones simples y por convergencia monótona para funciones no negativas.

Recíprocamente, para todo abierto $U$ la regularidad interior y el lema de corte dan
\[
 \mu(U)=\sup\left\{\int f\,d\mu: f\in\Cc(G),\ 0\leq f\leq1,
       \supp f\subset U\right\}.                     \tag{10.1}\label{v4:haar-abiertos}
\]
Una dirección es inmediata porque $f\leq1_U$. Para la otra, cada compacto $K\subset U$ admite una función de corte con $\int f\geq\mu(K)$. El argumento también cubre $\mu(U)=\infty$: entonces existen compactos de medida tan grande como se quiera. Las traslaciones establecen una biyección entre las funciones de prueba de $U$ y las de $xU$; la hipótesis integral da $\mu(xU)=\mu(U)$. Tomar ínfimos sobre abiertos que contienen un Borel $A$ prueba $\mu(xA)=\mu(A)$.
\end{proof}

\begin{proposicion}[Positividad]\label{v4:haar-positividad}
Todo abierto no vacío tiene medida de Haar positiva. Si $f\in\Cc^+(G)$ no es cero, entonces $\int f\,d\mu>0$.
\end{proposicion}
\begin{proof}
Como $\mu$ no es nula, $\mu(G)>0$; por regularidad interior del abierto $G$ existe un compacto $K$ con $\mu(K)>0$. Para un abierto no vacío $U$, las traslaciones $xU$ cubren $K$. Una subcubierta finita de $m$ elementos da $0<\mu(K)\leq m\mu(U)$, de modo que $\mu(U)>0$. Si $f(t_0)>0$, elijamos $0<a<f(t_0)$. El abierto no vacío $\{f>a\}$ tiene medida positiva y $\int f\geq a\mu(\{f>a\})>0$.
\end{proof}
\begin{observacion}
No se pueden omitir «no vacío» ni «no negativa». El vacío tiene medida cero, y una función continua no nula puede tener integral cero por cancelación. Estas dos precisiones corrigen el enunciado de positividad de la lectura 17.
\end{observacion}

\section{Construcción mediante recubrimientos}\label{lectura:v4-18}
\begin{estrategia}
Una medida que todavía no conocemos debería comparar dos funciones por el número ponderado de traslaciones necesario para que una cubra a la otra. Normalizamos esa comparación y hacemos que la función que cubre se concentre cerca de la identidad. La compacidad proporciona un límite; la concentración convierte una desigualdad subaditiva en aditividad.
\end{estrategia}

Para $f,\phi\in\Cc^+(G)\setminus\{0\}$ definimos
\[
 (f:\phi)=\inf\left\{\sum_{j=1}^m a_j:
  f\leq\sum_{j=1}^m a_jL_{x_j}\phi,\quad a_j>0,\ x_j\in G\right\}.
\]
\index{recubrimiento!de funciones por traslaciones}
El conjunto del ínfimo no es vacío. En efecto, $\phi>c>0$ en algún abierto no vacío $U$; un número finito de traslaciones de $U$ cubre $\supp f$. Multiplicar la suma de las correspondientes traslaciones de $\phi$ por $\norm f_\infty/c$ produce un mayorante de $f$.

\begin{lema}\label{v4:haar-recubrimientos}
Para funciones no negativas de soporte compacto se cumplen
\begin{align*}
 (L_xf:\phi)&=(f:\phi), & (af:\phi)&=a(f:\phi)\quad(a>0),\\
 (f+g:\phi)&\leq(f:\phi)+(g:\phi), & f\leq g&\Longrightarrow(f:\phi)\leq(g:\phi),\\
 (f:\phi)&\geq\norm f_\infty/\norm\phi_\infty,&
 (f:\phi)&\leq(f:g)(g:\phi).
\end{align*}
\end{lema}
\begin{proof}
Trasladar un recubrimiento cambia $L_{x_j}\phi$ por $L_{xx_j}\phi$ y conserva sus coeficientes; aplicar también $x^{-1}$ da la igualdad. Multiplicar o dividir los coeficientes prueba la homogeneidad. Unir dos recubrimientos prueba la subaditividad, tomando recubrimientos cuyas sumas estén a menos de $\eta$ de los respectivos ínfimos y dejando $\eta\downarrow0$. La monotonía se obtiene usando para $f$ cualquier recubrimiento de $g$. La cota inferior se deduce de $\norm f_\infty\leq\norm\phi_\infty\sum a_j$.

Para la última desigualdad, si $f\leq\sum a_iL_{x_i}g$ y $g\leq\sum b_jL_{y_j}\phi$, entonces
\[
 f\leq\sum_{i,j}a_ib_jL_{x_iy_j}\phi.
\]
La suma de los nuevos coeficientes es $(\sum a_i)(\sum b_j)$. Elegir ambas sumas arbitrariamente próximas a sus ínfimos prueba la afirmación. No se ha supuesto que un ínfimo se alcance.
\end{proof}

Fijemos $f_0\in\Cc^+(G)\setminus\{0\}$ y pongamos
\[
 I_\phi(f)=\frac{(f:\phi)}{(f_0:\phi)}.
\]
El denominador es positivo. La última desigualdad del lema, usada en ambos sentidos, da
\[
 \frac1{(f_0:f)}\leq I_\phi(f)\leq(f:f_0),\qquad I_\phi(f_0)=1.       \tag{10.2}\label{v4:haar-cotas}
\]
Cada $I_\phi$ es homogéneo, monótono, subaditivo e invariante por traslaciones izquierdas.

\begin{lema}[Aditividad aproximada]\label{v4:haar-casiaditivo}
Dados $f,g\in\Cc^+(G)\setminus\{0\}$ y $\eps>0$, existe una vecindad $U$ de $e$ tal que, para \emph{toda} $\phi\in\Cc^+(G)\setminus\{0\}$ con $\supp\phi\subset U$,
\[
 I_\phi(f)+I_\phi(g)\leq I_\phi(f+g)+\eps.
\]
\end{lema}
\begin{proof}
Elijamos $\psi\in\Cc^+(G)$ igual a uno en $S=\supp(f+g)$ y un número $\delta>0$ que fijaremos al final. Sea $h=f+g+\delta\psi$. Las funciones
\[
 h_1=f/h,\qquad h_2=g/h
\]
se definen como cero donde $h=0$. Son continuas: cerca de $S$ el denominador es positivo y fuera de $S$ los numeradores son cero. Tienen soporte compacto y $0\leq h_1+h_2\leq1$. Por el lema \ref{v4:haar-uniforme}, existe $U$ tal que
\[
 |h_i(xu)-h_i(x)|<\delta\quad(x\in G,\ u\in U,\ i=1,2).
\]
Si $h\leq\sum a_jL_{x_j}\phi$ y $L_{x_j}\phi(x)>0$, entonces $x_j^{-1}x\in U$, por lo que $h_i(x)\leq h_i(x_j)+\delta$. Multiplicando el recubrimiento por $h_i(x)$ obtenemos
\[
 (f:\phi)\leq\sum a_j(h_1(x_j)+\delta),\qquad
 (g:\phi)\leq\sum a_j(h_2(x_j)+\delta).
\]
Sumar y aproximar el ínfimo de los recubrimientos de $h$ da
\[
 I_\phi(f)+I_\phi(g)\leq(1+2\delta)I_\phi(h)
 \leq(1+2\delta)\bigl(I_\phi(f+g)+\delta I_\phi(\psi)\bigr).
\]
Si $A=(f+g:f_0)$ y $B=(\psi:f_0)$, la diferencia entre el último miembro e $I_\phi(f+g)$ es a lo más
$2\delta A+\delta(1+2\delta)B$. Elegimos $\delta$ tan pequeño que esta cantidad sea menor que $\eps$. Las cotas son independientes de $\phi$, que es precisamente la uniformidad necesaria para el paso al límite.
\end{proof}

\section{Existencia y unicidad}
\begin{teorema}[Haar]\label{v4:haar-teorema}
Todo grupo localmente compacto y Hausdorff admite una medida de Haar izquierda. Dos de estas medidas son múltiplos positivos una de otra.
\end{teorema}
\begin{proof}[Demostración de la existencia]
Consideremos el producto compacto
\[
 P=\prod_{f\in\Cc^+(G)\setminus\{0\}}
       \left[\frac1{(f_0:f)},(f:f_0)\right].
\]
Su compacidad es el teorema de Tychonoff. Por \eqref{v4:haar-cotas}, cada $I_\phi$ es un punto de $P$. Para cada vecindad $U$ de $e$, sea $K_U$ la clausura en $P$ del conjunto de los $I_\phi$ con $\supp\phi\subset U$. Este conjunto no es vacío, por el lema de corte. Además,
$K_{U_1\cap\cdots\cap U_n}\subset\bigcap K_{U_j}$. Los compactos $K_U$ tienen la propiedad de intersección finita; existe $I\in\bigcap_U K_U$.

Una igualdad o desigualdad que involucra un número finito de coordenadas y es cerrada se conserva al tomar clausuras. Por ello $I$ es homogéneo, monótono, invariante, subaditivo e $I(f_0)=1$. Fijados $f,g$ y $\eps>0$, el lema anterior proporciona $U$ para el cual todos los puntos que generan $K_U$ satisfacen
$I_\phi(f)+I_\phi(g)-I_\phi(f+g)\leq\eps$. La desigualdad define un cerrado del producto; la satisface también $I$. Como $\eps$ es arbitrario, obtenemos la desigualdad inversa a la subaditividad. Así
\[
 I(f+g)=I(f)+I(g).
\]
Declaramos $I(0)=0$ y extendemos a funciones reales mediante $I(u-v)=I(u)-I(v)$ para $u,v\geq0$. La extensión está bien definida: si $u-v=u'-v'$, entonces $u+v'=u'+v$ y la aditividad da la misma diferencia. Extendemos después por linealidad compleja. Si una función real $f$ está soportada en un compacto $K$ y $\psi\in\Cc^+(G)$ vale uno en $K$, la desigualdad $-\norm f_\infty\psi\leq f\leq\norm f_\infty\psi$ muestra que
$|I(f)|\leq I(\psi)\norm f_\infty$. Para funciones complejas, multiplicar por una constante de módulo uno y tomar parte real da la misma cota.

El teorema de representación de Riesz para funcionales positivos sobre $\Cc(G)$, desarrollado en el volumen II, proporciona una medida finita en compactos, exteriormente regular y regular interior en abiertos tal que $I(f)=\int f\,d\mu$. No es nula porque $I(f_0)=1$. La invariancia de $I$ y la proposición \ref{v4:haar-integral} prueban que $\mu$ es de Haar.
\end{proof}

\begin{proof}[Demostración de la unicidad]\label{lectura:v4-19}
Sean $I(f)=\int f\,d\mu$ y $J(f)=\int f\,d\nu$ para dos medidas de Haar izquierdas. Probaremos que $I(f)/J(f)$ tiene el mismo valor para todas las $f\in\Cc^+(G)\setminus\{0\}$. Los denominadores son positivos por la proposición \ref{v4:haar-positividad}.

Fijemos dos funciones de ese tipo $f,g$, una vecindad compacta y simétrica $K$ de $e$ y $\eps>0$. Por la continuidad uniforme, podemos escoger una vecindad simétrica $V\subset K$ tal que, para $t\in V$ y todo $x$,
\[
 |f(tx)-f(xt)|<\eps,\qquad |g(tx)-g(xt)|<\eps.
\]
En efecto, cada diferencia está acotada por la suma de las diferencias respecto de $f(x)$ o $g(x)$. Elijamos $h\in\Cc^+(G)\setminus\{0\}$, simétrica, con soporte en $V$. Se obtiene tomando una función de corte $j$ suficientemente localizada y poniendo $h(t)=j(t)+j(t^{-1})$.

La invariancia izquierda, la simetría de $h$ y Fubini dan
\begin{align*}
 I(f)J(h)&=\int_G\int_G f(yx)h(y)\,d\mu(x)\,d\nu(y),\\
 I(h)J(f)&=\int_G\int_G f(xy)h(y)\,d\mu(x)\,d\nu(y).
\end{align*}
Para comprobar la segunda igualdad, escribimos primero
$I(h)J(f)=\int f(y)\int h(y^{-1}x)\,d\mu(x)\,d\nu(y)$,
reemplazamos $h(y^{-1}x)$ por $h(x^{-1}y)$, cambiamos el orden de integración y hacemos $y=xz$ en la integral respecto de $\nu$. Todos los integrandos están soportados en un producto de compactos y son acotados; por tanto las integrales son absolutamente convergentes. La conmutación de las integrales puede justificarse sin ninguna hipótesis de numerabilidad: en el producto de los dos compactos, Stone--Weierstrass aproxima uniformemente el integrando continuo por sumas finitas de funciones de cada variable por separado. La igualdad es inmediata para esas sumas y pasa al límite por la finitud de ambas medidas.

Para $y\in K$, las funciones $x\mapsto f(yx)$ y $x\mapsto f(xy)$ están soportadas en el compacto
$C_f=K\supp f\cup(\supp f)K$. Al restar las dos identidades obtenemos
\[
 |I(f)J(h)-I(h)J(f)|\leq\eps\mu(C_f)J(h).
\]
Dividiendo por $J(f)J(h)$ resulta
\[
 \left|\frac{I(f)}{J(f)}-\frac{I(h)}{J(h)}\right|
 \leq\eps\frac{\mu(C_f)}{J(f)}.
\]
La misma estimación para $g$, con la misma $h$, muestra que la diferencia entre $I(f)/J(f)$ e $I(g)/J(g)$ es a lo más
$\eps(\mu(C_f)/J(f)+\mu(C_g)/J(g))$. Los compactos no dependen de $\eps$; dejar $\eps\downarrow0$ prueba que las razones coinciden. Por descomposición en partes positivas y negativas, $I=cJ$ en todo $\Cc(G)$ con $c>0$. La fórmula \eqref{v4:haar-abiertos} y la regularidad exterior implican $\mu=c\nu$ en todos los Borel.
\end{proof}

\section{Continuidad en norma y ejercicios de la lectura 17}
\begin{proposicion}\label{v4:haar-Lp}
Si $1\leq p<\infty$ y $f\in L^p(G)$, entonces $L_x$ es una isometría y la aplicación $x\mapsto L_xf$ es continua en la norma de $L^p$.
\end{proposicion}
\begin{proof}
La invariancia de Haar da $\int|L_xf|^p=\int|f|^p$, primero para funciones medibles no negativas y luego para sus clases. Para $f\in\Cc(G)$, los soportes de $L_xf-f$ quedan en un compacto fijo $C$ cuando $x$ está cerca de $e$. Por el lema \ref{v4:haar-uniforme},
$\norm{L_xf-f}_p\leq\mu(C)^{1/p}\norm{L_xf-f}_\infty\to0$.
La densidad de $\Cc(G)$ en $L^p(G)$, consecuencia de la regularidad y las funciones de corte, permite elegir $g\in\Cc(G)$ con $\norm{f-g}_p<\eps$. Entonces
\[
 \norm{L_xf-f}_p\leq2\eps+\norm{L_xg-g}_p.
\]
Esto prueba continuidad en $e$. Como $L_a^{-1}L_x=L_{a^{-1}x}$ y $L_a$ es una isometría, también prueba continuidad en cada $a\in G$.
\end{proof}

\begin{ejercicio}[Lectura 17, ejercicio 1]\label{ej:v4-l17-1}
Sea $\mu$ una medida de Haar izquierda de $G$. Demuestre que $\mu(G)<\infty$ si y sólo si $G$ es compacto.
\end{ejercicio}
\begin{solucion}
Si $G$ es compacto, la definición da $\mu(G)<\infty$. Para la otra dirección, supongamos que $G$ no es compacto. Elijamos un abierto no vacío $U$ de cierre compacto $K$. Construiremos $x_n$ con las traslaciones $x_nU$ disjuntas. Elegido $x_1=e$, si ya elegimos $x_1,\ldots,x_n$, la unión finita $\bigcup_{j=1}^n x_jKK^{-1}$ es compacta y no puede ser todo $G$. Tomemos $x_{n+1}$ fuera de ella. Si $x_{n+1}U\cap x_jU\ne\varnothing$, existirían $u,v\in U$ con $x_{n+1}u=x_jv$, de donde $x_{n+1}=x_jvu^{-1}\in x_jKK^{-1}$, contradicción. La aditividad y la invariancia dan
$\mu(G)\geq\sum_n\mu(x_nU)=\sum_n\mu(U)=\infty$, pues $\mu(U)>0$.

La hipótesis «$\mu$ es de Haar» es esencial y no aparece explícita en el ejercicio original. Una probabilidad gaussiana sobre $\R$ es finita, aunque $\R$ no sea compacto; no es invariante por traslaciones.
\end{solucion}

\begin{ejercicio}[Lectura 17, ejercicio 2]\label{ej:v4-l17-2}
Si $A$ es Borel y $0<\mu(A)<\infty$, pruebe que $a(x)=\mu(A\cap xA)$ es continua.
\end{ejercicio}
\begin{solucion}
Tenemos $1_A\in L^2(G)$ y $L_x1_A=1_{xA}$. Por tanto
$a(x)=\ip{1_A}{L_x1_A}$. Para $x$ cercano a $y$, Cauchy--Schwarz da
\[
 |a(x)-a(y)|\leq\norm{1_A}_2\norm{L_x1_A-L_y1_A}_2\longrightarrow0
\]
por la proposición \ref{v4:haar-Lp}. La fórmula con $x^{-1}$ sugerida en la lectura también vale: trasladar $A\cap x^{-1}A$ por $x$ produce $xA\cap A$.
\end{solucion}

\begin{ejercicio}[Lectura 17, ejercicio 3]\label{ej:v4-l17-3}
Con las mismas hipótesis, pruebe que $K=\{x:\mu(A\cap xA)=\mu(A)\}$ es un subgrupo cerrado.
\end{ejercicio}
\begin{solucion}
La invariancia da $\mu(xA)=\mu(A)$, y de aquí
\[
 \norm{L_x1_A-1_A}_2^2=\mu(A\mathbin\triangle xA)
       =2\mu(A)-2\mu(A\cap xA).
\]
Así $K$ es el estabilizador del vector $1_A$ para la representación $x\mapsto L_x$ en $L^2$. Contiene $e$; si $x,y\in K$, entonces $L_{xy}1_A=L_xL_y1_A=1_A$; si $x\in K$, aplicar $L_{x^{-1}}$ a $L_x1_A=1_A$ da $L_{x^{-1}}1_A=1_A$. Es, pues, un subgrupo. Es cerrado porque la función $x\mapsto\norm{L_x1_A-1_A}_2$ es continua y $K$ es su conjunto de ceros.
\end{solucion}

\begin{ejercicio}[Lectura 17, ejercicio 4]\label{ej:v4-l17-4}
Sea $H$ un subgrupo denso de $G$ y $\mu$ una medida no nula, finita en compactos y con la regularidad de la definición \ref{v4:haar-def}. Si $\mu(hA)=\mu(A)$ para $h\in H$, pruebe que $\mu$ es de Haar izquierda.
\end{ejercicio}
\begin{solucion}
Fijemos $f\in\Cc(G)$. La aplicación $x\mapsto\int L_xf\,d\mu$ es continua: cerca de un punto fijo los soportes quedan en un compacto $C$ y
$|\int(L_xf-L_yf)\,d\mu|\leq\mu(C)\norm{L_xf-L_yf}_\infty$, cantidad que tiende a cero. El conjunto donde esta integral vale $\int f\,d\mu$ es cerrado y contiene al subgrupo denso $H$; es todo $G$. Aplicando la proposición \ref{v4:haar-integral} concluimos la invariancia de la medida. No hemos supuesto que un punto de $G$ sea límite de una sucesión de $H$: la argumentación con un cerrado funciona también en grupos no metrizables.
\end{solucion}

\chapter{Función modular e integración en cocientes}\label{v4:modular}\label{cap:v4-11}
La invariancia por la izquierda no garantiza invariancia por la derecha. La función modular mide esa diferencia. Al integrar en un espacio de clases laterales, la misma función determina si existe una medida compatible con la acción del grupo. Desarrollaremos las lecturas 20 y 21 de Valencia Bucio \cite{valencia4}, manteniendo una sola convención de signos y justificando cuidadosamente el paso de funciones continuas a funciones integrables.

\section{La función modular y los cambios de variable}\label{lectura:v4-20}
Fijemos una medida de Haar izquierda $\mu$ de $G$. Para $a\in G$, la medida $\mu_a(A)=\mu(Aa)$ es de Haar izquierda: la traslación derecha es un homeomorfismo, preserva la regularidad y los compactos, y $\mu_a(xA)=\mu(xAa)=\mu(Aa)$. La unicidad de Haar proporciona un único número $\Delta(a)>0$ tal que
\begin{equation}\label{v4:mod-convencion}
 \mu(Aa)=\Delta(a)\mu(A).
\end{equation}
\index{funcion@función!modular}\index{grupo!unimodular}
\begin{definicion}
La aplicación $\Delta:G\to(0,\infty)$ determinada por \eqref{v4:mod-convencion} es la \emph{función modular}. El grupo es \emph{unimodular} si $\Delta\equiv1$.
\end{definicion}

\begin{proposicion}\label{v4:mod-prop}
La función modular no depende de la normalización de $\mu$ y es un homomorfismo continuo. Para toda función no negativa medible, y también para toda función integrable,
\begin{equation}\label{v4:mod-derecha}
 \int_G f(ta)\,d\mu(t)=\Delta(a)^{-1}\int_G f(t)\,d\mu(t).
\end{equation}
\end{proposicion}
\begin{proof}
Multiplicar $\mu$ por una constante positiva no cambia el factor de \eqref{v4:mod-convencion}. Si $A$ tiene medida finita y positiva, entonces
\[
 \mu(Aab)=\Delta(b)\mu(Aa)=\Delta(b)\Delta(a)\mu(A),
\]
lo que da $\Delta(ab)=\Delta(a)\Delta(b)$. Además, $\Delta(e)=1$ y $\Delta(a^{-1})=\Delta(a)^{-1}$.

Para un indicador, $1_A(ta)=1_{Aa^{-1}}(t)$; por ello su integral es $\mu(Aa^{-1})=\Delta(a)^{-1}\mu(A)$. Pasar a funciones simples y a límites monótonos prueba \eqref{v4:mod-derecha} para funciones no negativas. Al aplicarla a $|f|$, se comprueba primero que $R_af$ es integrable; después se aplica a las partes reales e imaginarias positivas y negativas. Así no se presupone la continuidad de $R_a$ en $L^1$ para demostrar la fórmula que precisamente la justifica.

Para probar continuidad, tomemos $f\in\Cc^+(G)$ con $\int f=\alpha>0$. La función $a\mapsto\int f(ta)\,d\mu(t)$ es continua: cerca de cualquier $a_0$ los soportes de $R_af$ están en un compacto fijo, y la convergencia uniforme de las traslaciones controla la integral de la diferencia. Por \eqref{v4:mod-derecha}, esa función es $\alpha\Delta(a)^{-1}$, positiva. Tomar su recíproco prueba la continuidad de $\Delta$.
\end{proof}
\begin{observacion}[Control de la convención]
El factor de la medida de $Aa$ es $\Delta(a)$; el factor de la integral de $f(ta)$ es su inverso. Confundirlos cambia las fórmulas de inversión y de cocientes. En la prueba de multiplicatividad de la lectura 20 se intercambian accidentalmente traslaciones izquierdas y derechas; en la lectura 21 aparecen factores modulares sin el inverso. Aquí todas las fórmulas se deducen de \eqref{v4:mod-convencion}.
\end{observacion}

\begin{proposicion}\label{v4:mod-Lp}
Si $1\leq p<\infty$, entonces
\[
 \norm{R_af}_p=\Delta(a)^{-1/p}\norm f_p,
\]
y las aplicaciones $a\mapsto R_af$ y $a\mapsto L_af$ son continuas en $L^p(G)$. En $L^\infty$, ambas traslaciones conservan la norma, pero la continuidad en norma puede fallar.
\end{proposicion}
\begin{proof}
Aplicar \eqref{v4:mod-derecha} a $|f|^p$ da la fórmula de la norma. Para $g\in\Cc(G)$, la continuidad uniforme y el compacto común de los soportes dan $\norm{R_ag-g}_p\to0$ cuando $a\to e$. Para $f\in L^p$, elijamos $g\in\Cc(G)$ próximo a $f$. Entonces
\[
 \norm{R_af-f}_p\leq
 (\Delta(a)^{-1/p}+1)\norm{f-g}_p+\norm{R_ag-g}_p.
\]
La continuidad de $\Delta$ ya está probada, de modo que el primer factor está acotado cerca de $e$. Primero elegimos $g$ y después $a$; obtenemos continuidad en $e$. Las identidades de composición y la acotación de cada traslación transportan el resultado a cualquier punto. El caso izquierdo se probó en \ref{v4:haar-Lp}.

Las traslaciones preservan los conjuntos nulos, por lo que conservan la norma esencial suprema. Sin embargo, para $f=1_{[0,\infty)}$ en $\R$ y $a\ne0$, la diferencia $R_af-f$ tiene módulo uno en un intervalo de longitud $|a|$. Por ello $\norm{R_af-f}_\infty=1$ aunque $a\to0$.
\end{proof}

\begin{teorema}[Cambio por inversión]\label{v4:mod-inversion}
Para $f\in L^1(G)$,
\begin{equation}\label{v4:mod-formula-inversion}
 \int_G f(t^{-1})\Delta(t^{-1})\,d\mu(t)=\int_G f(t)\,d\mu(t).
\end{equation}
La igualdad también vale para funciones no negativas, permitiendo el valor infinito.
\end{teorema}
\begin{proof}
Para $f\in\Cc(G)$ escribamos
$Tf(t)=f(t^{-1})\Delta(t^{-1})$ y $J(f)=\int Tf\,d\mu$.
Es un funcional positivo no nulo. Verifiquemos su invariancia izquierda. Si
$F(u)=f(u^{-1})\Delta(u^{-1})$, entonces
\[
 f(a^{-1}t^{-1})\Delta(t^{-1})=\Delta(a)F(ta).
\]
La identidad usa $(ta)^{-1}=a^{-1}t^{-1}$ y la multiplicatividad de $\Delta$. Por \eqref{v4:mod-derecha}, integrar el miembro derecho produce $\int F$, de donde $J(L_af)=J(f)$. Riesz y la unicidad de Haar dan $J=cI$ con $c>0$, donde $I(f)=\int f\,d\mu$.

Ahora $T^2f=f$, porque $\Delta(t)\Delta(t^{-1})=1$. Por tanto
$I(f)=I(T^2f)=cI(Tf)=c^2I(f)$. Elegir $I(f)>0$ muestra $c=1$. Esto prueba \eqref{v4:mod-formula-inversion} en $\Cc(G)$.

La medida $A\mapsto\int 1_A(t^{-1})\Delta(t^{-1})\,d\mu(t)$ es finita en compactos y tiene la misma regularidad que $\mu$: la inversión es un homeomorfismo y el peso es continuo y estrictamente positivo. La igualdad en $\Cc$ implica igualdad de las medidas por regularidad. La aproximación por funciones simples extiende entonces la fórmula a funciones no negativas y a $L^1$.
\end{proof}

\begin{ejemplo}[Grupos unimodulares y un ejemplo no unimodular]\label{v4:mod-ejemplos}
Los grupos abelianos son unimodulares, pues $Aa=aA$. Los grupos compactos también: $\Delta(G)$ es un subgrupo compacto de $(0,\infty)$, y el único tal subgrupo es $\{1\}$. En efecto, si contiene $r>1$, contiene las potencias no acotadas $r^n$; si contiene $0<r<1$, contiene $r^{-1}>1$. Los grupos discretos son unimodulares porque la medida de conteo es invariante por ambos lados.

Para ver que la distinción no es vacía, consideremos el grupo afín
\[
 G=\{(a,b):a>0,\ b\in\R\},\qquad
 (a,b)(c,d)=(ac,b+ad).
\]
La medida $d\mu(a,b)=a^{-2}\,da\,db$ es invariante por la izquierda. La aplicación $(a,b)\mapsto(c,d)(a,b)=(ca,d+cb)$ tiene jacobiano $c^2$, que cancela el factor $(ca)^{-2}$. La traslación derecha $(a,b)\mapsto(ac,b+ad)$ tiene jacobiano $c$, de modo que
$\mu(A(c,d))=c^{-1}\mu(A)$. Por tanto $\Delta(c,d)=c^{-1}$ y el grupo no es unimodular. La medida derecha correspondiente es $a^{-1}\,da\,db$.
\end{ejemplo}

\section{Ejercicios de la lectura 20}
\begin{ejercicio}[Lectura 20, ejercicio 1]\label{ej:v4-l20-1}
Demuestre que, para $f\in\Cc(G)$ y $\eps>0$, existe una vecindad $U$ de $e$ tal que $x^{-1}y\in U$ o $yx^{-1}\in U$ implica $|f(x)-f(y)|<\eps$.
\end{ejercicio}
\begin{solucion}
Por el lema \ref{v4:haar-uniforme}, elijamos una vecindad simétrica $U$ tal que $\norm{R_uf-f}_\infty<\eps$ y $\norm{L_uf-f}_\infty<\eps$ para $u\in U$. Si $u=x^{-1}y$, entonces $y=xu$ y
$|f(y)-f(x)|=|R_uf(x)-f(x)|<\eps$. Si $u=yx^{-1}$, entonces $y=ux$ y
$|f(y)-f(x)|=|L_{u^{-1}}f(x)-f(x)|<\eps$.

La parte esencial del argumento es que la vecindad no depende de $x$. Esto procede de la cubierta finita de un compacto que contiene todos los soportes trasladados, demostrada en el lema citado; la continuidad puntual por sí sola no bastaría.
\end{solucion}

\begin{ejercicio}[Lectura 20, ejercicio 2]\label{ej:v4-l20-2}
Pruebe la equivalencia entre: $G$ es discreto; $\mu(\{e\})>0$; $\mu$ es un múltiplo positivo de la medida de conteo.
\end{ejercicio}
\begin{solucion}
Si $G$ es discreto, los compactos son finitos y la medida de conteo es de Haar. La unicidad da la tercera afirmación. Esta implica evidentemente la segunda. Supongamos ahora que $c=\mu(\{e\})>0$. Por invariancia, cada singleton tiene medida $c$. Una vecindad compacta $K$ de $e$ tiene medida finita; no puede contener más de $\mu(K)/c$ puntos, porque la medida de cualquier conjunto de $m$ puntos es $mc$. Así $K$ es finito. Existe un abierto $V$ con $e\in V\subset K$; al quitar los finitos puntos de $V\setminus\{e\}$, que son cerrados por la propiedad Hausdorff, obtenemos que $\{e\}$ es abierto. Todas sus traslaciones son abiertas, luego $G$ es discreto. La medida de conteo en un conjunto infinito, incluso no numerable, se entiende como infinito sobre todo subconjunto infinito.
\end{solucion}

\begin{ejercicio}[Lectura 20, ejercicio 3]\label{ej:v4-l20-3}
Sea $H$ un subgrupo del grupo multiplicativo $(0,\infty)$. Demuestre que es discreto o denso.
\end{ejercicio}
\begin{solucion}
El logaritmo es un isomorfismo topológico de $(0,\infty)$ sobre $(\R,+)$. Basta estudiar el subgrupo aditivo $A=\log H$. Si $A=\{0\}$, es discreto. En otro caso, $A$ tiene elementos positivos. Sea $d=\inf(A\cap(0,\infty))$.

Si $d=0$, dados $x\in\R$ y $\eps>0$, elijamos $0<a\in A$ con $a<\eps$. Existe un entero $n$ con $|x-na|\leq a/2<\eps$. Como $na\in A$, el grupo $A$ es denso.

Si $d>0$, probaremos que $d\in A$. Elijamos $a\in A$ con $d\leq a<2d$. Si $a>d$, la definición de ínfimo proporciona $b\in A$ con $d\leq b<a$ y entonces $0<a-b<d$, contradicción. Por tanto $a=d$. Para cualquier $x\in A$ hacemos división euclidiana real: $x=nd+r$ con $n\in\Z$ y $0\leq r<d$. Como $r=x-nd\in A$, la definición de $d$ fuerza $r=0$. Así $A=d\Z$, que es discreto. Al aplicar la exponencial obtenemos $H=\{e^{nd}:n\in\Z\}$ en este caso.
\end{solucion}

\section{Promediar sobre un subgrupo}\label{lectura:v4-21}
Sea ahora $H$ un subgrupo cerrado de $G$, con medida de Haar izquierda $\mu_H$. El espacio $Q=G/H$ es Hausdorff; la proyección $\pi:G\to Q$ es abierta porque $\pi^{-1}(\pi(U))=UH$ es unión de traslaciones abiertas. Si $K$ es una vecindad compacta de $x$, entonces $\pi(K)$ es una vecindad compacta de $xH$. Por tanto $Q$ es localmente compacto.

\begin{observacion}[La medida de un subgrupo cerrado]\label{v4:mod-subgrupo-nulo}
Si $H$ es abierto, la restricción de $\mu$ a $H$ es una medida de Haar no nula. Si no es abierto, todo compacto $K\subset H$ tiene medida cero. En efecto, si $\mu(K)>0$, la función $x\mapsto\mu(K\cap xK)$ es continua y positiva en $e$ por el ejercicio \ref{ej:v4-l17-2}. Sería positiva en una vecindad $V$, lo que implica $V\subset KK^{-1}\subset H$ y haría abierto a $H$. En particular, si $H$ es $\sigma$-compacto y no es abierto, entonces $\mu(H)=0$.

Sin $\sigma$-compacidad no se debe afirmar automáticamente $\mu(H)=0$ bajo nuestra convención exteriormente regular. El subgrupo $\{0\}\times D$ del ejemplo \ref{v4:haar-radon} no es abierto y tiene medida infinita. Su restricción no es una medida de Haar de $H$: no es regular interior en el abierto $H$ de su propia topología. La medida intrínseca $\mu_H$ y la restricción de $\mu$ son objetos distintos.
\end{observacion}

Para $f\in\Cc(G)$ definamos
\[
 T_Hf(xH)=\int_H f(xh)\,d\mu_H(h).
\]
\index{integracion@integración!sobre un cociente}
La definición no depende del representante: si se cambia $x$ por $xh_0$, la invariancia izquierda de $\mu_H$ permite sustituir $h_0h$ por $h$. La función de $h$ integrada tiene soporte contenido en el compacto $x^{-1}\supp f\cap H$.

\begin{lema}\label{v4:mod-promedio}
El operador $T_H:\Cc(G)\to\Cc(Q)$ es lineal, positivo y sobreyectivo. Más precisamente,
$\supp(T_Hf)\subset\pi(\supp f)$, y toda $g\in\Cc^+(Q)$ tiene una preimagen no negativa.
\end{lema}
\begin{proof}
La linealidad y positividad se siguen de la integral. Fuera del compacto $\pi(\supp f)$ todas las fibras evitan $\supp f$, luego el promedio es cero. Para la continuidad, fijemos $x\in G$ y una vecindad compacta $U$ de $x$. Los soportes de $h\mapsto f(yh)$ para $y\in U$ están en
$C=U^{-1}\supp f\cap H$, compacto. Cuando $y\to x$, la continuidad uniforme por traslación izquierda da
\[
 |T_Hf(yH)-T_Hf(xH)|\leq
 \mu_H(C)\sup_{h\in H}|f(yh)-f(xh)|\longrightarrow0.
\]
Así $T_Hf\circ\pi$ es continua; por la definición de topología cociente, $T_Hf$ es continua en $Q$.

Para la sobreyectividad, sea $S=\supp g$. Cubramos $S$ con un número finito de abiertos $\pi(U_j)$ donde cada $U_j$ es relativamente compacto en $G$. El compacto $A=\bigcup_j\overline{U_j}$ satisface $S\subset\pi(A)$. Tomemos $\phi\in\Cc^+(G)$ igual a uno en $A$. Para $q\in S$, existe $x\in A$ con $q=xH$; la función $h\mapsto\phi(xh)$ es positiva cerca de $e$, por lo que $T_H\phi(q)>0$. Como $S$ es compacto, $T_H\phi$ tiene un mínimo estrictamente positivo en $S$.

Definamos
\[
 f(x)=\begin{cases}
 g(\pi(x))\phi(x)/T_H\phi(\pi(x)),&T_H\phi(\pi(x))>0,\\
 0,&T_H\phi(\pi(x))=0.
 \end{cases}
\]
Esta función es continua: cerca de $\pi^{-1}(S)$ el denominador no se anula, y fuera de ese cerrado el numerador $g\circ\pi$ es cero. Su soporte está contenido en $\supp\phi$, y la constancia de $g\circ\pi$ y $T_H\phi\circ\pi$ en cada fibra da $T_Hf=g$. Si $g\geq0$, también $f\geq0$.
\end{proof}

\section{Cuándo existe una medida invariante en el cociente}
\begin{teorema}[Fórmula integral de Weil]\label{v4:mod-weil}\index{formula@fórmula!integral de Weil}
Existe una medida no nula en $Q=G/H$, finita en compactos, con la regularidad de Haar e invariante por la acción izquierda de $G$, si y sólo si
\begin{equation}\label{v4:mod-compatibilidad}
 \Delta_G(h)=\Delta_H(h)\qquad(h\in H).
\end{equation}
En ese caso es única salvo un múltiplo positivo. Fijadas $\mu_G$ y $\mu_H$, tiene una única normalización $\nu$ para la cual
\begin{equation}\label{v4:mod-weil-formula}
 \int_G f\,d\mu_G=\int_{G/H}T_Hf\,d\nu
   =\int_{G/H}\int_H f(xh)\,d\mu_H(h)\,d\nu(xH)
\end{equation}
para $f\in\Cc(G)$ y para $f\in L^1(G)$. En este último caso las integrales sobre las fibras existen para casi toda clase lateral, y la función resultante se considera en la completación de $\nu$.
\end{teorema}
\begin{proof}[Necesidad y unicidad]
Si existe $\nu$, el funcional $J(f)=\int_Q T_Hf\,d\nu$ es positivo. No es nulo, porque la sobreyectividad positiva de $T_H$ permite obtener una función no negativa cuya integral sea positiva. Además,
$T_H(L_af)(xH)=T_Hf(a^{-1}xH)$; la invariancia de $\nu$ da $J(L_af)=J(f)$. Por Riesz y unicidad de Haar, $J(f)=c\int f\,d\mu_G$ con $c>0$.

Para $h\in H$, la traslación derecha dentro de $H$ da
$T_H(R_hf)=\Delta_H(h)^{-1}T_Hf$. La misma integral en $G$ satisface
$\int R_hf=\Delta_G(h)^{-1}\int f$. Elegir $f$ de integral no nula y comparar ambas expresiones de $J(R_hf)$ fuerza \eqref{v4:mod-compatibilidad}.

Si $\nu_1$ y $\nu_2$ son invariantes, sus funcionales elevados a $\Cc(G)$ son proporcionales por unicidad de Haar. Como $T_H$ es sobreyectivo, las integrales respecto de $\nu_1,\nu_2$ son proporcionales en todo $\Cc(Q)$; la regularidad extiende la igualdad a medidas. Reescalar $\nu$ fija el factor $c$ igual a uno de manera única.
\end{proof}

\begin{proof}[Suficiencia y fórmula en funciones continuas]
Supongamos \eqref{v4:mod-compatibilidad}. El paso decisivo es demostrar
\begin{equation}\label{v4:mod-kernel}
 T_Hf=0\quad\Longrightarrow\quad\int_G f\,d\mu_G=0.
\end{equation}
Para $f,\phi\in\Cc(G)$ y $T_Hf=0$, Fubini y el cambio por traslación derecha dan
\begin{align*}
 0&=\int_G\phi(x)\int_H f(xh)\,dh\,dx\\
  &=\int_H\Delta_G(h)^{-1}\int_G f(u)\phi(uh^{-1})\,du\,dh\\
  &=\int_G f(u)\int_H\Delta_H(h)^{-1}\phi(uh^{-1})\,dh\,du\\
  &=\int_G f(u)T_H\phi(uH)\,du.
\end{align*}
La última igualdad es la fórmula de inversión en $H$. Las aplicaciones de Fubini son legítimas: en el integrando inicial $x\in\supp\phi$ y $h\in(\supp\phi)^{-1}\supp f\cap H$, dos compactos; el integrando es acotado. La conmutación de las integrales también se obtiene aproximando uniformemente el integrando, sobre ese producto compacto, por sumas de productos de funciones continuas de cada variable. Los cambios de variable conservan la integrabilidad.

Elijamos $g\in\Cc^+(Q)$ igual a uno en $\pi(\supp f)$ y una preimagen $\phi\geq0$ con $T_H\phi=g$. La última integral es entonces $\int f$, lo que prueba \eqref{v4:mod-kernel}.

Para $g\in\Cc(Q)$, elijamos $f\in\Cc(G)$ con $T_Hf=g$ y definamos
$\Lambda(g)=\int_G f\,d\mu_G$. La implicación \eqref{v4:mod-kernel} muestra que la definición no depende de la preimagen. Es lineal y positiva porque toda $g\geq0$ admite preimagen positiva. No es cero: si $f\geq0$ no es nula, entonces $\Lambda(T_Hf)=\int f>0$. El teorema de Riesz proporciona la medida $\nu$ y la igualdad \eqref{v4:mod-weil-formula} en $\Cc(G)$. La identidad para traslaciones prueba que $\nu$ es $G$-invariante, usando nuevamente la determinación de una medida por sus integrales de funciones de soporte compacto.
\end{proof}

\subsection{Por qué la fórmula sigue siendo válida en \texorpdfstring{$L^1$}{L1}}
El paso siguiente evita dos afirmaciones incorrectas que aparecen en el argumento original: una función integrable no tiene necesariamente soporte compacto ni soporte topológico $\sigma$-compacto; y un conjunto nulo de $G$ no puede descartarse sobre cada fibra sin una justificación.

\begin{lema}[Los nulos son nulos en casi toda fibra]\label{v4:mod-nulos}
Si $N$ es un Borel de medida $\mu_G(N)=0$, entonces
$\mu_H(\{h:xh\in N\})=0$ para $\nu$-casi toda $xH$.
\end{lema}
\begin{proof}
Primero consideremos un abierto $U\subset G$. Las funciones
$T_H\phi$, con $\phi\in\Cc(G)$, $0\leq\phi\leq1$ y $\supp\phi\subset U$, forman una familia dirigida hacia arriba: $\max(\phi_1,\phi_2)$ domina ambas. Su supremo puntual es $T_H1_U$. Para comprobarlo en una fibra, aproximamos el abierto $x^{-1}U\cap H$ por compactos de $H$; para cada compacto $K$ de ese abierto existe una función de corte en $G$ igual a uno en $xK$ y soportada en $U$. Por regularidad interior de $\mu_H$, las integrales de tales funciones se acercan a $\mu_H(x^{-1}U\cap H)$, incluso si es infinita. Por tanto $T_H1_U$ es semicontinua inferior y medible.

Para una familia dirigida de funciones continuas no negativas sobre un espacio con medida regular interior en abiertos, la integral del supremo es el supremo de las integrales. He aquí la justificación que permite usar una familia no numerable. Los abiertos de nivel del supremo son uniones de los abiertos de nivel de los miembros de la familia. Todo compacto dentro de una de esas uniones está cubierto por un número finito de ellos; la dirección de la familia permite dominarlos con un solo miembro. La regularidad interior da la misma igualdad para las medidas de los abiertos de nivel. Finalmente, una aproximación de la integral por sumas inferiores con un número finito de niveles permite escoger un único miembro que aproxime simultáneamente todos los niveles; tomar el supremo de esas sumas prueba la afirmación.

Aplicamos este hecho y la fórmula de Weil en $\Cc(G)$:
\[
 \int_Q T_H1_U\,d\nu
 =\sup_\phi\int_Q T_H\phi\,d\nu
 =\sup_\phi\int_G\phi\,d\mu_G=\mu_G(U).
\]
Por regularidad exterior, elijamos abiertos $U_n\supset N$ con $\mu_G(U_n)<2^{-n}$. La suma $\sum_nT_H1_{U_n}$ tiene integral finita por convergencia monótona, así que es finita casi en todas partes. En esas clases laterales, $T_H1_{U_n}\to0$. Como $0\leq T_H1_N\leq T_H1_{U_n}$ para todo $n$, concluimos $T_H1_N=0$ allí. La función $T_H1_N$, aun antes de probar su medibilidad global, está soportada en un conjunto $\nu$-nulo y se puede interpretar en la medida completada.
\end{proof}

\begin{proof}[Final de la demostración de Weil en $L^1$]
Para $\phi\in\Cc(G)$,
\[
 \norm{T_H\phi}_{L^1(Q)}\leq\int_QT_H|\phi|\,d\nu
       =\norm\phi_{L^1(G)}.
\]
Por densidad, $T_H$ se extiende a una contracción de $L^1(G)$ en $L^1(Q)$. Falta identificarla con la integral sobre las fibras, y ese paso es el que no debe omitirse.

Elijamos $\phi_n\in\Cc(G)$ con $\norm{\phi_n-f}_1<2^{-n}$. La suma de esas normas es finita; por convergencia monótona, $\phi_n\to f$ fuera de un conjunto Borel nulo $N$, después de elegir un representante Borel de $f$. Además,
\[
 \sum_n\norm{\phi_{n+1}-\phi_n}_1<\infty.
\]
La fórmula en funciones continuas y la convergencia monótona implican que, para casi toda $xH$,
\[
 \int_H\left(|\phi_1(xh)|+
          \sum_n|\phi_{n+1}(xh)-\phi_n(xh)|\right)dh<\infty.
\]
Por el lema anterior, en casi toda fibra la convergencia $\phi_n(xh)\to f(xh)$ vale para casi todo $h$. En la intersección de ambos conjuntos de clases laterales, la expresión entre paréntesis es un mayorante integrable de todos los $|\phi_n(xh)|$ y de $|f(xh)|$. La convergencia dominada da
\[
 T_H\phi_n(xH)\longrightarrow\int_H f(xh)\,dh.
\]
Por otro lado, las contracciones $T_H\phi_n$ convergen en $L^1(Q)$ a la extensión de $T_Hf$. Una subsucesión converge casi en todas partes al mismo límite, de modo que este coincide con la integral sobre la fibra. Finalmente, pasamos al límite en $\int_G\phi_n=\int_QT_H\phi_n$ usando convergencia en norma $L^1$ en ambos lados. Esto prueba \eqref{v4:mod-weil-formula}. El lema de los nulos muestra también que el resultado no depende del representante de $f$.
\end{proof}

\begin{corolario}\label{v4:mod-weil-positivo}
Si $G$ es $\sigma$-compacto, la fórmula de Weil vale para toda función Borel $f\geq0$, permitiendo el valor infinito. Lo mismo vale si $f$ se anula fuera de una unión numerable de compactos.
\end{corolario}
\begin{proof}
Elijamos compactos crecientes $K_n$ que cubran el conjunto donde $f$ puede ser no nula y pongamos $f_n=\min(f,n)1_{K_n}$. Cada $f_n$ es integrable y $f_n\uparrow f$. Aplicamos la fórmula recién demostrada a $f_n$ y usamos convergencia monótona primero en $H$ y después en $Q$. Los conjuntos excepcionales de las fórmulas para $f_n$ tienen unión nula porque son numerables.
\end{proof}

\section{Ejercicios de la lectura 21 y alcance de las hipótesis}
\begin{ejercicio}[Lectura 21, ejercicio 1; versión precisa]\label{ej:v4-l21-1}
Suponga $\Delta_G|_H=\Delta_H$. Si $f:G\to\C$ es Borel y se anula fuera de una unión numerable de compactos, pruebe que
\[
 \int_Q\int_H|f(xh)|\,dh\,d\nu(xH)<\infty
 \quad\Longrightarrow\quad f\in L^1(G).
\]
Explique por qué las hipótesis requieren cuidado en el enunciado original.
\end{ejercicio}
\begin{solucion}
El corolario \ref{v4:mod-weil-positivo}, aplicado a $|f|$, identifica la integral iterada con $\int_G|f|$. La finitud de la primera implica la de la segunda, que es exactamente $f\in L^1(G)$. En particular, la afirmación vale para toda función Borel cuando $G$ es $\sigma$-compacto.

El original no exige medibilidad. Esta no se deduce de la medibilidad de $|f|$: en el grupo compacto $\T$, tomemos un conjunto no medible $A$ y $f=1$ en $A$, $f=-1$ fuera. Entonces $|f|=1$, de integral finita, pero $f$ no es medible. Podemos tomar $H=G$ para que la integral iterada del módulo esté perfectamente definida.

Incluso añadiendo medibilidad, en grupos generales y con la convención exteriormente regular aparece otra dificultad real. Sea $G=\R\times D$ con $D$ discreto no numerable, $H=\R\times\{e\}$ y $f=1_{\{0\}\times D}$. Ambas funciones modulares son uno y $G/H$ se identifica con $D$, con medida de conteo. Sobre cada fibra, $f$ es el indicador de un punto de $\R$, así que la integral interior es cero y la integral iterada también. Pero la observación \ref{v4:haar-radon} demuestra que $\int_Gf=\infty$. El conjunto donde $f$ es no nula no está contenido en una unión numerable de compactos. El contraejemplo muestra por qué no se extiende sin más la fórmula $L^1$ a todos los Borel positivos de un grupo no $\sigma$-compacto. La versión demostrada conserva el contenido útil del ejercicio y declara la hipótesis adicional.
\end{solucion}

\begin{ejercicio}[Lectura 21, ejercicio 2]\label{ej:v4-l21-2}
Si $H$ es cerrado y normal en $G$, demuestre que $\Delta_G|_H=\Delta_H$. Deduzca que $\ker\Delta_G$ es unimodular.
\end{ejercicio}
\begin{solucion}
Cuando $H$ es normal, $G/H$ es un grupo localmente compacto Hausdorff. Su medida de Haar izquierda es invariante bajo la acción izquierda de $G$, pues esta acción es precisamente la traslación por la clase $xH$. La necesidad en el teorema \ref{v4:mod-weil} da $\Delta_G|_H=\Delta_H$. No estamos suponiendo la fórmula cociente para obtener la medida: usamos la existencia independiente de Haar en el grupo $G/H$.

El núcleo $N=\ker\Delta_G$ es cerrado porque $\Delta_G$ es continua y $\{1\}$ es cerrado, y es normal porque es el núcleo de un homomorfismo. Aplicando lo anterior a $N$, obtenemos $\Delta_N=\Delta_G|_N=1$. Esta es la unimodularidad de $N$, aunque $G$ no sea unimodular.
\end{solucion}

\begin{ejercicio}[Lectura 21, ejercicio 3]\label{ej:v4-l21-3}
Sean $H$ cerrado y $K$ compacto, subgrupos de $G$, con $G=HK$. Demuestre que se pueden normalizar sus medidas de Haar de modo que
\[
 \int_Gf(g)\,dg=\int_H\int_K f(hk)\,dk\,dh
 \qquad(f\in L^1(G)).
\]
\end{ejercicio}
\begin{solucion}
Normalicemos la medida de $K$ como probabilidad y pongamos $P=H\times K$. El grupo $P$ actúa transitivamente en $G$ por $(h,k)\cdot g=hgk^{-1}$: la hipótesis $G=HK$ permite escribir cualquier elemento como $hk^{-1}$. El estabilizador de $e$ es
$L=\{(\ell,\ell):\ell\in H\cap K\}$, que es compacto. La aplicación $P/L\to G$, $(h,k)L\mapsto hk^{-1}$, es una biyección continua. Es un homeomorfismo: el mapa $P\to G$ es propio, pues la preimagen de un compacto $C$ es cerrada y está contenida en el compacto $(CK\cap H)\times K$; al descender al cociente, una biyección continua propia entre espacios localmente compactos Hausdorff es un homeomorfismo.

La función modular de un grupo es uno en todo subgrupo compacto. Así $\Delta_P|_L=1=\Delta_L$, y el teorema de Weil proporciona una medida $P$-invariante en $P/L$. Normalicemos también la medida de $L$ como probabilidad. Para $f\in\Cc(G)$, la función $F(h,k)=f(hk^{-1})$ tiene soporte compacto en $P$ por la propiedad anterior; es constante en las clases de $L$, de modo que $T_LF$ se identifica con $f$. La fórmula de Weil en $P$ da
\[
 \int_{P/L}f\,d\eta=\int_H\int_K f(hk^{-1})\,dk\,dh
                  =\int_H\int_K f(hk)\,dk\,dh,
\]
porque la medida de Haar de un grupo compacto es invariante por inversión.

La medida de Haar de $G$ también es $P$-invariante: es invariante bajo la izquierda por $H$, y bajo la derecha por $K$ porque $\Delta_G|_K=1$. La unicidad de la medida invariante en $P/L$ muestra que, después de reescalar la medida de $G$, esta coincide con $\eta$. Obtenemos la igualdad deseada en $\Cc(G)$. Para precisar la extensión a $L^1$, la imagen de la medida de Haar de $P$ bajo el mapa propio $p(h,k)=hk^{-1}$ es una medida finita en compactos y con la regularidad adoptada. La regularidad exterior se conserva porque $p$ es cerrado: si un abierto $U$ contiene $p^{-1}(A)$, el abierto $G\setminus p(P\setminus U)$ contiene $A$ y su preimagen está en $U$. La regularidad interior en abiertos se obtiene proyectando compactos. La igualdad ya probada en $\Cc(G)$ identifica esta medida imagen con $\mu_G$. Por aproximación simple, $f\circ p\in L^1(P)$ y su integral es $\int_Gf$. Finalmente aplicamos Weil en $P$ al subgrupo normal y compacto $\{e\}\times K$; el cociente es $H$ y la normalización de Haar da la integral iterada escrita arriba. La inversión en $K$ no altera la integral ni la integrabilidad.
\end{solucion}

\parte{Convolución y representaciones}{El álgebra de convolución codifica las traslaciones del grupo. Las representaciones unitarias se integran para producir representaciones del álgebra, y la construcción recíproca explica el papel de las aproximaciones de la identidad.}
\chapter{El álgebra de convolución de un grupo}\label{cap:v4-12}
La medida de Haar permite repetir la definición de convolución sin coordenadas. La ausencia de conmutatividad distingue el orden de los factores; la posible diferencia entre medidas izquierda y derecha introduce la función modular. Mantendremos visibles ambas cuestiones en las pruebas y en los ejercicios.

\section{Convolución e involución}\label{lectura:v4-22}
\index{convolucion@convolución!en grupos}\index{algebra@álgebra!L1 de un grupo}
Sea $G$ localmente compacto y Hausdorff, con medida de Haar izquierda $\mu$. Escribimos $dx=d\mu(x)$, y fijamos $\int h(xa)dx=\Delta(a)^{-1}\int h(x)dx$. Las clases $L^p$ se toman respecto de la completación de Haar. Para $f,g\in L^1(G)$ ponemos
\begin{equation}\label{v4-l24-d:algebra}
 (f*g)(x)=\int_Gf(y)g(y^{-1}x)dy,\qquad
 f^*(x)=\Delta(x^{-1})\overline{f(x^{-1})}.
\end{equation}
Recordamos la fórmula de inversión de Haar:
\begin{equation}\label{v4-l24-eq:inv}
 \int_Gh(x^{-1})\Delta(x^{-1})dx=\int_Gh(x)dx.
\end{equation}
En un grupo abeliano $\Delta=1$. En $\R^d$ recuperamos la convolución aditiva; en $\T$ usamos Haar de masa uno; en un grupo discreto con medida de conteo obtenemos $(a*b)(x)=\sum_y a(y)b(y^{-1}x)$. La suma puede tener índices no numerables, pero una función de $\ell^1$ tiene soporte numerable: para cada $n$ solo hay finitos términos de módulo al menos $1/n$, y la unión de esos conjuntos contiene todos los términos no nulos.

\begin{teorema}\label{v4:conv-grupo}\label{v4-l24-p:banach}
La convolución existe casi en todas partes y hace de $L^1(G)$ un álgebra de Banach con involución isométrica. En particular,
\[
 \norm{f*g}_1\leq\norm f_1\norm g_1,\quad
 \norm{f^*}_1=\norm f_1,\quad
 (f*g)^*=g^**f^*.
\]
\end{teorema}
\begin{proof}
Primero precisemos la cuestión de medida para grupos que no son $\sigma$-compactos. Cada función de $L^1$ tiene, salvo un nulo, soporte contenido en una unión numerable de compactos: se aproxima en $L^1$ por funciones de $C_c(G)$ y se elige una subsucesión con errores sumables. Una vecindad relativamente compacta y esa colección numerable de compactos generan un subgrupo abierto $\sigma$-compacto. En él la medida es $\sigma$-finita y se pueden aplicar Tonelli y Fubini ordinarios. Fuera de ese subgrupo las funciones elegidas y su convolución se anulan. Esto no afirma que el soporte topológico de todo representante sea $\sigma$-compacto.

Para representantes borelianos, $(y,x)\mapsto f(y)g(y^{-1}x)$ es medible. La transformación $(y,x)\mapsto(y,y^{-1}x)$ preserva la medida producto en el subgrupo anterior: integrar primero en $x$ y usar invariancia izquierda prueba la identidad para indicadoras y después para no negativas. Por ello sus preimágenes de nulos son nulas, y el argumento se extiende a representantes de la medida completada. Tonelli da
\[
 \int_G\int_G|f(y)g(y^{-1}x)|dydx
 =\int_G|f(y)|\int_G|g(z)|dzdy=\norm f_1\norm g_1.
\]
Fubini y la desigualdad triangular producen existencia, integrabilidad y la cota de norma. También prueban independencia de representantes y $\int f*g=(\int f)(\int g)$.

Para asociatividad, la integral triple absoluta es $\norm f_1\norm g_1\norm h_1$. El cambio $y=zt$ y Fubini dan
\[
 ((f*g)*h)(x)=\int\int f(z)g(z^{-1}y)h(y^{-1}x)dzdy
 =\int f(z)\int g(t)h(t^{-1}z^{-1}x)dtdz.
\]
La última expresión es $f*(g*h)(x)$ casi en todas partes. La bilinealidad es inmediata y la completitud es la de $L^1$.

La fórmula~\eqref{v4-l24-eq:inv} aplicada a $|f|$ prueba $\norm{f^*}_1=\norm f_1$. La antilinealidad y $f^{**}=f$ se obtienen de la conjugación y $\Delta(x^{-1})\Delta(x)=1$. Finalmente,
\begin{align*}
 (g^**f^*)(x)
 &=\int\Delta(y^{-1})\overline{g(y^{-1})}\Delta(x^{-1}y)\overline{f(x^{-1}y)}dy\\
 &=\Delta(x^{-1})\int\overline{g(y^{-1})}\,\overline{f(x^{-1}y)}dy\\
 &=\Delta(x^{-1})\int\overline{f(z)}\,\overline{g(z^{-1}x^{-1})}dz=(f*g)^*(x),
\end{align*}
donde en la última línea usamos $y=xz$. Las identidades se interpretan casi en todas partes.
\end{proof}

\section{Traslaciones y aproximaciones de la identidad}
\index{identidad!aproximada en un grupo}\index{red!de Dirac}
Definimos $L_af(x)=f(a^{-1}x)$ y $R_af(x)=f(xa)$. Un cambio de variable da
\[
 L_a(f*g)=(L_af)*g,\qquad R_a(f*g)=f*(R_ag),\qquad
 \norm{L_af}_1=\norm f_1,\quad\norm{R_af}_1=\Delta(a)^{-1}\norm f_1.
\]
\begin{lema}\label{v4-l24-l:translation}
Para cada $f\in L^1(G)$, $a\mapsto L_af$ es continua en norma. También lo es $a\mapsto\widetilde R_af$, donde $\widetilde R_af(x)=\Delta(a)^{-1}f(xa^{-1})$; estos últimos operadores son isometrías de $L^1$.
\end{lema}
\begin{proof}
Para $f\in C_c$ y $a$ en una vecindad compacta de $e$, los soportes de $f$ y de sus traslaciones quedan dentro de un compacto fijo. La continuidad conjunta, seguida de un recubrimiento finito de ese compacto, muestra que las diferencias tienden uniformemente a cero. Multiplicar la cota uniforme por la medida finita del compacto da continuidad en $L^1$. Para $f$ general, aproximamos por $g\in C_c$ y usamos
\[
 \norm{L_af-f}_1\leq2\norm{f-g}_1+\norm{L_ag-g}_1.
\]
En un punto distinto de $e$ se usa la ley de las traslaciones. Para $\widetilde R_a$, el mismo argumento usa continuidad de $\Delta$ y
$\norm{\widetilde R_af}_1=\Delta(a)^{-1}\Delta(a)\norm f_1=\norm f_1$.
\end{proof}
\begin{definicion}
Una red de Dirac es una red $(k_\alpha)\subset L^1(G)$ con $k_\alpha\geq0$, $\int k_\alpha=1$ y, para \emph{cada} vecindad $U$ de $e$,
\[
 \int_{G\setminus U}k_\alpha(x)dx\longrightarrow0.
\]
La condición de la lectura 22 que permite elegir alguna vecindad $U$ para cada $\eps$ es insuficiente: tomar $U=G$ la haría vacía de contenido.
\end{definicion}
\begin{teorema}\label{v4:dirac-grupo}
Para una red de Dirac y $f\in L^1(G)$, $k_\alpha*f\to f$ y $f*k_\alpha\to f$ en $L^1$.
\end{teorema}
\begin{proof}
Por Tonelli,
\[
 \norm{k_\alpha*f-f}_1\leq\int k_\alpha(y)\norm{L_yf-f}_1dy.
\]
La norma de la diferencia es pequeña en una vecindad de $e$ y está acotada por $2\norm f_1$ en todo $G$. La concentración de masa da el límite. Para el otro lado, el cambio $z=y^{-1}x$ y la fórmula de inversión dan
\[
 f*k_\alpha(x)=\int k_\alpha(z)\Delta(z)^{-1}f(xz^{-1})dz
 =\int k_\alpha(z)\widetilde R_zf(x)dz.
\]
La misma estimación con $\widetilde R_z$ prueba el segundo límite. El factor modular no puede omitirse en este paso.
\end{proof}
\begin{lema}\label{v4-l24-l:approx}
Existe una red $(e_U)\subset C_c(G)$, indexada por vecindades simétricas relativamente compactas de $e$ ordenadas por inclusión inversa, tal que
\[
 e_U\geq0,\quad\supp e_U\subset U,\quad e_U^*=e_U,
 \quad\int e_U=\norm{e_U}_1=1.
\]
Es una identidad aproximada por ambos lados.
\end{lema}
\begin{proof}
La local compacidad y la forma local de Urysohn proporcionan $h\in C_c(G)$ no negativa, no nula y apoyada en $U$. Su integral $a$ es positiva porque Haar da medida positiva a todo abierto no vacío. Sea $e_U=(h+h^*)/(2a)$. La involución preserva positividad e integral por~\eqref{v4-l24-eq:inv}, y $\supp h^*\subset U^{-1}=U$. Así se verifican las propiedades. Si $U$ se hace más pequeña que una vecindad fijada, el soporte queda dentro de ella, de modo que es una red de Dirac. El teorema anterior concluye.
Además se obtiene la cota útil
\begin{equation}\label{v4-l24-eq:approxbound}
 \norm{e_U*f-f}_1\leq\sup_{y\in U}\norm{L_yf-f}_1.
\end{equation}
\end{proof}
\begin{teorema}
El álgebra $L^1(G)$ es conmutativa si y solo si el grupo $G$ es abeliano.
\end{teorema}
\begin{proof}
Si $G$ es abeliano, $\Delta=1$ y el cambio $z=y^{-1}x$ transforma $f*g$ en $g*f$. Recíprocamente, sea conmutativa el álgebra. Para cada $a\in G$, los operadores de convolución por $L_ae_U$ son contracciones en $L^1$ y convergen fuertemente a $L_a$, porque $(L_ae_U)*f=L_a(e_U*f)\to L_af$. Todos esos operadores conmutan por hipótesis. Pasar a límites fuertes, primero en uno de los índices y después en el otro, usando sus cotas uniformes, da $L_aL_b=L_bL_a$, es decir, $L_{ab}=L_{ba}$.

La acción izquierda sobre $L^1$ es fiel: si $z\ne w$, podemos escoger $f\in C_c$ con $f(e)=1$ y $f(w^{-1}z)=0$ por separación local. En el punto $z$, $L_zf=1$ y $L_wf=0$; por continuidad difieren en un abierto de medida positiva. Por tanto $L_z\ne L_w$. Aplicar esto a $z=ab,w=ba$ prueba que $ab=ba$.
\end{proof}

\section{Ejercicios: modularidad y convolución}
\begin{ejercicio}[Lectura 22, ejercicio 1: Young con factor modular]\label{ej:v4-l22-1}
Sean $1/p+1/q=1+1/r$, con exponentes en $[1,\infty]$. Probar la versión válida en un grupo general:
\[
 \norm{f*g}_r\leq\norm f_p\norm{\Delta^{-1/p'}g}_q.
\]
En un grupo unimodular se reduce a la desigualdad usual. Explicar por qué la versión sin peso no vale siempre.
\end{ejercicio}
\begin{solucion}
Pongamos $h=\Delta^{-1/p'}g$. Para $r<\infty$ escribimos, con $z=y^{-1}x$,
\[
 |f(y)g(z)|=(|f(y)|^p|h(z)|^q)^{1/r}
 |f(y)|^{1-p/r}|h(z)|^{1-q/r}\Delta(z)^{1/p'}.
\]
Aplicamos Hölder con recíprocos $1/r$, $1/p-1/r$ y $1/q-1/r=1/p'$. Para el tercer factor, el cambio $y=xz^{-1}$ da
\[
 \int |h(y^{-1}x)|^q\Delta(y^{-1}x)dy=\int|h(z)|^qdz.
\]
Así
\[
 |f*g(x)|^r\leq\norm f_p^{r-p}\norm h_q^{r-q}(|f|^p*|h|^q)(x).
\]
El teorema~\ref{v4:conv-grupo}, seguido de integración, da la cota. Los factores de potencia cero se omiten. El caso $p=1$ se obtiene directamente por Jensen para $|f|/\norm f_1$; el caso $r=\infty$ es Hölder con el mismo cambio de variable y $q=p'$. De este modo se cubren los extremos sin dividir entre exponentes nulos.

El peso es real, no un artificio de la prueba. Para $1<p<\infty$ y $a$ con $\Delta(a)<1$, sean $g_U=L_ae_U$, de norma $L^1$ uno. Para $f\in C_c$ no nula, el cambio de variable anterior y continuidad de traslaciones muestran
\[
 f*g_U\longrightarrow\Delta(a)^{-1}f(\cdot a^{-1})\quad\text{en }L^p.
\]
Su norma límite es $\Delta(a)^{-1/p'}\norm f_p>\norm f_p$. La desigualdad sin peso, con $q=1,r=p$, fallaría para $U$ suficientemente pequeña. La convergencia se verifica expresando la integral como promedio de las funciones $\Delta(au)^{-1}f(\cdot u^{-1}a^{-1})$, $u\in U$, que convergen en $L^p$ a la función indicada; sobre una vecindad fija de $e$, las normas y factores modulares están uniformemente acotados.
\end{solucion}
\begin{ejercicio}[Lectura 22, ejercicio 2]\label{ej:v4-l22-2}
Si $G$ es discreto, demostrar que $L^1(G)$ tiene unidad.
\end{ejercicio}
\begin{solucion}
Por invariancia, todos los puntos tienen una misma masa $c>0$; es finita porque los puntos son compactos y positiva porque son abiertos. La función $u=c^{-1}\mathbf1_{\{e\}}$ satisface $u*f=f$ directamente en la suma de convolución. También $f*u=f$, pues solo contribuye el término $y=x$, de masa $c$. Con medida de conteo $c=1$ y la unidad es $\delta_e$. No se confunde esta función con una masa puntual como medida en un grupo no discreto.
\end{solucion}
\begin{proposicion}[El recíproco anunciado en la lectura]\label{v4:luno-unitaria}
Si $L^1(G)$ tiene unidad, $G$ es discreto.
\end{proposicion}
\begin{proof}
Primero, si $\mu(\{e\})=c>0$, una vecindad de clausura compacta tiene a lo sumo $\mu(K)/c$ puntos: de lo contrario contendría subconjuntos finitos con medida arbitrariamente grande. Esa vecindad es finita y abierta; quitando sus otros puntos, cerrados por Hausdorff, obtenemos que $\{e\}$ es abierto. Así $G$ es discreto.

Supongamos, para obtener una contradicción, que $G$ no es discreto y que $u\in L^1(G)$ es unidad. Entonces $\mu(\{e\})=0$. Por regularidad exterior y continuidad absoluta de la integral de $|u|$, hay una vecindad abierta $V$ de $e$ con $\int_V|u|<1$. Elijamos $\phi\in\Cc(G)$ con $0\le\phi\le1$, $\phi(e)=1$ y soporte contenido en $V^{-1}$; la existencia de estos cortes se probó al preparar Haar. La identidad $u*\phi=\phi$ en $L^1$ vale en cada punto: ambos lados tienen representantes continuos, por continuidad de convolución de $L^1$ con una función continua de soporte compacto, y una función continua cero casi en todas partes es cero en todas partes porque todo abierto no vacío tiene medida positiva. Evaluando en $e$,
\[
1=\left|\int_Gu(y)\phi(y^{-1})dy\right|\le\int_V|u(y)|dy<1,
\]
contradicción. La continuidad utilizada también se obtiene directamente por convergencia uniforme de las traslaciones de $\phi$, dominando la diferencia por $|u|$ multiplicado por el error uniforme.
\end{proof}
\begin{ejercicio}[Lectura 22, ejercicio 3: control de soportes]\label{ej:v4-l22-3}
Para $f\in C(G)$, demostrar que $e_U*f$ y $f*e_U$ existen en todo punto y convergen puntualmente a $f$, usando la red con soportes $\supp e_U\subset U$. Mostrar que, para una red de Dirac arbitraria, no basta que cada núcleo tenga soporte compacto.
\end{ejercicio}
\begin{solucion}
En la primera convolución, la variable de integración está en $\supp e_U$; en la segunda, está en $x(\supp e_U)^{-1}$. Son compactos y $f$ es acotada en ellos, así que ambas integrales existen. Para $x$ fijo,
\[
 |e_U*f(x)-f(x)|\leq\sup_{y\in U}|f(y^{-1}x)-f(x)|\longrightarrow0.
\]
Para la otra convolución usamos
\[
 f*e_U(x)=\int e_U(z)\Delta(z)^{-1}f(xz^{-1})dz.
\]
La función de $z$ dentro del promedio es continua y vale $f(x)$ en $e$; el mismo argumento prueba su límite.

La concentración de masa sola no controla una función no acotada. En $\R$, sea $f(x)=e^{x^2}$, sea $b$ una función suave no negativa de integral uno apoyada en $[-1,1]$, y sea $h_n$ una identidad aproximada de soporte $[-1/n,1/n]$. Entonces $k_n=(1-1/n)h_n+n^{-1}b(\cdot-n)$ es una sucesión de Dirac, cada término de soporte compacto, pero
\[
 (k_n*f)(0)\geq n^{-1}e^{(n-1)^2}\longrightarrow\infty.
\]
Por tanto se corrige el enunciado usando soportes que se contraen, condición que cumple la red construida en este capítulo.
\end{solucion}
\begin{ejercicio}[Lectura 22, ejercicio 4]\label{ej:v4-l22-4}
Si $f\in L^1(G)$ y $g\in L^\infty(G)$, probar que $f*g$ está definida en todo punto, es acotada y satisface $\norm{L_a(f*g)-f*g}_\infty\to0$ cuando $a\to e$.
\end{ejercicio}
\begin{solucion}
Elijamos un representante acotado de $g$. Entonces la integral absoluta es a lo sumo $\norm f_1\norm g_\infty$ para cada $x$. Otro representante equivalente da la misma integral: inversión y traslaciones preservan nulos porque el factor modular es positivo. La identidad $L_a(f*g)=(L_af)*g$ da
\[
 \norm{L_a(f*g)-f*g}_\infty\leq\norm{L_af-f}_1\norm g_\infty\longrightarrow0.
\]
Esta es la continuidad uniforme por traslaciones izquierdas que se afirma; en particular implica continuidad ordinaria. En un grupo no abeliano conviene especificar la uniformidad utilizada en lugar de intercambiar sin explicación izquierda y derecha.
\end{solucion}
\begin{ejercicio}[Lectura 22, ejercicio 5]\label{ej:v4-l22-5}
Si $f,g\in C_c(G)$, demostrar $\supp(f*g)\subseteq\supp f\,\supp g$ y que $f*g\in C_c(G)$.
\end{ejercicio}
\begin{solucion}
Si $x$ no pertenece al producto de soportes, no puede ocurrir simultáneamente $y\in\supp f$ y $y^{-1}x\in\supp g$; el integrando es cero. El producto de soportes es compacto por continuidad de multiplicación. La convolución es continua por el ejercicio anterior, pues $g$ es acotada y $f\in L^1$. Así se anula fuera de un compacto y pertenece a $C_c$.
\end{solucion}

\paragraph{Los operadores que produce la convolución.}
La unidad puede faltar en $L^1(G)$, pero la red $(e_U)$ permite aproximar su acción. Esa red hará posible reconstruir operadores del grupo a partir de una representación del álgebra, que es el objetivo del capítulo siguiente.

\chapter{Representaciones unitarias y su forma integrada}\label{cap:v4-13}
Una traslación puede mirarse como un movimiento del grupo o como un operador sobre funciones. Una representación convierte la primera descripción en la segunda, preservando el producto y la geometría de un espacio de Hilbert. Estudiaremos sus subespacios invariantes y después construiremos, en ambas direcciones, el enlace con el álgebra de convolución. La medida de Haar y las identidades aproximadas del capítulo anterior son las herramientas que hacen posible esa correspondencia.

\section{Continuidad, equivalencia y representaciones regulares}\label{lectura:v4-23}
En este capítulo $G$ es localmente compacto y Hausdorff, $dx$ es Haar izquierda y los espacios de Hilbert son complejos. El producto interno es lineal en el primer argumento.
\begin{definicion}
Una \emph{representación unitaria fuertemente continua}\index{representacion@representación!unitaria} es un homomorfismo $\pi:G\to\mathcal U(H)$ tal que $x\mapsto\pi(x)u$ es continua para cada $u\in H$. En particular,
\[
\pi(e)=I,\qquad \pi(x^{-1})=\pi(x)^{-1}=\pi(x)^*.
\]
Un \emph{entrelazador}\index{entrelazador} entre $\pi_1$ en $H_1$ y $\pi_2$ en $H_2$ es un operador acotado $T:H_1\to H_2$ con $T\pi_1(x)=\pi_2(x)T$ para todo $x$. Las representaciones son equivalentes si admiten un entrelazador unitario.
\end{definicion}
La topología fuerte y la débil de operadores inducen la misma topología \emph{sobre} $\mathcal U(H)$. En efecto, si $T_\alpha\to T$ débilmente y tanto $T_\alpha$ como $T$ son unitarios,
\[
\norm{(T_\alpha-T)u}^2=2\norm u^2-2\Rea\ip{T_\alpha u}{Tu}\longrightarrow0.
\]
La implicación inversa es Cauchy--Schwarz. Es esencial que el límite también esté en $\mathcal U(H)$; el argumento no afirma que todo límite débil de unitarios sea unitario.

\begin{proposicion}[Representaciones regulares]\label{v4:rep-regulares}
En $L^2(G)$ las fórmulas
\[
[\lambda(x)f](y)=f(x^{-1}y),\qquad
[\rho(x)f](y)=\Delta(x)^{1/2}f(yx)
\]
definen representaciones unitarias fuertemente continuas. Son equivalentes mediante la involución unitaria
\[
(Jf)(y)=\Delta(y)^{-1/2}f(y^{-1}),\qquad J\lambda(x)J=\rho(x).
\]
\end{proposicion}
\begin{proof}
Las leyes de homomorfismo se verifican sustituyendo las fórmulas y usando que $\Delta$ es multiplicativa. La invariancia izquierda da $\norm{\lambda(x)f}_2=\norm f_2$; la identidad $\int h(yx)dy=\Delta(x)^{-1}\int h$ da la misma igualdad para $\rho(x)$. Sus inversas corresponden a $x^{-1}$, de modo que son sobreyectivas y unitarias.

Para $f\in\Cc(G)$, las traslaciones próximas a $e$ tienen soportes en un compacto común y convergen uniformemente a $f$. La norma $L^2$ de la diferencia queda acotada por la norma uniforme multiplicada por la raíz de la medida de ese compacto. Para $\rho$ usamos además $\Delta(x)^{1/2}\to1$. La densidad de $\Cc(G)$ y la estimación $\norm{\pi(x)f-f}_2\leq2\norm{f-h}_2+\norm{\pi(x)h-h}_2$ extienden la continuidad a todo $f\in L^2$. La ley de grupo la traslada de $e$ a cualquier punto.

La fórmula de inversión de Haar prueba
\[
\norm{Jf}_2^2=\int_G\Delta(y)^{-1}|f(y^{-1})|^2dy
=\int_G|f(z)|^2dz.
\]
Como $J^2=I$, $J$ es unitario. Finalmente,
\[
[J\lambda(x)Jf](y)
=\Delta(y)^{-1/2}\Delta(x^{-1}y^{-1})^{-1/2}f(yx)
=\Delta(x)^{1/2}f(yx).
\]
\end{proof}
\begin{observacion}[Corrección de la lectura 23]
La multiplicación por $\Delta^{1/2}$, sin inversión de la variable, no proporciona el entrelazador anunciado en la fuente. Tampoco $f(y)\mapsto f(y^{-1})$ es unitario con Haar izquierda en un grupo general. La fórmula de $J$ incorpora ambos cambios y se reduce a la inversión cuando el grupo es unimodular.
\end{observacion}

\section{Subespacios invariantes y descomposición cíclica}
Un subespacio $M$ es $\pi$-invariante si $\pi(x)M\subseteq M$ para todo $x$. Aplicar también $x^{-1}$ muestra que en realidad $\pi(x)M=M$. Su clausura es invariante por continuidad de cada operador. Si $M$ es cerrado, la restricción es una representación en $M$ y $M^\perp$ también es invariante, pues
\[
\ip{\pi(x)v}{m}=\ip{v}{\pi(x^{-1})m}=0\quad(v\in M^\perp,\ m\in M).
\]
Por tanto, la proyección $P_M$ conmuta con todos los $\pi(x)$: basta comparar ambos operadores sobre $M$ y sobre $M^\perp$. Recíprocamente, si $P_M$ conmuta, para $m=P_Mm$ se tiene $\pi(x)m=P_M\pi(x)m\in M$.

\begin{definicion}
Una representación en $H\ne\{0\}$ es \emph{irreducible}\index{representacion@representación!irreducible} si sus únicos subespacios \emph{cerrados} invariantes son $0$ y $H$. El subespacio cíclico generado por $u$ es
\[
H_u=\overline{\spanop\{\pi(x)u:x\in G\}}.
\]
La representación es \emph{cíclica}\index{representacion@representación!ciclica@cíclica} si $H_u=H$ para algún $u$. La clausura es parte de esta definición: una igualdad de clausuras no equivale a que cada vector sea una combinación lineal finita.
\end{definicion}
El espacio $H_u$ es invariante porque $\pi(y)\pi(x)u=\pi(yx)u$, primero en el span y después en su clausura. Para una familia de representaciones $(\pi_j,H_j)$, su suma directa actúa coordenada a coordenada sobre la suma directa hilbertiana. Preserva la norma y es sobreyectiva. También es fuertemente continua: dado $v=(v_j)$ y $\varepsilon>0$, elegimos un conjunto finito $F$ con $\sum_{j\notin F}\norm{v_j}^2<\varepsilon^2/16$. La contribución de esas coordenadas a $\norm{\pi(x)v-\pi(x_0)v}$ es menor que $\varepsilon/2$, uniformemente en $x$; la de $F$ tiende a cero por continuidad de sus finitas representaciones. Este argumento cubre familias no numerables, pues todo vector tiene a lo sumo un conjunto numerable de coordenadas no nulas.

\begin{teorema}\label{v4:rep-ciclicas}
Toda representación unitaria es una suma directa ortogonal de representaciones cíclicas.
\end{teorema}
\begin{proof}
Ordenemos por inclusión las familias de subespacios cíclicos no nulos, cerrados y ortogonales dos a dos. La familia vacía pertenece a este conjunto. La unión de una cadena sigue siendo una familia de ese tipo: dos de sus miembros pertenecen a una misma familia de la cadena. El lema de Zorn proporciona una familia maximal $(H_j)$. Sea $M$ la clausura de su suma algebraica, identificada isométricamente con $\bigoplus_jH_j$. Es invariante por lo anterior, y su complemento ortogonal también lo es. Si $M\ne H$, un vector $0\ne u\in M^\perp$ genera un subespacio cíclico no nulo contenido en $M^\perp$. Añadirlo contradice maximalidad. Así, $M=H$, y la identificación de la suma directa entrelaza las representaciones.
\end{proof}

\section{El lema de Schur y sus consecuencias}
\begin{teorema}[Lema de Schur]\label{v4:schur}\index{Schur, lema de}
Una representación unitaria $\pi$ es irreducible si y solo si su conmutante es $\C I$. Entre dos representaciones irreducibles, todo entrelazador no nulo es un múltiplo de un unitario. Si existe uno, fijado un entrelazador unitario $V$, todos los entrelazadores son $cV$; si no son equivalentes, el único entrelazador es cero.
\end{teorema}
\begin{proof}
Si $M$ es cerrado invariante no trivial, $P_M$ pertenece al conmutante y no es escalar. Para la dirección contraria supongamos irreducibilidad. Si $T$ conmuta con la representación, tomar adjuntas de la relación para $x^{-1}$ muestra que también $T^*$ conmuta. Entonces $A=(T+T^*)/2$ y $B=(T-T^*)/(2i)$ son autoadjuntos del conmutante. Basta demostrar que cada uno es escalar.

Si $A$ no fuera escalar, su espectro tendría mínimo $m$ y máximo $M$ distintos: para un autoadjunto, espectro reducido a $\{a\}$ implicaría $\norm{A-aI}=0$ por cálculo funcional. Elijamos $m<t<M$. Las funciones $r(s)=\max(s-t,0)$ y $q(s)=\max(t-s,0)$ dan operadores $r(A),q(A)$ no nulos con producto cero. Conmutan con $\pi(x)$, porque son límites en norma de polinomios en $A$. El espacio $\overline{\operatorname{ran}r(A)}$ es entonces cerrado, invariante y no nulo. Es propio, ya que la imagen no nula de $q(A)$ le es ortogonal: $\ip{r(A)u}{q(A)v}=\ip{u}{r(A)q(A)v}=0$. Contradice irreducibilidad. Este argumento usa únicamente el cálculo funcional continuo de autoadjuntos desarrollado en el volumen III. Luego $A,B$ son escalares y también $T=A+iB$.

Sea ahora $T\pi(x)=\pi'(x)T$. Tomando adjuntas para $x^{-1}$ obtenemos $\pi(x)T^*=T^*\pi'(x)$. Por ello $T^*T=\alpha I$ y $TT^*=\beta I$, con $\alpha,\beta\geq0$. Si $T\ne0$, ambas son positivas y la igualdad $(TT^*)T=T(T^*T)$ da $\beta=\alpha$. Por tanto $\alpha^{-1/2}T$ es unitario y entrelaza. Si $V$ es un entrelazador unitario fijo y $S$ otro entrelazador, $V^*S$ conmuta con $\pi$, luego es $cI$ y $S=cV$. Esto incluye $S=0$ sin usar su inversa.
\end{proof}
\begin{corolario}
Si $G$ es abeliano, toda representación unitaria irreducible tiene dimensión uno.
\end{corolario}
\begin{proof}
Cada $\pi(x)$ conmuta con todos los $\pi(y)$, luego es escalar por Schur. Todos los subespacios cerrados son entonces invariantes. Si $\dim H>1$, cualquier recta sería un subespacio cerrado propio no nulo, contradicción.
\end{proof}

\section{Ejercicios de entrelazado e invariancia}
\begin{ejercicio}[Lectura 23, ejercicio 1]\label{ej:v4-l23-1}
Demuestre que el conjunto $C(\pi,\pi')$ de entrelazadores es un subespacio de $\Bcal(H,H')$.
\end{ejercicio}

\begin{ejercicio}[Lectura 23, ejercicio 2]\label{ej:v4-l23-2}
Demuestre que el conmutante $C(\pi)$ es un álgebra de von Neumann.
\end{ejercicio}

\begin{ejercicio}[Lectura 23, ejercicio 3]\label{ej:v4-l23-3}
Si $M$ es cerrado e invariante, demuestre que $M^\perp$ es invariante.
\end{ejercicio}

\begin{ejercicio}[Lectura 23, ejercicio 4]\label{ej:v4-l23-4}
Para $M$ cerrado, pruebe que $M$ es invariante si y solo si $P_M\in C(\pi)$.
\end{ejercicio}


\section{Representaciones del álgebra de convolución}\label{lectura:v4-24}
La lectura 24 disponible es una reconstrucción editorial independiente, no un original recuperado de Valencia Bucio. Se integra aquí su desarrollo completo, enlazado con las identidades aproximadas del capítulo \ref{cap:v4-12}; su función es justificar los pasos que conectan las lecturas originales 23 y 25--27. La correspondencia que sigue es un resultado clásico de la teoría de representaciones \cite{reconstruccion24,folland2016}.

\begin{definicion}\label{v4-l24-d:rep}
Una representación unitaria fuertemente continua de $G$ en $H$ es una aplicación $U:G\to\mathcal U(H)$ que satisface $U(xy)=U(x)U(y)$ y para la cual $x\mapsto U(x)u$ es continua para todo $u\in H$. Una $*$-representación de $L^1(G)$ es una aplicación lineal $\rho:L^1(G)\to\Bcal(H)$ que satisface
\[
\rho(f*g)=\rho(f)\rho(g),\qquad \rho(f^*)=\rho(f)^*.
\]
Se llama no degenerada si
\[
\overline{\spanop\{\rho(f)u:f\in L^1(G),\ u\in H\}}=H.
\]
\end{definicion}

\begin{lema}\label{v4-l24-l:contractive}
Toda $*$-representación $\rho:L^1(G)\to\Bcal(H)$ es contractiva:
\[
\norm{\rho(f)}\leq\norm f_1.
\]
En particular, es continua, aun si no se incluyó continuidad en su definición.
\end{lema}
\begin{proof}
Usaremos un argumento de espectros, cuyos ingredientes recordamos. La unitarización de $A=L^1(G)$ es $\widetilde A=A\oplus\C$, con
\[
(a,z)(b,w)=(a*b+zb+wa,zw),\qquad\norm{(a,z)}=\norm a_1+|z|.
\]
Tiene identidad $(0,1)$ y es un álgebra de Banach: la multiplicación es submultiplicativa por la proposición \ref{v4-l24-p:banach}, y una sucesión de Cauchy converge en cada componente. La aplicación $\widetilde\rho(a,z)=\rho(a)+zI$ es un homomorfismo unital.

Si $\lambda(0,1)-(a,0)$ tiene inverso en $\widetilde A$, su imagen es un inverso de $\lambda I-\rho(a)$. Por tanto,
\[
\sigma_{\Bcal(H)}(\rho(a))\subseteq\sigma_{\widetilde A}((a,0)).
\]
Cuando $|\lambda|>\norm a_1$, la serie geométrica
$\lambda^{-1}\sum_{n\geq0}((a,0)/\lambda)^n$ converge en $\widetilde A$ y, multiplicando sus sumas parciales, se comprueba que es el inverso requerido. En consecuencia, el espectro de $\rho(a)$ está dentro del disco de radio $\norm a_1$.

Aplicamos esto a $a=f^**f$. El operador $\rho(a)=\rho(f)^*\rho(f)$ es positivo. Para un operador positivo acotado, el teorema espectral da
\[
\norm{\rho(f)}^2=\norm{\rho(f)^*\rho(f)}
=\sup\{|\lambda|:\lambda\in\sigma(\rho(a))\}.
\]
La inclusión anterior y la proposición \ref{v4-l24-p:banach} implican
\[
\norm{\rho(f)}^2\leq\norm{f^**f}_1\leq\norm f_1^2.
\]
Tomando raíces obtenemos la conclusión. Aquí el único resultado espectral utilizado es la igualdad entre norma y radio espectral para operadores positivos, estudiada en el volumen III; no se supone continuidad de $\rho$ para trasladar inversos.
\end{proof}

\begin{proposicion}\label{v4-l24-p:nd}
Para una $*$-representación $\rho$, son equivalentes:
\begin{enumerate}
\item $\rho$ es no degenerada.
\item Si $\rho(f)u=0$ para toda $f\in L^1(G)$, entonces $u=0$.
\item $\rho(e_U)u\to u$ para todo $u\in H$.
\end{enumerate}
\end{proposicion}
\begin{proof}
Sea $D=\spanop\{\rho(f)v\}$. Un vector $u$ es ortogonal a $D$ si y solo si
\[
0=\ip{\rho(f)v}{u}=\ip{v}{\rho(f)^*u}
\quad\text{para todos }f,v.
\]
Como un vector ortogonal a todo $v$ es cero, esto equivale a $\rho(f^*)u=0$ para toda $f$. La involución es biyectiva; por ello equivale a (ii). Además, $\overline D=H$ si y solo si $D^\perp=\{0\}$. Se obtiene (i)$\Leftrightarrow$(ii).

Supongamos (i). Si $w=\rho(f)v$, entonces
\[
\norm{\rho(e_U)w-w}
=\norm{\rho(e_U*f-f)v}\leq\norm{e_U*f-f}_1\norm v\longrightarrow0.
\]
Por linealidad vale para $w\in D$. Dado $u\in H$ y $\varepsilon>0$, tomamos $w\in D$ con $\norm{u-w}<\varepsilon/3$. Como $\norm{\rho(e_U)}\leq1$,
\[
\norm{\rho(e_U)u-u}\leq2\norm{u-w}+\norm{\rho(e_U)w-w}<\varepsilon
\]
para índices suficientemente posteriores. Así obtenemos (iii). Finalmente, (iii) implica (ii): si todas las $\rho(f)$ anulan $u$, también $\rho(e_U)u=0$, y al tomar el límite resulta $u=0$.
\end{proof}

\section{Integración de una representación}\index{representacion@representación!integrada}
\begin{proposicion}\label{v4-l24-p:integral}
Sea $U$ una representación unitaria fuertemente continua. Para cada $f\in L^1(G)$ existe un único operador $\Pi_U(f)\in\Bcal(H)$ tal que
\begin{equation}\label{v4-l24-eq:weak}
\ip{\Pi_U(f)u}{v}=\int_Gf(x)\ip{U(x)u}{v}\,dx
\quad(u,v\in H).
\end{equation}
Además, $\norm{\Pi_U(f)}\leq\norm f_1$.
\end{proposicion}
\begin{proof}
El coeficiente $x\mapsto\ip{U(x)u}{v}$ es continuo y está acotado por $\norm u\norm v$, por Cauchy--Schwarz y unitariedad. La integral de \eqref{v4-l24-eq:weak} existe absolutamente y su módulo es menor o igual que $\norm f_1\norm u\norm v$. Para $u$ fijo es una funcional antilineal continua de $v$. El teorema de representación de Riesz produce un único vector $w$ que la representa como $\ip{w}{v}$. Definimos $\Pi_U(f)u=w$.

La linealidad respecto de $u$ se verifica tomando producto interno con todo $v$: ambas expresiones para una combinación lineal tienen los mismos coeficientes y por unicidad son iguales. Asimismo,
\[
\norm{\Pi_U(f)u}=\sup_{\norm v=1}|\ip{\Pi_U(f)u}{v}|
\leq\norm f_1\norm u.
\]
Así, el operador es acotado con la norma anunciada. Si otro operador tiene los mismos coeficientes, su diferencia aplicada a $u$ es ortogonal a todo $v$ y por ello es cero. Esto prueba unicidad.
\end{proof}

Escribiremos $\Pi_U(f)=\int_Gf(x)U(x)dx$, entendiendo siempre esta integral por \eqref{v4-l24-eq:weak}. No se está suponiendo continuidad en norma de operadores. Una representación fuertemente continua puede no ser continua en esa norma. Este cuidado evita introducir una hipótesis que no pertenece al teorema.

\begin{teorema}[Representación integrada]\label{v4-l24-t:direct}
La aplicación $f\mapsto\Pi_U(f)$ de la proposición \ref{v4-l24-p:integral} es una $*$-representación no degenerada de $L^1(G)$. Además,
\begin{equation}\label{v4-l24-eq:cov}
U(x)\Pi_U(f)=\Pi_U(L_xf).
\end{equation}
\end{teorema}
\begin{proof}
La linealidad en $f$ se obtiene de \eqref{v4-l24-eq:weak} y unicidad. Para el producto, fijamos $u,v\in H$ y calculamos
\begin{align*}
\ip{\Pi_U(f)\Pi_U(g)u}{v}
&=\int_G\int_G f(x)g(y)\ip{U(xy)u}{v}\,dy\,dx\\
&=\int_G\int_G f(x)g(x^{-1}z)\ip{U(z)u}{v}\,dz\,dx\\
&=\int_G(f*g)(z)\ip{U(z)u}{v}\,dz.
\end{align*}
Para obtener la primera línea, pasamos $U(x)$ al segundo argumento como $U(x)^*$ y aplicamos \eqref{v4-l24-eq:weak} a $g$. El valor absoluto de la integral doble está acotado por $\norm f_1\norm g_1\norm u\norm v$; Fubini es aplicable. El cambio $z=xy$ usa invariancia izquierda. La igualdad de todos los coeficientes prueba multiplicatividad.

Para la involución, \eqref{v4-l24-eq:inv} da
\begin{align*}
\ip{\Pi_U(f^*)u}{v}
&=\int_G\Delta(x^{-1})\overline{f(x^{-1})}\ip{U(x)u}{v}\,dx\\
&=\int_G\overline{f(y)}\ip{U(y^{-1})u}{v}\,dy\\
&=\int_G\overline{f(y)}\ip{u}{U(y)v}\,dy
=\ip{u}{\Pi_U(f)v}
=\ip{\Pi_U(f)^*u}{v}.
\end{align*}
En la penúltima igualdad se utiliza la antilinealidad del segundo argumento. Por unicidad, $\Pi_U(f^*)=\Pi_U(f)^*$.

Por otra parte, para todo $u$,
\begin{equation}\label{v4-l24-eq:strongapprox}
\norm{\Pi_U(e_U)u-u}
\leq\int_G e_U(y)\norm{U(y)u-u}\,dy
\leq\sup_{y\in U}\norm{U(y)u-u}\longrightarrow0.
\end{equation}
La primera desigualdad se demuestra tomando coeficientes con vectores unitarios y usando la proposición \ref{v4-l24-p:integral}. El último límite es la continuidad fuerte en $e$. La proposición \ref{v4-l24-p:nd} implica no degeneración.

Finalmente,
\[
\ip{U(x)\Pi_U(f)u}{v}
=\int_G f(y)\ip{U(xy)u}{v}\,dy
=\int_G f(x^{-1}z)\ip{U(z)u}{v}\,dz.
\]
Este es el coeficiente de $\Pi_U(L_xf)$. Queda probada \eqref{v4-l24-eq:cov}.
\end{proof}

Para demostrar el recíproco necesitaremos pasar una integral a través de un operador lineal acotado. Damos el argumento para evitar un intercambio de límites injustificado. En particular, no utilizaremos un supuesto teorema de convergencia dominada para redes arbitrarias.

\section{Integrales vectoriales y reconstrucción}\index{integral!vectorial}
\begin{lema}[Integral de una función vectorial ponderada]\label{v4-l24-l:vector}
Sea $X$ un espacio de Banach, sea $F:G\to X$ continua y acotada, y sea $f\in L^1(G)$. Existe un vector, denotado $\int_Gf(x)F(x)dx$, tal que
\[
\ell\left(\int_G f(x)F(x)dx\right)=\int_G f(x)\ell(F(x))dx
\quad(\ell\in X^*).
\]
Su norma es menor o igual que $\norm f_1\sup_x\norm{F(x)}$. Si $A:X\to Y$ es lineal acotado entre espacios de Banach, entonces
\[
A\left(\int_G f(x)F(x)dx\right)=\int_G f(x)AF(x)dx.
\]
\end{lema}
\begin{proof}
Supongamos primero que $f$ se anula fuera de un compacto $K$. Por continuidad, $F(K)$ es compacto. Para cada $n$ cubrimos $K$ por una cantidad finita de abiertos relativos en cada uno de los cuales los valores de $F$ distan menos de $1/n$ de algún valor $w_j=F(x_j)$. Quitando de cada abierto los anteriores, obtenemos una partición Borel $E_1,\ldots,E_m$ de $K$. La función simple $F_n=\sum_j w_j\chi_{E_j}$ satisface $\norm{F_n-F}<1/n$ en $K$. Pongamos
\[
v_n=\sum_j w_j\int_{E_j}f(x)dx.
\]
Al usar la partición común a $F_n$ y $F_m$, la desigualdad triangular da
\[
\norm{v_n-v_m}\leq(1/n+1/m)\norm f_1.
\]
Por completitud existe $v=\lim v_n$. Para toda $\ell\in X^*$, la estimación uniforme permite pasar al límite en las integrales escalares y obtener $\ell(v)=\int f\ell(F)$. La misma estimación prueba $\norm v\leq\norm f_1\sup\norm F$. La independencia de las particiones se obtiene comparando dos aproximaciones mediante su refinamiento común. La unicidad también se sigue de que $X^*$ separa puntos, por Hahn--Banach.

Para $f$ general, elijamos $f_n\in\Cc(G)$ con $\norm{f_n-f}_1\to0$. Las integrales ya construidas satisfacen
\[
\norm{\int f_nF-\int f_mF}\leq\norm{f_n-f_m}_1\sup\norm F.
\]
Su límite existe y no depende de la sucesión elegida. Las integrales escalares convergen por la misma cota, de modo que el límite tiene la propiedad anunciada y conserva la estimación. Por último, para cada $\eta\in Y^*$, la funcional $\eta A$ pertenece a $X^*$. Aplicar la identidad ya probada a $\eta A$ muestra que los dos vectores de la última fórmula tienen los mismos valores bajo toda $\eta$. Hahn--Banach implica que son iguales.
\end{proof}

\begin{observacion}\label{v4-l24-o:convvector}
El lema anterior se aplica a $F(x)=L_xg$ como función con valores en $L^1(G)$: el lema \ref{v4-l24-l:translation} da continuidad y $\norm{L_xg}_1=\norm g_1$. En este sentido,
\begin{equation}\label{v4-l24-eq:convvector}
f*g=\int_Gf(x)L_xg\,dx\quad\text{en }L^1(G).
\end{equation}
Para comprobar la identidad sin identificar indebidamente valores puntuales de clases $L^1$, sea $q\in L^\infty(G)$. La funcional $h\mapsto\int h(y)q(y)dy$ es acotada en $L^1$. El lema \ref{v4-l24-l:vector} y Fubini muestran que su valor sobre el lado derecho de \eqref{v4-l24-eq:convvector} es
\[
\int_G\int_G f(x)g(x^{-1}y)q(y)\,dy\,dx
=\int_G(f*g)(y)q(y)\,dy.
\]
La integral absoluta es menor o igual que $\norm f_1\norm g_1\norm q_\infty$. Estas funcionales separan puntos de $L^1$: si $h\ne0$, elegir $q=\overline h/|h|$ donde $h\ne0$ y $q=0$ en otro caso da $\int hq=\norm h_1>0$. Por tanto, los dos lados de \eqref{v4-l24-eq:convvector} son iguales.
\end{observacion}

\begin{teorema}[Reconstrucción de la representación del grupo]\label{v4-l24-t:inverse}
Sea $\rho:L^1(G)\to\Bcal(H)$ una $*$-representación no degenerada. Existe una única representación unitaria fuertemente continua $U$ de $G$ tal que $\rho=\Pi_U$. Para todo $x\in G$ se cumple
\begin{equation}\label{v4-l24-eq:reconstruct}
U(x)\rho(f)u=\rho(L_xf)u,
\qquad
U(x)=\operatorname*{s-lim}_{U_0}\rho(L_xe_{U_0}).
\end{equation}
El subíndice $U_0$ designa aquí una vecindad, para distinguirlo de la representación $U$.
\end{teorema}
\begin{proof}
\textit{Construcción sobre un subespacio denso.} Sea $D=\spanop\{\rho(f)u\}$, denso por hipótesis. Para $x\in G$ fijo, definimos provisionalmente
\[
V_x\left(\sum_{j=1}^n\rho(f_j)u_j\right)
=\sum_{j=1}^n\rho(L_xf_j)u_j.
\]
Debemos demostrar que el valor no depende de la expresión del vector. La identidad
\[
(L_xe_{U_0})*f=L_x(e_{U_0}*f)
\]
se verifica directamente mediante el cambio $y=xz$ en la convolución. Por los lemas \ref{v4-l24-l:approx} y \ref{v4-l24-l:contractive},
\[
\rho(L_xe_{U_0})\sum_j\rho(f_j)u_j
\longrightarrow\sum_j\rho(L_xf_j)u_j.
\]
Si la suma inicial es cero, cada término de la red a la izquierda es cero; su límite también lo es. Dos expresiones de un mismo vector tienen diferencia cero, de modo que sus imágenes coinciden. Esto prueba que $V_x$ está bien definida y es lineal en $D$.

\textit{Extensión acotada.} Como $\norm{\rho(L_xe_{U_0})}\leq\norm{L_xe_{U_0}}_1=1$, tomando el límite anterior obtenemos $\norm{V_xw}\leq\norm w$ para $w\in D$. Para $u\in H$, elegimos $w_n\in D$ con $w_n\to u$. La desigualdad hace que $(V_xw_n)$ sea de Cauchy. Su límite no depende de la sucesión, pues $\norm{V_xw_n-V_xz_n}\leq\norm{w_n-z_n}$. Definimos $U(x)u=\lim V_xw_n$. Esta extensión es lineal, tiene norma a lo sumo uno y es única por densidad.

\textit{Ley de grupo y unitariedad.} La fórmula sobre generadores da
\[
U(x)U(y)\rho(f)u=\rho(L_xL_yf)u
=\rho(L_{xy}f)u=U(xy)\rho(f)u.
\]
Por linealidad vale en $D$, y por continuidad en $H$. Del mismo modo $U(e)=I$. Así, $U(x^{-1})$ es el inverso de $U(x)$. Ambos son contracciones; por ello
\[
\norm u=\norm{U(x^{-1})U(x)u}\leq\norm{U(x)u}\leq\norm u.
\]
Entonces $U(x)$ es una isometría sobreyectiva. Una isometría lineal preserva el producto interno por la identidad de polarización; su inversa coincide con su adjunta. Por tanto, $U(x)$ es unitaria.

\textit{Continuidad fuerte.} Para $w=\sum_j\rho(f_j)u_j\in D$,
\[
\norm{U(x)w-U(x_0)w}
\leq\sum_j\norm{L_xf_j-L_{x_0}f_j}_1\norm{u_j}\longrightarrow0
\]
cuando $x\to x_0$, por el lema \ref{v4-l24-l:translation}. Para $u\in H$, tomamos $w\in D$ próximo a $u$ y usamos unitariedad:
\[
\norm{U(x)u-U(x_0)u}
\leq2\norm{u-w}+\norm{U(x)w-U(x_0)w}.
\]
Primero hacemos pequeño el primer término y después elegimos $x$ suficientemente cerca de $x_0$. Así obtenemos continuidad para todo $u$.

\textit{Recuperación de $\rho$.} Por \eqref{v4-l24-eq:convvector}, para $f,g\in L^1(G)$ y $u\in H$,
\begin{align*}
\Pi_U(f)\rho(g)u
&=\int_G f(x)U(x)\rho(g)u\,dx\\
&=\int_G f(x)\rho(L_xg)u\,dx
=\rho\left(\int_G f(x)L_xg\,dx\right)u\\
&=\rho(f*g)u=\rho(f)\rho(g)u.
\end{align*}
El paso de $\rho$ a través de la integral es el lema \ref{v4-l24-l:vector}, aplicado al operador acotado $h\mapsto\rho(h)u$; su acotación procede del lema \ref{v4-l24-l:contractive}. La primera integral coincide con la definición débil de $\Pi_U$ al tomar coeficientes. Los operadores $\Pi_U(f)$ y $\rho(f)$ coinciden sobre $D$ y, al ser acotados, sobre $H$.

\textit{Límite fuerte y unicidad.} En $D$ ya probamos la convergencia de $\rho(L_xe_{U_0})$ a $U(x)$. La cota uniforme uno extiende la convergencia a $H$ mediante la misma estimación de aproximación usada para continuidad. Si $W$ es otra representación con $\Pi_W=\rho$, \eqref{v4-l24-eq:cov} y \eqref{v4-l24-eq:strongapprox} dan
\[
W(x)u=\lim_{U_0}W(x)\rho(e_{U_0})u
=\lim_{U_0}\rho(L_xe_{U_0})u=U(x)u.
\]
Esto prueba unicidad y concluye el teorema.
\end{proof}

\section{Entrelazado, invariancia y ciclicidad}
\begin{teorema}[Entrelazado e invariancia]\label{v4-l24-t:intertwine}
Sean $U_j:G\to\mathcal U(H_j)$ representaciones unitarias fuertemente continuas, y sean $\Pi_j$ sus representaciones integradas. Para $T\in\Bcal(H_1,H_2)$ son equivalentes
\[
TU_1(x)=U_2(x)T\quad(x\in G),
\qquad T\Pi_1(f)=\Pi_2(f)T\quad(f\in L^1(G)).
\]
Además, un subespacio cerrado $M\subseteq H$ es invariante bajo $U(G)$ si y solo si lo es bajo $\Pi_U(L^1(G))$.
\end{teorema}
\begin{proof}
Supongamos la identidad de entrelazado para $G$. Tomando coeficientes y usando $T^*$, la integral de $f(x)TU_1(x)u$ tiene los mismos coeficientes que $T\Pi_1(f)u$. Sustituir $TU_1(x)$ por $U_2(x)T$ da $\Pi_2(f)Tu$. Por unicidad de coeficientes se obtiene la identidad para $L^1(G)$.

Recíprocamente, fijemos $x\in G$. Aplicamos la identidad para funciones a $L_xe_{U_0}$ y tomamos límites fuertes del teorema \ref{v4-l24-t:inverse}. Para cada $u$,
\[
TU_1(x)u=\lim T\Pi_1(L_xe_{U_0})u
=\lim\Pi_2(L_xe_{U_0})Tu=U_2(x)Tu.
\]
La continuidad de $T$ permite el primer paso.

Si $M$ es $U$-invariante, $u\in M$ y $v\in M^\perp$, entonces $\ip{U(x)u}{v}=0$ para todo $x$; integrar da $\ip{\Pi_U(f)u}{v}=0$. Como $M$ es cerrado, $(M^\perp)^\perp=M$, de modo que $\Pi_U(f)u\in M$.

Si $M$ es invariante bajo todas las $\Pi_U(f)$, también contiene $\Pi_U(L_xe_{U_0})u$ cuando $u\in M$. Estos vectores convergen a $U(x)u$. La cerradura de $M$ implica que el límite pertenece a $M$. Esto prueba la otra dirección.
\end{proof}

\begin{corolario}\label{v4-l24-c:invariants}
La correspondencia $U\leftrightarrow\Pi_U$ conserva equivalencia unitaria, irreducibilidad y vectores cíclicos. Además,
\[
\{T:TU(x)=U(x)T\ \forall x\}
=\{T:T\Pi_U(f)=\Pi_U(f)T\ \forall f\}.
\]
\end{corolario}
\begin{proof}
La igualdad es el teorema \ref{v4-l24-t:intertwine} con $H_1=H_2$. Si el entrelazador es unitario, ese mismo teorema prueba equivalencia en ambas direcciones. La irreducibilidad significa que los únicos subespacios \emph{cerrados} invariantes son $\{0\}$ y $H$; la equivalencia de invariancia demuestra la afirmación.

Para un vector $u$ pongamos $M_G=\overline{\spanop\{U(x)u:x\in G\}}$ y $M_A=\overline{\spanop\{\Pi_U(f)u:f\in L^1(G)\}}$. Si $v\perp M_G$, todos los coeficientes $\ip{U(x)u}{v}$ son cero y su integral también; entonces $v\perp M_A$. Esto implica $M_A\subseteq M_G$. Por otro lado, $U(x)u=\lim\Pi_U(L_xe_{U_0})u\in M_A$, de donde $M_G\subseteq M_A$. Por tanto, $M_G=M_A$, y uno es $H$ si y solo si el otro lo es.
\end{proof}

\section{Representación regular, caracteres y completaciones}
\begin{ejemplo}[Representación regular]\label{v4-l24-e:regular}
La representación regular izquierda en $L^2(G)$ viene dada por
\[
[\lambda(x)u](y)=u(x^{-1}y).
\]
Es unitaria por invariancia izquierda; su continuidad fuerte se prueba como en el lema \ref{v4-l24-l:translation}, usando densidad de $\Cc(G)$ en $L^2(G)$ y la norma $L^2$. Su representación integrada satisface
\[
\Pi_\lambda(f)u=f*u,\qquad
\norm{f*u}_2\leq\norm f_1\norm u_2.
\]
Para $f,u\in\Cc(G)$, la fórmula se obtiene sustituyendo la definición de $\lambda$ en los coeficientes e intercambiando integrales absolutamente convergentes sobre compactos. Para $f\in L^1$ y $u\in\Cc$, aproximamos $f$ en $L^1$ y usamos la cota de la proposición \ref{v4-l24-p:integral}. Después extendemos por densidad en $u$ a $L^2$; esta extensión define el operador de convolución en $L^2$. Es la construcción que utilizaremos en el capítulo \ref{cap:v4-16}.
\end{ejemplo}

\begin{ejemplo}[Caracteres y Fourier]\label{v4-l24-e:characters}
Sea $\chi:G\to\T$ un homomorfismo continuo. La fórmula $U_\chi(x)z=\chi(x)z$ define una representación unitaria en $\C$. Su forma integrada es la funcional
\[
\Phi_\chi(f)=\int_G f(x)\chi(x)\,dx,
\]
que satisface $\Phi_\chi(f*g)=\Phi_\chi(f)\Phi_\chi(g)$ y $\Phi_\chi(f^*)=\overline{\Phi_\chi(f)}$. Esto justifica el paso usado en la observación 4 de la lectura 26.

Si se adopta $\widehat f(\chi)=\int f(x)\overline{\chi(x)}dx$, entonces $\widehat f(\chi)=\Phi_{\overline\chi}(f)$. La aparición o ausencia del conjugado corresponde a elegir $\chi$ o su inverso; debe mantenerse una sola convención durante un cálculo. El teorema de correspondencia da el recíproco para funcionales no nulas que preservan la involución. Para funcionales multiplicativas arbitrarias hay que demostrar además esa compatibilidad; no se debe atribuir ese paso al teorema \ref{v4-l24-t:inverse} sin comprobar su hipótesis.
\end{ejemplo}

\begin{observacion}[Alcance de la reconstrucción]
En un grupo discreto, $\delta_x$ es una función de $L^1(G)$ y $U(x)=\Pi_U(\delta_x)$ cuando se usa medida de conteo. En un grupo no discreto, la masa de Dirac $\delta_x$ es una medida, no un elemento de $L^1(G)$. El límite de \eqref{v4-l24-eq:reconstruct} reemplaza esa evaluación inexistente. Se trata de convergencia fuerte, no necesariamente de convergencia en norma de operadores.
\end{observacion}

La lectura 27 parte de la representación regular para construir un álgebra $C^*$. Conviene registrar la distinción necesaria para grupos no abelianos.

\begin{definicion}\label{v4-l24-d:cstar}
El álgebra reducida de grupo es
\[
C_r^*(G)=\overline{\Pi_\lambda(L^1(G))}^{\,\norm{\cdot}_{\Bcal(L^2(G))}}.
\]
La norma universal de grupo sobre $L^1(G)$ es
\[
\norm f_{\mathrm{univ}}=\sup_U\norm{\Pi_U(f)},
\]
donde basta tomar un representante de cada clase de representaciones cíclicas. En efecto, la descomposición del teorema \ref{v4:rep-ciclicas} convierte toda representación en una suma directa, cuya norma es el supremo de las normas de sus sumandos. Además, si $D$ es un subconjunto denso de $L^1(G)$ de cardinal $\kappa$, una representación cíclica tiene un conjunto total de cardinal a lo sumo $\max(\kappa,\aleph_0)$, formado por combinaciones racionales complejas de $\Pi_U(D)u$. Sus espacios de Hilbert pueden por ello elegirse entre los espacios $\ell^2(I)$ con cardinal acotado por ese número. Sobre dichos espacios las representaciones forman un conjunto, de modo que el supremo no recorre una clase propia. La completación en esa norma se denota $C^*(G)$.
\end{definicion}

\begin{observacion}
La cota $\norm{\Pi_U(f)}\leq\norm f_1$ hace finito el supremo. La representación regular es fiel sobre $L^1(G)$: se demuestra en el ejercicio \ref{ej:v4-l24-5}. Por ello el supremo es una norma y no solo una seminorma. Las identidades de representación dan
\[
\norm{f^**f}_{\mathrm{univ}}
=\sup_U\norm{\Pi_U(f)^*\Pi_U(f)}
=\norm f_{\mathrm{univ}}^2.
\]
Así, su completación es un álgebra $C^*$. Una representación integrada se extiende de manera única a $C^*(G)$: si $f_n$ es de Cauchy en norma universal, la desigualdad $\norm{\Pi_U(f_n)-\Pi_U(f_m)}\leq\norm{f_n-f_m}_{\mathrm{univ}}$ proporciona el límite, y garantiza independencia y continuidad.

Para grupos localmente compactos abelianos, el teorema de Gelfand y el análisis de Fourier identifican ambas completaciones con $C_0(\widehat G)$. Este resultado pertenece al desarrollo posterior de las lecturas 25--27 y no se utiliza en las demostraciones anteriores. En general, la clausura de la representación regular debe llamarse $C_r^*(G)$; la correspondencia con \emph{todas} las representaciones unitarias se formula para $C^*(G)$.
\end{observacion}

\section{Ejercicios de la correspondencia}

Los ocho ejercicios siguientes verifican consecuencias y casos concretos del teorema. Sus soluciones forman parte del apunte y no requieren completar argumentos omitidos.

\begin{ejercicio}[Lectura 24 (reconstrucción), ejercicio 1]\label{ej:v4-l24-1}
\textbf{La representación trivial y una representación no fiel.}
Demuestre que $\varepsilon(f)=\int_G f(x)dx$ es una $*$-representación no degenerada en $\C$. Para $G=\R$, encuentre una función no nula en su núcleo.

\end{ejercicio}


\begin{ejercicio}[Lectura 24 (reconstrucción), ejercicio 2]\label{ej:v4-l24-2}
\textbf{Reconstrucción en $\mathbb Z$.}
Sea $T$ un operador unitario en $H$. Para $a=(a_n)_{n\in\mathbb Z}\in\ell^1(\mathbb Z)$ defina $\rho_T(a)=\sum_{n\in\mathbb Z}a_nT^n$. Pruebe que es una $*$-representación no degenerada y recupere su representación de grupo.

\end{ejercicio}


\begin{ejercicio}[Lectura 24 (reconstrucción), ejercicio 3]\label{ej:v4-l24-3}
\textbf{Una identidad aproximada explícita en $\R$.}
Con medida de Lebesgue, sea $e_n(x)=\tfrac n2\chi_{[-1/n,1/n]}(x)$. Demuestre que $e_n*f\to f$ en $L^1(\R)$ y que $\Pi_U(e_n)u\to u$ para toda representación unitaria fuertemente continua de $\R$.

\end{ejercicio}


\begin{ejercicio}[Lectura 24 (reconstrucción), ejercicio 4]\label{ej:v4-l24-4}
\textbf{Evaluación de Fourier como representación.}
Para $\xi\in\R$ ponga $\chi_\xi(x)=e^{-i\xi x}$ y $\Phi_\xi(f)=\int_\R f(x)e^{-i\xi x}dx$. Verifique multiplicatividad, compatibilidad con involución y $\norm{\Phi_\xi}=1$.

\end{ejercicio}


\begin{ejercicio}[Lectura 24 (reconstrucción), ejercicio 5]\label{ej:v4-l24-5}
\textbf{Fidelidad de la representación regular en $L^1(G)$.}
Si $\Pi_\lambda(f)=0$ como operador sobre $L^2(G)$, demuestre que $f=0$ en $L^1(G)$. Explique por qué esto no identifica automáticamente $C^*(G)$ con $C_r^*(G)$.

\end{ejercicio}


\begin{ejercicio}[Lectura 24 (reconstrucción), ejercicio 6]\label{ej:v4-l24-6}
\textbf{Por qué es indispensable la no degeneración.}
Sea $U_0$ una representación en $H_0$, sea $K\ne\{0\}$ y defina $\rho(f)=\Pi_{U_0}(f)\oplus0$ en $H_0\oplus K$. Pruebe que $\rho$ es una $*$-representación que no puede ser la representación integrada de una representación unitaria en todo $H_0\oplus K$.

\end{ejercicio}


\begin{ejercicio}[Lectura 24 (reconstrucción), ejercicio 7]\label{ej:v4-l24-7}
\textbf{Conmutante e irreducibilidad.}
Pruebe que $U$ es irreducible si y solo si todo operador que conmuta con todas las $\Pi_U(f)$ es un múltiplo escalar de $I$. Precise qué significa aquí ``invariante''.

\end{ejercicio}


\begin{ejercicio}[Lectura 24 (reconstrucción), ejercicio 8]\label{ej:v4-l24-8}
\textbf{Convergencia fuerte que no es convergencia en norma.}
Para la representación regular de $\R$, pruebe que los operadores del ejercicio \ref{ej:v4-l24-3} convergen fuertemente a $I$, pero no en norma de operadores. Puede utilizar el teorema de Plancherel y la convención $\widehat f(\xi)=\int_\R f(x)e^{-i\xi x}dx$.

\end{ejercicio}




\paragraph{Cierre.} La representación del grupo y su forma integrada contienen la misma información cuando se exige no degeneración. La continuidad fuerte permite recuperar cada operador del grupo mediante funciones integrables concentradas cerca de un punto; no hace falta convertir una masa de Dirac en una función. En los capítulos siguientes, los caracteres serán el caso unidimensional de esta correspondencia y la representación regular proporcionará el paso de Fourier a las álgebras $C^*$.

\parte{Caracteres y dualidad}{Para grupos abelianos, las representaciones irreducibles son caracteres. La topología del dual, la transformada abstracta y las álgebras asociadas reúnen los modelos finitos, periódicos y euclidianos en una misma teoría.}
\chapter{Caracteres y topología del grupo dual}\label{cap:v4-14}
Una frecuencia en la recta es la función $x\mapsto e^{i\xi x}$: transforma una suma de posiciones en un producto de números de módulo uno. Esa propiedad tiene sentido en cualquier grupo topológico y permite definir frecuencias sin coordenadas. El conjunto de esas frecuencias también tendrá una estructura de grupo y una topología. La lectura 25 de \cite{valencia4} proporciona el punto de partida; precisaremos la topología y las hipótesis de sus ejercicios.

Trabajamos con grupos de Hausdorff. Escribimos $\T=\{z\in\C:|z|=1\}$ y usamos notación multiplicativa, salvo en grupos como $\R$ y $\Z$. Un grupo \emph{LCA} es un grupo localmente compacto, abeliano y de Hausdorff. Para demostrar que el dual de un grupo LCA vuelve a ser localmente compacto necesitaremos el capítulo~\ref{cap:v4-15}. El ejercicio de exactitud requiere además la dualidad que se demostrará en el capítulo~\ref{cap:v4-16}; se indica esta dependencia expresamente para no esconder un teorema profundo bajo una verificación elemental.

\section{Frecuencias como homomorfismos}\label{lectura:v4-25}
\index{caracter@carácter!de un grupo}\index{grupo!dual}
\begin{definicion}
Un \emph{carácter} de $G$ es un homomorfismo continuo $\chi:G\to\T$. El conjunto de todos ellos se denota $\widehat G$. El producto de caracteres es puntual: $(\chi\eta)(x)=\chi(x)\eta(x)$.
\end{definicion}

\begin{proposicion}
El conjunto $\widehat G$ es un grupo abeliano. Su identidad es la función constante $1$ y $\chi^{-1}(x)=\overline{\chi(x)}=\chi(x^{-1})$. Sus elementos son exactamente las representaciones unitarias unidimensionales de $G$.
\end{proposicion}
\begin{proof}
El producto de funciones continuas es continuo, y
\[
 (\chi\eta)(xy)=\chi(x)\chi(y)\eta(x)\eta(y)
 =\chi(x)\eta(x)\chi(y)\eta(y).
\]
Se ha usado que los valores escalares conmutan, aunque $G$ no sea abeliano. La conjugación también preserva continuidad y la propiedad de homomorfismo. Las demás identidades del grupo se comprueban punto por punto en $\T$. Todo operador lineal sobre $\C$ es multiplicación por su valor en $1$, y preserva el producto interno si y sólo si ese valor tiene módulo uno. Por tanto una representación unitaria continua sobre $\C$ equivale precisamente a un carácter.
\end{proof}

\begin{lema}[No hay subgrupos pequeños en el círculo]\label{v4:arco-sin-subgrupos}
El único subgrupo de $\T$ contenido en $\{z:\operatorname{Re}z>0\}$ es $\{1\}$.
\end{lema}
\begin{proof}
Si $z=e^{i\theta}\ne1$, podemos elegir $0<|\theta|\le\pi$. Si $|\theta|\ge\pi/2$, el propio $z$ no pertenece al semiplano. Si $0<|\theta|<\pi/2$, elija el primer entero $n$ con $n|\theta|\ge\pi/2$. Entonces $\pi/2\le n|\theta|<\pi$, y $\operatorname{Re}(z^n)\le0$. Un subgrupo que contenga $z$ contiene esa potencia, de modo que no puede estar contenido en el semiplano indicado.
\end{proof}

\section{Convergencia uniforme sobre compactos}
\index{topologia@topología!compacto-abierta}
Para un espacio localmente compacto $X$, un compacto $K\subset X$ y un abierto $U\subset\C$, ponemos
\[
 [K,U]=\{f\in C(X):f(K)\subset U\}.
\]
La \emph{topología compacto-abierta} sobre $C(X)$ es la generada por estos conjuntos como subbase. La topología se define en el espacio de funciones $C(X)$, no en $X$; esto corrige una errata del enunciado original. La familia cubre $C(X)$, pues $[K,\C]=C(X)$ para cualquier $K$.

\begin{teorema}\label{v4:compacto-abierta}
El espacio $C(X)$ es Hausdorff con esta topología. Una red $f_\alpha$ converge a $f$ si y sólo si, para cada compacto $K\subset X$, $\sup_K|f_\alpha-f|\to0$.
\end{teorema}
\begin{proof}
Si $f\ne g$, existe $x$ con $f(x)\ne g(x)$. Dos discos abiertos disjuntos alrededor de esos valores producen vecindades disjuntas $[\{x\},U]$ y $[\{x\},V]$.

Supongamos convergencia uniforme en compactos y sea $f\in[K,U]$. Si $U=\C$ no hay nada que probar. En otro caso, el compacto $f(K)$ tiene distancia positiva al cerrado $\C\setminus U$: si esa distancia fuera cero, una sucesión en $f(K)$ con distancia tendiendo a cero tendría una subsucesión cuyo límite pertenecería simultáneamente a ambos conjuntos. Para un error uniforme menor que esa distancia, $f_\alpha(K)\subset U$. Esto vale para cada miembro de una intersección finita de subbásicos, y el conjunto de índices de una red permite elegir un umbral común para un número finito de condiciones.

Recíprocamente, sea $\eps>0$. La compacidad de $f(K)$ permite elegir $x_1,\ldots,x_m\in K$ cuyos discos de radio $\eps/4$ cubran $f(K)$. Los conjuntos
\[
 K_j=\{x\in K:|f(x)-f(x_j)|\le\eps/3\}
\]
son compactos, cubren $K$ y satisfacen $f(K_j)\subset B(f(x_j),2\eps/3)$. Así
\[
 \mathcal U=\bigcap_{j=1}^m[K_j,B(f(x_j),2\eps/3)]
\]
es una vecindad de $f$. Eventualmente $f_\alpha\in\mathcal U$. Para $x\in K_j$ obtenemos
$|f_\alpha(x)-f(x)|<2\eps/3+\eps/3=\eps$. Para asegurar una desigualdad estricta después de tomar supremos se repite el argumento con $\eps/2$. Resulta convergencia uniforme en $K$.
\end{proof}

En la lectura se intersectan ciertos abiertos con $K$ y se afirma que los conjuntos resultantes son compactos. Eso no es cierto en general. Los conjuntos cerrados $K_j$ de la prueba anterior evitan ese problema. Las redes permiten trabajar sin suponer una base numerable de vecindades: sus índices forman un conjunto dirigido, es decir, cada par tiene una cota superior, y «eventualmente» significa después de algún índice.

Dotamos a $\widehat G$ de la topología inducida por $C(G)$. Sus vecindades pueden describirse mediante errores uniformes en un número finito de compactos.
\begin{proposicion}\label{v4:dual-topologico}
El grupo $\widehat G$ es un grupo topológico abeliano y Hausdorff. Además, si $G,H$ son localmente compactos,
\[
 \widehat{G\times H}\simeq\widehat G\times\widehat H
\]
como grupos topológicos.
\end{proposicion}
\begin{proof}
Para caracteres, $|\chi(x)|=1$, y por tanto
\[
 |\chi(x)\overline{\eta(x)}-\chi_0(x)\overline{\eta_0(x)}|
 \le|\chi(x)-\chi_0(x)|+|\eta(x)-\eta_0(x)|.
\]
Tomar supremos en compactos demuestra continuidad conjunta del producto por el inverso; en particular da continuidad de producto e inversión. La separación Hausdorff se hereda de $C(G)$.

Para $\gamma\in\widehat{G\times H}$, las restricciones $\chi(x)=\gamma(x,e_H)$ y $\eta(y)=\gamma(e_G,y)$ son continuas y multiplicativas, y
$\gamma(x,y)=\chi(x)\eta(y)$. Recíprocamente, esa fórmula produce un carácter a partir de cualquier par $(\chi,\eta)$. Son aplicaciones inversas y homomorfismos. La restricción es continua porque $K\times\{e_H\}$ y $\{e_G\}\times L$ son compactos. La aplicación del par al producto es continua porque un compacto $S\subset G\times H$ tiene proyecciones compactas $K,L$, y el error uniforme del producto en $S$ está acotado por la suma de los errores de los factores en $K,L$.
\end{proof}

\begin{lema}[Compacidad de frecuencias y control común]\label{v4:dual-equicontinuidad}
Si $G$ es localmente compacto, la evaluación $(\chi,x)\mapsto\chi(x)$ en $\widehat G\times G$ es continua. Si $E\subset\widehat G$ es compacto, para todo $\eps>0$ existe una vecindad $U$ de $e$ tal que
\[
 \sup_{\chi\in E,\ x\in U}|\chi(x)-1|<\eps.
\]
\end{lema}
\begin{proof}
Fijados $(\chi_0,x_0)$, elija una vecindad de $x_0$ con clausura compacta $K$. En esa vecindad,
\[
 |\chi(x)-\chi_0(x_0)|\le\sup_K|\chi-\chi_0|+|\chi_0(x)-\chi_0(x_0)|.
\]
Se controla el primer término con una vecindad de $\chi_0$ y el segundo con una vecindad de $x_0$. Para la última afirmación, la continuidad en cada $(\chi,e)$ proporciona vecindades $V_\chi$ de $\chi$ y $U_\chi$ de $e$ donde el error es menor que $\eps/2$. Una subcubierta finita de los $V_\chi$ cubre $E$; la intersección de los correspondientes $U_\chi$ sirve simultáneamente para todo $E$.
\end{proof}

\section{Los modelos compacto, discreto y euclidiano}
\begin{teorema}\label{v4:dual-compacto-discreto}
Si $G$ es compacto, $\widehat G$ es discreto. Si $G$ es discreto, $\widehat G$ es compacto.
\end{teorema}
\begin{proof}
Si $G$ es compacto, el conjunto de caracteres cuya imagen está contenida en el semiplano derecho es una vecindad abierta de $1$ para la topología compacto-abierta. Esa imagen es un subgrupo de $\T$, luego el lema~\ref{v4:arco-sin-subgrupos} muestra que la vecindad contiene sólo el carácter trivial. Las traslaciones del grupo dual prueban que todos sus puntos son abiertos.

Si $G$ es discreto, toda función $G\to\T$ es continua y los compactos de $G$ son finitos. Por ello la topología compacto-abierta coincide con la convergencia puntual, esto es, con la topología inducida por el producto $\T^G$. El subconjunto de homomorfismos es cerrado: está definido por las ecuaciones cerradas $z_e=1$ y $z_{xy}=z_xz_y$, una por cada pareja. Tychonoff hace compacto a $\T^G$, y su cerrado $\widehat G$ es compacto. No se supone aquí que el dual sea discreto; en la lectura aparece esa palabra donde debe decir «$G$ es discreto».
\end{proof}

\begin{proposicion}\label{v4:duales-modelos}
Con sus topologías usuales,
\[
 \widehat\R\simeq\R,\quad \widehat\T\simeq\Z,\quad
 \widehat\Z\simeq\T,\quad
 \widehat{\Z/N\Z}\simeq\Z/N\Z.
\]
En el primer caso el parámetro $\xi$ corresponde a $x\mapsto e^{i\xi x}$.
\end{proposicion}
\begin{proof}
Para un carácter $\chi$ de $\R$, la continuidad en cero permite escribir $\chi(t)=e^{i\theta(t)}$ para $|t|<\delta$, con $\theta$ continua y $|\theta(t)|<\pi/3$. Siempre que $s,t,s+t$ pertenezcan a ese intervalo, el número $\theta(s+t)-\theta(s)-\theta(t)$ es un múltiplo de $2\pi$ y tiene valor absoluto menor que $\pi$, por lo que es cero. Esta aditividad local y continuidad implican $\theta(t)=\xi t$: fijando $0<a<\delta/2$, la subdivisión de $a$ en partes iguales da la fórmula en los múltiplos racionales de $a$ que permanecen en el intervalo, y la densidad de esos múltiplos da la fórmula en todo él. Para un $t$ cualquiera, elija $n$ con $t/n$ en ese intervalo; entonces $\chi(t)=\chi(t/n)^n=e^{i\xi t}$.

El parámetro es único, pues si $\xi\ne\eta$ basta elegir $t=\pi/(\xi-\eta)$ para distinguir los caracteres. La continuidad del parámetro hacia el dual se sigue de
$|e^{i\xi t}-e^{i\eta t}|\le|\xi-\eta||t|$. Para la continuidad inversa en cero, si $|\xi|\ge\eps$, el punto $t=\pi/|\xi|$ pertenece a $[-\pi/\eps,\pi/\eps]$ y $|e^{i\xi t}-1|=2$. Así un error uniforme menor que $1$ en ese compacto obliga a $|\xi|<\eps$. Las traslaciones de los dos grupos concluyen el argumento.

Un carácter de $\T$ compuesto con $t\mapsto e^{it}$ es $e^{i\xi t}$ y debe tener período $2\pi$. Por ello $\xi\in\Z$, y los caracteres son $z\mapsto z^n$. Su topología es discreta por compacidad de $\T$. La aplicación $t+2\pi\Z\mapsto e^{it}$ identifica $\R/(2\pi\Z)$ con $\T$: es una biyección continua, su dominio es compacto porque es imagen de $[0,2\pi]$, y el codominio es Hausdorff; por ello es un homeomorfismo. Esto recupera también el modelo del cociente empleado en la lectura. Un carácter de $\Z$ queda determinado por $z=\chi(1)$ y vale $\chi(n)=z^n$. La evaluación en $1$ es continua; la aplicación inversa es continua uniformemente en los subconjuntos finitos de $\Z$, luego es un homeomorfismo. Finalmente, los caracteres de $\Z/N\Z$ corresponden a los $z$ con $z^N=1$, que son $e^{2\pi i k/N}$, $0\le k<N$; sus productos suman los índices módulo $N$.
\end{proof}

\section{Ejercicios sobre productos, cocientes y duales}
\begin{ejercicio}[Lectura 25, ejercicio 1]\label{ej:v4-l25-1}
Sea $G$ un grupo abeliano finito de orden $n$. Demuestre que $\widehat G\simeq G$. Explique por qué la afirmación no proporciona, en general, una identificación canónica de $G$ con su dual.
\end{ejercicio}


\begin{ejercicio}[Lectura 25, ejercicio 2]\label{ej:v4-l25-2}
Sea $A\xrightarrow{f}B\xrightarrow{g}C$ una sucesión exacta corta de grupos LCA: $f$ identifica $A$ con $\ker g$ con su topología de subespacio, y $g$ identifica $C$ con $B/f(A)$ con su topología cociente. Pruebe que
\[
 \widehat C\xrightarrow{\widehat g}\widehat B
 \xrightarrow{\widehat f}\widehat A,
 \qquad \widehat g(\eta)=\eta\circ g,\quad
 \widehat f(\chi)=\chi\circ f,
\]
es exacta corta en el mismo sentido topológico.
\end{ejercicio}


\begin{ejercicio}[Lectura 25, ejercicio 3: topologías e hipótesis necesarias]\label{ej:v4-l25-3}
Demuestre las siguientes versiones precisas de los isomorfismos de la lectura:
\begin{enumerate}
\item Si cada $G_n$ es compacto, el dual de $\prod_{n\ge1}G_n$, con topología producto, es isomorfo a $\bigoplus_{n\ge1}\widehat G_n$, dotado de la topología discreta.
\item Si cada $G_n$ es discreto, el dual de $\bigoplus_{n\ge1}G_n$, dotado de la topología discreta, es isomorfo a $\prod_{n\ge1}\widehat G_n$, dotado de la topología producto.
\end{enumerate}
Explique por qué no basta suponer que los $G_n$ sean localmente compactos sin especificar las topologías de las sumas.
\end{ejercicio}


El grupo dual ya tiene operaciones continuas y una descripción concreta en los ejemplos básicos. El capítulo siguiente probará su compacidad local mediante el espectro del álgebra de convolución. Más adelante, la inversión y Plancherel permitirán recuperar el grupo original a partir de sus frecuencias y completar el ejercicio de exactitud.

\chapter{Fourier en grupos abelianos localmente compactos}\label{cap:v4-15}
En los grupos concretos, Fourier asigna a una función un número por cada frecuencia. La definición abstracta conserva exactamente esa idea: integramos la función contra un carácter. El problema nuevo es demostrar que esas frecuencias forman el espacio correcto para estudiar el álgebra de convolución, con la topología correcta. La lectura 26 de \cite{valencia4} contiene esta identificación; desarrollaremos sus pasos, en particular la recuperación de un carácter a partir de un funcional multiplicativo.

En todo el capítulo $G$ es LCA y $\mu$ es una medida de Haar fija. Usamos $L_xf(y)=f(x^{-1}y)$, $f^*(x)=\overline{f(x^{-1})}$ y la convención
\[
 \widehat f(\chi)=\int_G f(x)\chi(x)\,d\mu(x).
\]
No hay conjugado en el carácter de la integral. Para $\chi_\xi(x)=e^{i\xi x}$ coincide con el signo positivo empleado antes; una normalización unitaria euclidiana se incorpora eligiendo la escala de Haar correspondiente. La medida dual se determinará en el capítulo siguiente, y no debe adivinarse multiplicando factores sin especificar ambas medidas.

\section{El espectro de un álgebra sin unidad}\label{lectura:v4-26}
\index{espectro!de un algebra de Banach@de un álgebra de Banach}\index{Gelfand!transformada de}
\begin{definicion}
Para un álgebra de Banach conmutativa $A$, con unidad o sin ella, un carácter es un funcional lineal continuo, multiplicativo y no nulo $\varphi:A\to\C$. Escribimos $\Delta(A)$ para el conjunto de estos caracteres, con la topología débil estrella heredada de $A^*$: convergencia de $\varphi_\alpha$ significa convergencia de $\varphi_\alpha(a)$ para cada $a\in A$.
\end{definicion}
La palabra «carácter» se aplica aquí a un álgebra, mientras que en el capítulo anterior se aplicaba a un grupo. Los dominios y las propiedades son diferentes; Fourier establecerá una correspondencia entre ambas nociones.

\begin{proposicion}\label{v4:espectro-no-unitario}
Todo $\varphi\in\Delta(A)$ satisface $\norm\varphi\le1$. El conjunto $\Delta(A)\cup\{0\}$ es compacto Hausdorff en la topología débil estrella, y $\Delta(A)$ es localmente compacto Hausdorff.
\end{proposicion}
\begin{proof}
Para $a\in A$ y $n\ge1$, multiplicatividad y continuidad dan
\[
 |\varphi(a)|^n=|\varphi(a^n)|\le\norm\varphi\norm{a^n}
 \le\norm\varphi\norm a^n.
\]
Si $a\ne0$, tomar raíces $n$-ésimas y dejar $n\to\infty$ da $|\varphi(a)|\le\norm a$; para $a=0$ es inmediato. Así todos los caracteres están en la bola unidad del dual, compacta débil estrella por Banach--Alaoglu, desarrollado en el volumen III, capítulo 9, «Compacidad débil y reflexividad».

Dentro de esa bola, las ecuaciones
\[
 \varphi(ab)=\varphi(a)\varphi(b)\qquad(a,b\in A)
\]
son condiciones cerradas para la convergencia puntual. Sus soluciones son exactamente los caracteres y el funcional cero. Por tanto $\Delta(A)\cup\{0\}$ es cerrado en una bola compacta Hausdorff, y es compacto Hausdorff. Al quitar el punto cerrado $0$ queda un abierto de un compacto Hausdorff, que es localmente compacto: dado un punto, la regularidad permite elegir una vecindad cuya clausura evite $0$ y sea compacta. Si $A$ tiene unidad, los caracteres satisfacen $\varphi(1)=1$; entonces el punto cero es incluso aislado en este compacto ampliado.
\end{proof}

\begin{teorema}\label{v4:gelfand-cero-infinito}
La transformada de Gelfand $a\mapsto\widehat a$, $\widehat a(\varphi)=\varphi(a)$, es un homomorfismo contractivo de $A$ en $C_0(\Delta(A))$.
\end{teorema}
\begin{proof}
La definición de topología débil estrella hace continua a cada evaluación $\widehat a$. Linealidad y multiplicatividad se verifican evaluando en cada $\varphi$, y la proposición anterior da $\norm{\widehat a}_\infty\le\norm a$. Para $\eps>0$, el conjunto
\[
 K_\eps=\{\varphi\in\Delta(A)\cup\{0\}:|\varphi(a)|\ge\eps\}
\]
es cerrado en el compacto ampliado y no contiene al funcional cero. Es, por tanto, un compacto de $\Delta(A)$. Esa es exactamente la condición de que $\widehat a$ se anule en el infinito. No basta decir que $K_\eps$ es cerrado en un espacio localmente compacto: se necesita el compacto ampliado para justificar su compacidad.
\end{proof}

\section{La integral que asigna una frecuencia a una función}
\index{Fourier!transformada en grupos LCA}
\begin{proposicion}\label{v4:fourier-identidades-lca}
Para $f,g\in L^1(G)$ y $\chi\in\widehat G$, la transformada satisface
\[
 |\widehat f(\chi)|\le\norm f_1,\qquad
 \widehat{f*g}=\widehat f\,\widehat g,\qquad
 \widehat{f^*}=\overline{\widehat f},
\]
\[
 \widehat{L_xf}(\chi)=\chi(x)\widehat f(\chi),\qquad
 \widehat{\eta f}(\chi)=\widehat f(\eta\chi)\quad(\eta\in\widehat G).
\]
Además, $\Phi_\chi:f\mapsto\widehat f(\chi)$ es un carácter de $L^1(G)$.
\end{proposicion}
\begin{proof}
El módulo del integrando es $|f|$, por lo que la integral existe y satisface la primera cota. Para la convolución, la integral absoluta doble queda acotada por
\[
 \int_G\int_G|f(y)||g(y^{-1}x)|\,d\mu(x)d\mu(y)
 =\norm f_1\norm g_1.
\]
Fubini y el cambio $x=yz$ dan
\[
 \widehat{f*g}(\chi)=\int_G\int_Gf(y)g(z)\chi(yz)\,d\mu(z)d\mu(y)
 =\widehat f(\chi)\widehat g(\chi).
\]
Los cambios de variable en un grupo LCA no introducen función modular: es unimodular. En particular, $x\mapsto x^{-1}$ preserva Haar, y
\[
 \widehat{f^*}(\chi)=\int_G\overline{f(y)}\chi(y^{-1})\,d\mu(y)
 =\overline{\widehat f(\chi)}.
\]
El cambio $y=xz$ en la traslación da el factor $\chi(x)$; la identidad de modulación se sigue de multiplicar los dos caracteres dentro de la integral.

Queda verificar que el funcional no sea cero. Elija $h\in C_c(G)$ no negativa con $\int h=1$, posible por la positividad de Haar en abiertos no vacíos. Para $f=h\overline\chi$ se tiene $\widehat f(\chi)=1$. La linealidad, continuidad y multiplicatividad ya probadas completan la verificación.
\end{proof}

En grupos LCA generales, Haar no tiene por qué ser sigmafinita en todo $G$. Las aplicaciones de Fubini anteriores se justifican con la versión para Haar y funciones integrables del capítulo de convolución: cada función de $L^1$ se aproxima por funciones de soporte compacto y queda concentrada, salvo un conjunto nulo, en una unión numerable de compactos. Se reducen así las integrales pertinentes a piezas sigmafinitas; no se está añadiendo una hipótesis de segunda numerabilidad al grupo.

\begin{lema}[Una integral con valores en $L^1$]\label{v4:convolucion-funcional}
Si $f\in C_c(G)$, $g\in L^1(G)$ y $\varphi\in(L^1(G))^*$, entonces
\[
 \varphi(f*g)=\int_G f(x)\varphi(L_xg)\,d\mu(x).
\]
\end{lema}
\begin{proof}
La aplicación $x\mapsto L_xg$ es continua en norma $L^1$ por la continuidad de traslaciones ya demostrada. En el compacto $K=\supp f$ podemos cubrir por un número finito de conjuntos en los que $\norm{L_xg-L_{x_j}g}_1<\eps$. Refinando la cubierta en una partición boreliana finita $E_j$ de $K$, aproximamos la convolución por
\[
 v_\eps=\sum_j\left(\int_{E_j}f\,d\mu\right)L_{x_j}g.
\]
La desigualdad de la norma integral, obtenida por Tonelli, da
$\norm{f*g-v_\eps}_1\le\eps\norm f_1$. A su vez, la diferencia entre la integral del enunciado y $\varphi(v_\eps)$ tiene módulo a lo sumo $\norm\varphi\eps\norm f_1$. Al dejar $\eps\downarrow0$ y usar continuidad de $\varphi$ obtenemos la identidad. Esta prueba explicita por qué puede pasarse el funcional a través de la integral, sin presuponer una teoría general de integración vectorial.
\end{proof}

\section{Todo carácter del álgebra procede de una frecuencia}
\begin{teorema}\label{v4:espectro-luno-biyeccion}
La aplicación $\chi\mapsto\Phi_\chi$ es una biyección de $\widehat G$ sobre $\Delta(L^1(G))$.
\end{teorema}
\begin{proof}
Para la inyectividad, suponga $\Phi_\chi=\Phi_\eta$. Si $\chi(x_0)-\eta(x_0)\ne0$, la continuidad permite elegir un escalar $c$ de módulo uno y una vecindad relativamente compacta $U$ de $x_0$ en la que $\operatorname{Re}(c(\chi-\eta))>a>0$. Elija $h\in C_c(U)$ no negativa, no nula. Entonces
\[
 \operatorname{Re}\int_G ch(x)(\chi(x)-\eta(x))\,d\mu(x)>0,
\]
porque $\int h>0$. Esto contradice la igualdad de los dos funcionales aplicada a $ch$. Así los caracteres coinciden en todo punto.

Para la sobreyectividad, sea $\varphi\in\Delta(L^1(G))$. Por densidad de $C_c(G)$ y continuidad de $\varphi$, existe $g\in C_c(G)$ con $\varphi(g)\ne0$. Definimos
\[
 \chi(x)=\frac{\varphi(L_xg)}{\varphi(g)}.
\]
Es continua por continuidad de las traslaciones en $L^1$, y $\chi(e)=1$. En un grupo abeliano,
$L_xg*L_yg=L_{xy}(g*g)=(L_{xy}g)*g$; esto se comprueba en la integral de convolución cambiando la variable por una traslación. Aplicar $\varphi$ y dividir por $\varphi(g)^2$ da $\chi(x)\chi(y)=\chi(xy)$.

Es importante verificar el módulo. La cota del funcional y la isometría de las traslaciones dan
\[
 |\chi(x)|\le\frac{\norm\varphi\norm g_1}{|\varphi(g)|}=:C
 \qquad(x\in G).
\]
Por multiplicatividad, $\chi(x)\chi(x^{-1})=1$, de modo que ningún valor es cero. Para cada entero positivo $n$, $|\chi(x)|^n=|\chi(x^n)|\le C$; al tomar raíces se obtiene $|\chi(x)|\le1$. Aplicar lo mismo a $x^{-1}$ da la desigualdad contraria. Así $|\chi(x)|=1$ y $\chi\in\widehat G$.

Para $f\in C_c(G)$, el lema anterior y la multiplicatividad dan
\[
 \int_Gf(x)\chi(x)\,d\mu(x)
 =\frac{\varphi(f*g)}{\varphi(g)}=\varphi(f).
\]
Ambos lados son funcionales continuos de $f$ en norma $L^1$, por lo que la densidad extiende la igualdad a todo $L^1(G)$. Esto demuestra sobreyectividad. El argumento de las potencias evita inferencias no justificadas sobre convergencia puntual de redes de Dirac.
\end{proof}

\section{Por qué coinciden las dos topologías}
La convergencia débil estrella controla las integrales contra cada función de $L^1$. La topología compacto-abierta controla simultáneamente los valores de un carácter en cada compacto de $G$. La implicación de las integrales al control uniforme requiere el siguiente argumento cuantitativo.

\begin{lema}\label{v4:dual-vecindades-integrales}
Dados $\chi_0\in\widehat G$, un compacto $K\subset G$ y $\eps>0$, existen $f_0,\ldots,f_m\in L^1(G)$ y $\delta>0$ tales que
\[
 |\widehat f_j(\chi)-\widehat f_j(\chi_0)|<\delta\quad(0\le j\le m)
 \quad\Longrightarrow\quad
 \sup_K|\chi-\chi_0|<\eps.
\]
\end{lema}
\begin{proof}
Basta tratar $0<\eps\le1$ y $\chi_0=1$. El caso general se reduce a éste reemplazando $\chi$ por $\chi\overline{\chi_0}$ y cada función de prueba $h$ por $h\overline{\chi_0}$.

Elija $h\in C_c(G)$ no negativa con $\int h=1$. La continuidad en $L^1$ de las traslaciones permite elegir una vecindad $U$ de $e$ con
$\norm{L_uh-h}_1<\eps/8$ para $u\in U$. Cubra $K$ por $x_1U,\ldots,x_mU$ y tome $f_0=h$, $f_j=L_{x_j}h$. Sus transformadas en el carácter trivial valen $1$. Elija $\delta=\eps/16$ y suponga las desigualdades del enunciado. Escribiendo $a=\widehat h(\chi)$ tenemos $|a-1|<\delta$, de donde $|a|>1/2$. Además,
\[
 |\chi(x_j)-1|\,|a|
 \le|\chi(x_j)a-1|+|1-a|<2\delta,
\]
por lo que $|\chi(x_j)-1|<4\delta=\eps/4$. Para $u\in U$,
\[
 |\chi(u)-1|\,|a|
 =|\widehat{L_uh}(\chi)-\widehat h(\chi)|
 \le\norm{L_uh-h}_1<\eps/8,
\]
y $|\chi(u)-1|<\eps/4$. Si $x=x_ju\in K$, entonces
$|\chi(x)-1|\le|\chi(x_j)-1|+|\chi(u)-1|<\eps/2$. El supremo es a lo sumo $\eps/2<\eps$, como se quería.
\end{proof}

\begin{teorema}\label{v4:dual-espectro-homeomorfismo}
La biyección $\widehat G\to\Delta(L^1(G))$ del teorema~\ref{v4:espectro-luno-biyeccion} es un homeomorfismo. En particular, $\widehat G$ es LCA y $\widehat f\in C_0(\widehat G)$ para toda $f\in L^1(G)$.
\end{teorema}
\begin{proof}
Si $\chi_\alpha\to\chi$ uniformemente en compactos, dado $f\in L^1$ elija $g\in C_c$ con $\norm{f-g}_1<\eps/4$. Entonces
\[
 |\widehat f(\chi_\alpha)-\widehat f(\chi)|
 \le2\norm{f-g}_1+\norm g_1\sup_{\supp g}|\chi_\alpha-\chi|.
\]
El último término acaba siendo menor que $\eps/2$; si $g=0$ se omite. Esto prueba continuidad hacia la topología débil estrella. Para la continuidad inversa, toda vecindad definida por un compacto y un error contiene, por el lema anterior, una vecindad dada por un número finito de evaluaciones del espectro. Las intersecciones finitas permiten tratar vecindades generales.

El espectro es localmente compacto Hausdorff por la proposición~\ref{v4:espectro-no-unitario}; el dual ya es un grupo topológico abeliano por el capítulo anterior. Finalmente, bajo este homeomorfismo la transformada de Fourier es precisamente la de Gelfand de $L^1(G)$, y el teorema~\ref{v4:gelfand-cero-infinito} demuestra la propiedad $C_0$.
\end{proof}

Esta última conclusión es el teorema de Riemann--Lebesgue para un grupo LCA. Cuando el grupo dual es discreto, significa que fuera de un conjunto finito las frecuencias tienen amplitud menor que cualquier tolerancia fijada. Cuando no es discreto, la palabra «finito» se reemplaza por «compacto».

\section{Ejercicios complementarios de interpretación}
La lectura 26 no contiene una lista numerada de ejercicios; sus verificaciones propuestas se han incluido en las pruebas anteriores. Los siguientes ejercicios son propios de esta edición y ayudan a distinguir las normalizaciones y las topologías.

\begin{ejercicio}[Ampliación: dos escalas de Haar]\label{ej:v4-ampl-15-1}
Si se reemplaza $\mu$ por $c\mu$, con $c>0$, determine cómo cambian la norma $L^1$, la convolución y la transformada de Fourier. Compruebe que la identidad de convolución sigue siendo válida con esas tres modificaciones simultáneas.
\end{ejercicio}


\begin{ejercicio}[Ampliación: traslación, modulación y signo]\label{ej:v4-ampl-15-2}
En $G=\R$ con Lebesgue y $\chi_\xi(x)=e^{i\xi x}$, calcule la transformada de $e^{i\eta x}f(x-a)$ en términos de $\widehat f$. Determine también la transformada de $\overline{f(-x)}$.
\end{ejercicio}


\begin{ejercicio}[Ampliación: una frecuencia que escapa al infinito]\label{ej:v4-ampl-15-3}
En $\widehat\R\simeq\R$, considere $\chi_n(x)=e^{inx}$. Pruebe que los funcionales $\Phi_{\chi_n}$ convergen débil estrella al funcional cero, aunque ninguno es cero, y que $(\chi_n)$ no converge en $\widehat\R$.
\end{ejercicio}


La transformada ya es un homomorfismo contractivo hacia funciones que se anulan en el infinito. Falta probar que distingue funciones y construir una transformación unitaria en $L^2$. El álgebra de operadores de convolución permitirá resolver ambas cuestiones sin suponer de antemano la dualidad de Pontryagin.

\chapter{Álgebras del grupo, inversión y dualidad}\label{cap:v4-16}
La convolución puede estudiarse como una operación entre funciones o como un operador sobre $L^2(G)$. La segunda perspectiva permite completar $L^1(G)$ con una norma distinta y recuperar todas las funciones continuas que se anulan en el infinito del grupo dual. Después construiremos una medida de Haar dual y demostraremos la dualidad de Pontryagin. Así quedará resuelto, con sus dependencias completas, el ejercicio~\ref{ej:v4-l25-2} sobre sucesiones exactas.

Las lecturas 27 y 28 de \cite{valencia4} contienen el núcleo de la construcción. Desarrollamos sus verificaciones, corregimos dos problemas de soporte y de signo en el último teorema, y añadimos el tramo de inversión y dualidad. Como referencia general pueden consultarse los teoremas 4.32 y 4.40 de \cite{folland2016}. Todo el capítulo supone que $G$ es LCA y que se ha fijado una medida de Haar $\mu$. Escribimos
\[
 f^*(x)=\overline{f(x^{-1})},\qquad
 \widehat f(\chi)=\int_G f(x)\chi(x)\,d\mu(x).
\]
La unimodularidad permite estas fórmulas sin factores adicionales. Usaremos tres resultados de los volúmenes anteriores: la representación positiva sobre $C_c(X)$ del volumen II, capítulo 16, «El teorema de representación de Riesz»; Stone--Weierstrass del volumen I, capítulo 13, «Extensión y aproximación de funciones continuas»; y Gelfand del volumen III, capítulo 13, «Álgebras $C^*$ y cálculo funcional continuo», incluido allí el ejercicio sobre álgebras no unitarias de la lectura 26. Recordaremos sus conclusiones justo donde intervengan. La prueba de Stone--Weierstrass por interpolaciones y máximos y mínimos utiliza sólo compacidad y separación Hausdorff; por ello vale en los compactos no necesariamente métricos que aparecerán aquí. No se presupone la dualidad de Pontryagin.

En los intercambios de integrales no suponemos segunda numerabilidad. Las identidades se prueban primero para funciones de $C_c$: sobre los productos de compactos pertinentes, Stone--Weierstrass aproxima uniformemente los núcleos continuos por sumas de productos de funciones de cada variable, para las cuales el intercambio es inmediato. Después se pasa a funciones $L^1$ o $L^2$ por sus aproximaciones en norma y las cotas absolutas indicadas en cada prueba. Cada una de esas funciones admite, salvo un conjunto nulo, un representante concentrado en una unión numerable de compactos: tome aproximaciones por $C_c$ y una subsucesión que converja casi en todas partes. Para un número finito de funciones, todos esos compactos están contenidos en un subgrupo abierto sigma compacto: se añade una vecindad compacta de $e$ y se toman todos los productos finitos de esos compactos y sus inversos. En ese subgrupo Haar es sigmafinita y se aplican Tonelli y Fubini; la aproximación previa por núcleos continuos resuelve también la cuestión de las sigma álgebras producto en espacios no metrizables. Esta reducción, compatible con el capítulo de Haar, justifica las aplicaciones siguientes sin restringir el grupo completo.

\section{La representación regular y su norma}\label{lectura:v4-27}
\index{representacion@representación!regular}\index{algebra@álgebra!reducida de un grupo}
Para $f\in L^1(G)$ y $u\in L^2(G)$ queremos definir $\lambda(f)u=f*u$. Conviene justificar primero que la integral produce un elemento de $L^2$.

\begin{proposicion}\label{v4:regular-l2}
La fórmula $\lambda(f)u=f*u$ define un operador lineal acotado en $L^2(G)$, con
\[
 \norm{\lambda(f)}\le\norm f_1,
 \qquad
 \ip{\lambda(f)u}{v}
 =\int_G f(y)\ip{L_yu}{v}\,d\mu(y).
\]
Además, $\lambda(f*g)=\lambda(f)\lambda(g)$ y $\lambda(f^*)=\lambda(f)^*$.
\end{proposicion}
\begin{proof}
Si $f=0$ no hay nada que demostrar. La desigualdad de Cauchy--Schwarz para la medida $|f(y)|\,d\mu(y)$ da, donde la integral de la derecha es finita,
\[
 \left|\int_G f(y)u(y^{-1}x)\,d\mu(y)\right|^2
 \le\norm f_1\int_G|f(y)|\,|u(y^{-1}x)|^2\,d\mu(y).
\]
Por Tonelli e invariancia de Haar, la integral en $x$ del último miembro es $\norm f_1^2\norm u_2^2$. Por tanto la convolución existe absolutamente para casi todo $x$, pertenece a $L^2$ y satisface la cota afirmada. Los representantes distintos casi en todas partes dan el mismo elemento de $L^2$, por la misma estimación.

Para $u,v\in L^2$,
\[
 \int_G|f(y)|\int_G|u(y^{-1}x)v(x)|\,d\mu(x)\,d\mu(y)
 \le\norm f_1\norm u_2\norm v_2.
\]
Fubini permite intercambiar las integrales y obtiene la identidad del producto interno. Ésta también construye $\lambda(f)$ mediante el teorema de representación de Riesz del espacio de Hilbert: el funcional conjugado lineal en $v$ está acotado por la cantidad indicada. Ambas construcciones coinciden.

Si $u\in C_c(G)$, las identidades de asociatividad en $L^1$ muestran que
$(f*g)*u=f*(g*u)$. Ambos miembros también pertenecen a $L^2$ por la cota anterior. La densidad de $C_c$ en $L^2$ y la acotación de los operadores extienden la igualdad a todo $u\in L^2$.
Finalmente, $L_y^*=L_{y^{-1}}$, y el producto interno es lineal en su primera variable. Por ello
\[
 \ip{\lambda(f)u}{v}
 =\int_G f(y)\ip{u}{L_{y^{-1}}v}\,d\mu(y)
 =\ip{u}{\lambda(f^*)v}.
\]
En la última igualdad se usa inversión en Haar y la conjugación del escalar que entra en la segunda variable. Esto prueba la identidad para el adjunto.
\end{proof}

\begin{proposicion}\label{v4:regular-fiel}
El homomorfismo $\lambda:L^1(G)\to\mathcal B(L^2(G))$ es inyectivo.
\end{proposicion}
\begin{proof}
Si $\lambda(f)=0$, entonces $f*\xi=0$ como elemento de $L^2$ para toda $\xi\in C_c(G)$. Esa convolución también es un elemento de $L^1$, y es la misma función casi en todas partes, de modo que es cero en $L^1$. Elija una identidad aproximada $(\xi_U)$ en $C_c(G)$, no negativa y de integral uno. La convergencia $f*\xi_U\to f$ en $L^1$ obliga a $f=0$. No hace falta extraer subredes con convergencia puntual: la convergencia en $L^1$ ya resuelve el argumento.
\end{proof}

\begin{definicion}
El \emph{álgebra $C^*$ reducida} de $G$ es
\[
 C_r^*(G)=\overline{\lambda(L^1(G))}^{\,\norm{\cdot}_{\mathcal B(L^2)}}.
\]
\end{definicion}
Es un álgebra cerrada de operadores, estable bajo adjuntos, y hereda la identidad $\norm{A^*A}=\norm A^2$. Es conmutativa porque la convolución lo es. La lectura la denota $C^*(G)$; usamos inicialmente el subíndice $r$ para hacer visible la construcción. Al final de la siguiente sección justificaremos que, en grupos abelianos, coincide con el álgebra universal que también suele recibir ese nombre.

\section{El espectro del álgebra del grupo}
Recordemos el teorema de Gelfand para un álgebra $C^*$ conmutativa $B$: su transformada $A\mapsto\widetilde A$, dada por $\widetilde A(\rho)=\rho(A)$, es un isomorfismo isométrico y preservador del adjunto entre $B$ y $C_0(\Delta(B))$. En particular, si $B\ne0$, su espectro no es vacío y sus caracteres tienen norma uno. Se incluyen aquí álgebras sin unidad.

\begin{lema}[Un criterio de homeomorfismo de espectros]\label{v4:espectros-homeomorfismo}
Sean $A,B$ álgebras de Banach conmutativas y $T:A\to B$ un homomorfismo continuo. Suponga que cada carácter de $B$ tiene restricción $\rho\circ T\ne0$ y que $T^*: \Delta(B)\to\Delta(A)$, $\rho\mapsto\rho\circ T$, es biyectivo. Entonces $T^*$ es un homeomorfismo.
\end{lema}
\begin{proof}
La continuidad se sigue de $T^*(\rho)(a)=\rho(Ta)$: cada evaluación en $a$ se transforma en una evaluación en $Ta$. Para controlar la inversa añadimos formalmente una unidad a ambas álgebras, incluso si ya tenían una. Los caracteres del álgebra unitalizada $A^+=A\oplus\C$ son
\[
 (a,z)\longmapsto\phi(a)+z\quad(\phi\in\Delta(A)),
 \qquad (a,z)\longmapsto z.
\]
En efecto, la restricción a $A$ es un carácter o es cero, y el valor en la nueva unidad debe ser uno. El espectro de $A^+$ es compacto Hausdorff: es un subconjunto cerrado de la bola dual, definido por multiplicatividad y valor uno en la unidad; Banach--Alaoglu da compacidad. Lo mismo ocurre para $B^+$.

El homomorfismo $T^+(a,z)=(Ta,z)$ induce una aplicación continua entre esos compactos. Las hipótesis dicen exactamente que es una biyección, incluyendo el nuevo carácter de restricción cero. Una biyección continua de un compacto a un Hausdorff es un homeomorfismo. Al quitar esos dos caracteres nuevos quedan las topologías originales de $\Delta(B)$ y $\Delta(A)$ y el homeomorfismo buscado.
\end{proof}

\begin{teorema}\label{v4:cstar-fourier}
Cada carácter $\chi\in\widehat G$ determina un único carácter de $C_r^*(G)$ cuyo valor en $\lambda(f)$ es $\widehat f(\chi)$. Esto identifica homeomórficamente $\Delta(C_r^*(G))$ con $\widehat G$, y proporciona un isomorfismo isométrico de álgebras $C^*$
\[
 C_r^*(G)\simeq C_0(\widehat G),\qquad
 \norm{\lambda(f)}=\norm{\widehat f}_\infty.
\]
En particular, la transformada de Fourier de $L^1(G)$ es inyectiva.
\end{teorema}
\begin{proof}
Un carácter $\rho$ de $C_r^*(G)$ no puede anularse en toda $\lambda(L^1)$, porque este subespacio es denso y $\rho$ es continuo. Por el capítulo~\ref{cap:v4-15}, la restricción es $\widehat f(\chi)$ para algún $\chi\in\widehat G$. Dos caracteres con la misma restricción coinciden por densidad. Falta demostrar que se obtienen todos los $\chi$.

Como $C_r^*(G)\ne0$, elija un carácter $\rho_0$ y sea $\chi_0$ su restricción. Para $\eta\in\widehat G$ la modulación $M_\eta u=\eta u$ es unitaria en $L^2$ y un cálculo directo da
\[
 M_\eta\lambda(f)M_\eta^{-1}=\lambda(\eta f).
\]
En efecto, $\eta(x)\overline{\eta(y^{-1}x)}=\eta(y)$ en el integrando. Para un carácter arbitrario $\chi$, tome $\eta=\chi\chi_0^{-1}$. Entonces
\[
 |\widehat f(\chi)|
 =|\widehat{\eta f}(\chi_0)|
 =|\rho_0(\lambda(\eta f))|
 \le\norm{\lambda(\eta f)}=\norm{\lambda(f)}.
\]
Así $\lambda(f)\mapsto\widehat f(\chi)$ es un funcional acotado en el subespacio denso $\lambda(L^1)$. Se extiende de modo único a $C_r^*(G)$. La multiplicatividad pasa al límite porque la multiplicación es continua; la extensión es no nula, pues existe $f\in C_c$ con $\widehat f(\chi)\ne0$. El lema anterior, aplicado a $T=\lambda$, prueba el homeomorfismo de espectros. El teorema de Gelfand produce ahora el isomorfismo con $C_0(\widehat G)$ y la igualdad de normas. Si $\widehat f=0$, entonces $\lambda(f)=0$ y la proposición~\ref{v4:regular-fiel} da $f=0$.
\end{proof}

\begin{corolario}[Los caracteres separan puntos]\label{v4:caracteres-separan}
Si $x\ne y$ en $G$, existe $\chi\in\widehat G$ tal que $\chi(x)\ne\chi(y)$.
\end{corolario}
\begin{proof}
Basta tratar $a=y^{-1}x\ne e$. Si $\chi(a)=1$ para todos los caracteres, la fórmula $\widehat{L_af}(\chi)=\chi(a)\widehat f(\chi)$ y la inyectividad anterior implican $L_af=f$ en $L^1$ para toda $f\in C_c$. Ambas funciones son continuas; su igualdad casi en todas partes es igualdad en todos los puntos, porque todo abierto no vacío tiene medida de Haar positiva. Pero existe una función $f\in C_c$ con $f(e)=1$ y $f(a^{-1})=0$: se toma una función de corte apoyada en una vecindad relativamente compacta de $e$ que no contiene $a^{-1}$. Evaluar $L_af=f$ en $e$ contradice esos valores.
\end{proof}

\begin{observacion}[Por qué las dos álgebras de grupo coinciden aquí]
Si $\pi$ es cualquier representación unitaria continua de $G$, su forma integrada satisface $\norm{\pi(f)}\le\norm f_1$ y es un homomorfismo que preserva adjuntos. La clausura $B_\pi$ de su imagen es una $C^*$-álgebra conmutativa. Cada carácter de $B_\pi$, restringido a $L^1$, es no nulo por densidad y corresponde a un carácter de $G$. Por Gelfand,
\[
 \norm{\pi(f)}
 =\sup_{\rho\in\Delta(B_\pi)}|\rho(\pi(f))|
 \le\sup_{\chi\in\widehat G}|\widehat f(\chi)|
 =\norm{\lambda(f)}.
\]
Si $B_\pi=0$ la desigualdad es inmediata. El supremo sobre todas las representaciones unitarias, que define la norma universal, coincide por tanto con la norma regular, pues $\lambda$ es una de esas representaciones integradas. En este capítulo podemos escribir $C^*(G)=C_r^*(G)$. La prueba utiliza la conmutatividad y no afirma esa igualdad para grupos arbitrarios.
\end{observacion}

\section{Un espacio que controla ambos lados de Fourier}\label{lectura:v4-28}
\index{convolucion@convolución!de funciones de cuadrado integrable}
\begin{lema}\label{v4:l2-convolucion-cero}
Si $u,v\in L^2(G)$, entonces $u*v$ existe absolutamente en cada punto, pertenece a $C_0(G)$ y satisface
\[
 \norm{u*v}_\infty\le\norm u_2\norm v_2.
\]
\end{lema}
\begin{proof}
La inversión y la traslación preservan la norma $L^2$. Por tanto
\[
 (u*v)(x)=\ip{u}{L_xv^*}
\]
existe para cada $x$ por Cauchy--Schwarz, y satisface la cota uniforme. Si se cambia un representante en un conjunto nulo, su traslación o inversión sigue siendo nula; la integral no depende del representante. La continuidad fuerte de las traslaciones en $L^2$ implica continuidad de esta función de $x$.

Elija $u_n,v_n\in C_c$ con $u_n\to u$ y $v_n\to v$ en $L^2$. Entonces
\[
 \norm{u_n*v_n-u*v}_\infty
 \le\norm{u_n-u}_2\norm{v_n}_2
       +\norm u_2\norm{v_n-v}_2\longrightarrow0.
\]
Cada convolución $u_n*v_n$ tiene soporte compacto, contenido en $\supp(u_n)\supp(v_n)$. Como $C_0(G)$ es cerrado para la norma uniforme, el límite pertenece a $C_0(G)$.
\end{proof}

Pongamos $D=L^1(G)\cap C_0(G)$. La expresión
\[
 \norm f_{0,*}=\max\{\norm f_\infty,\norm{\widehat f}_\infty\}
 \qquad(f\in D)
\]
es una norma: homogeneidad y desigualdad triangular se deducen de las dos normas uniformes, y $\norm f_{0,*}=0$ implica $f=0$ por el primer término. Su completación será $C_0^*(G)$. El símbolo no afirma que este espacio sea un álgebra $C^*$ con esa norma.

\index{espacio!C0 estrella@$C_0^*(G)$}
Una realización concreta de la completación evita confundir sus elementos:
\[
 E=\overline{\{(f,\widehat f):f\in D\}}
 ^{\,C_0(G)\oplus C_0(\widehat G)},\qquad
 \norm{(u,h)}=\max\{\norm u_\infty,\norm h_\infty\}.
\]
El espacio producto es Banach, y la clausura indicada es precisamente la completación isométrica. Demostraremos que cada una de las dos coordenadas determina la otra.

\begin{proposicion}\label{v4:grafo-fourier-inyectivo}
Las proyecciones $E\to C_0(G)$ y $E\to C_0(\widehat G)$ son lineales, contractivas e inyectivas. Por ello podemos identificar $E=C_0^*(G)$ con un espacio de funciones $u\in C_0(G)$ y escribir $\widehat u=h$ para su segunda coordenada. Esta notación no presupone que $u$ sea integrable.
\end{proposicion}
\begin{proof}
La linealidad y las cotas son inmediatas de la norma máxima. Sea $(u,h)\in E$ y elija $f_n\in D$ con $f_n\to u$ uniformemente y $\widehat f_n\to h$ uniformemente. Por el teorema~\ref{v4:cstar-fourier}, existe $A\in C_r^*(G)$ con símbolo $h$, y $\lambda(f_n)\to A$ en norma de operadores.

Para $g,\psi\in C_c(G)$, tenemos simultáneamente
\[
 \ip{\lambda(f_n)g}{\psi}\longrightarrow\ip{Ag}{\psi},
 \qquad
 \int_G(f_n*g)\overline\psi\,d\mu
 \longrightarrow\int_G(u*g)\overline\psi\,d\mu.
\]
La segunda convergencia se sigue de
$\norm{(f_n-u)*g}_\infty\le\norm{f_n-u}_\infty\norm g_1$
y de $\psi\in L^1$. Ambas expresiones de partida son iguales, así que $Ag=u*g$ casi en todas partes. Además, $u*g\in C_0$: se aproxima $u$ uniformemente por funciones de $C_c$ y se usa la misma cota. La unicidad local de Riesz--Markov dice inicialmente que las dos funciones coinciden casi en todas partes en cada abierto relativamente compacto. Aquí eso implica igualdad global: $Ag\in L^2$ está concentrada, salvo un conjunto nulo, en una unión numerable de compactos por aproximación mediante $C_c$; la función $u*g\in C_0$ se anula fuera de $\bigcup_{n\ge1}\{|u*g|\ge1/n\}$, también una unión numerable de compactos. Cubrir cada uno de esos compactos por un número finito de abiertos relativamente compactos y tomar la unión numerable de los conjuntos excepcionales completa la justificación. No se usa una afirmación irrestricta de igualdad global para cualesquiera funciones sólo localmente integrables.

Si $u=0$, entonces $Ag=0$ para todo $g\in C_c$, y por densidad $A=0$, luego $h=0$. Si $h=0$, entonces $A=0$ y $u*g=0$ casi en todas partes para todo $g\in C_c$. Por continuidad y soporte pleno de Haar, esa igualdad vale en todos los puntos. Para una identidad aproximada $\xi_U\in C_c$, se tiene $u*\xi_U\to u$ uniformemente, por lo que $u=0$. Para verificar esta última convergencia, obsérvese primero que las traslaciones de una función de $C_c$ son continuas en norma uniforme; se extiende a $C_0$ por aproximación uniforme, y luego se promedia el error de traslación contra $\xi_U$. Así ambas proyecciones tienen núcleo cero.
\end{proof}

La realización por pares muestra además que $E$ es estable bajo traslación, modulación e involución: esas operaciones preservan $D$, actúan isométricamente sobre las correspondientes coordenadas uniformes, y por continuidad pasan a la clausura. En particular,
\[
 \widehat{L_xu}(\chi)=\chi(x)\widehat u(\chi),\quad
 \widehat{\eta u}(\chi)=\widehat u(\eta\chi),\quad
 \widehat{u^*}=\overline{\widehat u}
 \qquad(u\in E).
\]

\section{Positividad, regularización e identidades aproximadas}
\index{identidad!aproximada con transformada positiva}
\begin{lema}\label{v4:dirac-positivo}
Existe una identidad aproximada $(\xi_U)$ en $C_c(G)$ tal que $\xi_U\ge0$, $\int\xi_U=1$, $\supp\xi_U\subset U$ y $\widehat\xi_U\ge0$.
\end{lema}
\begin{proof}
Los índices son las vecindades $U$ de $e$, ordenadas por inclusión inversa. Elija una vecindad simétrica $W$ con $WW\subset U$. Una función de corte no negativa y no nula, apoyada en $W$, puede reemplazarse por la suma con su reflejo $x\mapsto h(x^{-1})$. Normalizando su integral obtenemos $h_U\in C_c$, $h_U\ge0$, $h_U=h_U^*$ y $\int h_U=1$. Defina $\xi_U=h_U*h_U^*$. El integrando de la convolución es no negativo, su integral total es uno y su soporte está contenido en $WW\subset U$. Además,
$\widehat\xi_U=|\widehat h_U|^2\ge0$.
La masa está enteramente en $U$, de modo que estas funciones forman una identidad aproximada.
\end{proof}

\begin{lema}[Positividad en la identidad]\label{v4:positividad-identidad}
Si $u\in E$ y $\widehat u$ es real, entonces $u(e)\in\R$. Si $\widehat u\ge0$, entonces $u(e)\ge0$.
\end{lema}
\begin{proof}
En el primer caso $\widehat{u^*}=\overline{\widehat u}=\widehat u$. La inyectividad de la segunda proyección da $u=u^*$; al evaluar en $e$, se obtiene $u(e)=\overline{u(e)}$.

En el segundo caso, el operador $A\in C_r^*(G)$ de símbolo $\widehat u$ es positivo: por Gelfand, la función $\sqrt{\widehat u}\in C_0(\widehat G)$ corresponde a un operador $B$ y $A=B^*B$. Para las funciones $h_U$ del lema anterior,
\[
 0\le\ip{Ah_U}{h_U}
 =\int_G(u*h_U)(x)\overline{h_U(x)}\,d\mu(x)
 =(u*h_U*h_U^*)(e)=(u*\xi_U)(e).
\]
La identidad intermedia se obtiene por inversión en Haar; todas las integrales son absolutamente convergentes porque $h_U$ tiene soporte compacto y $u$ es acotada. Al pasar al límite uniforme $u*\xi_U\to u$, resulta $u(e)\ge0$.
\end{proof}

\begin{lema}[Regularizar un operador]\label{v4:regularizacion}
Se cumplen las siguientes propiedades.
\begin{enumerate}
\item Si $f\in L^1(G)$ y $\phi\in C_c(G)$, entonces $f*\phi\in C_0(G)$.
\item Si $A\in C_r^*(G)$ tiene símbolo $a\in C_0(\widehat G)$ y $\phi,\psi\in C_c(G)$, entonces
\[
 u=A(\phi*\psi)\in E\cap L^2(G),\qquad
 \widehat u=a\widehat\phi\widehat\psi.
\]
Existe una sucesión $u_n\in D$ que converge a $u$ tanto en $E$ como en $L^2$.
\end{enumerate}
\end{lema}
\begin{proof}
Para (i), elija $f_n\in C_c$ con $f_n\to f$ en $L^1$. Como
$\norm{(f-f_n)*\phi}_\infty\le\norm{f-f_n}_1\norm\phi_\infty$,
la función $f*\phi$ es límite uniforme de funciones continuas de soporte compacto.

Para (ii), tome $f_n\in L^1$ tales que $\lambda(f_n)\to A$ en norma de operadores y ponga $u_n=f_n*\phi*\psi$. Cada $u_n$ pertenece a $L^1\cap C_0$ y a $L^2$. Por el lema~\ref{v4:l2-convolucion-cero},
\[
 \norm{u_n-u_m}_\infty
 \le\norm{(\lambda(f_n)-\lambda(f_m))\phi}_2\norm\psi_2
 \le\norm{\lambda(f_n)-\lambda(f_m)}\norm\phi_2\norm\psi_2.
\]
Además, $\widehat u_n=\widehat f_n\widehat\phi\widehat\psi$ converge uniformemente a $a\widehat\phi\widehat\psi$. Por tanto $(u_n)$ converge en $E$ a un elemento con ese símbolo. Por otro lado,
$u_n=\lambda(f_n)(\phi*\psi)\to A(\phi*\psi)$ en $L^2$.
El límite uniforme y el límite en $L^2$ coinciden casi en todas partes: sus integrales contra cualquier $\theta\in C_c$ coinciden por ambas convergencias. Esto prueba todas las afirmaciones. En particular,
\[
 \norm{A(\phi*\psi)}_{0,*}
 \le\max\{\norm\phi_2\norm\psi_2,
                 \norm{\widehat\phi\widehat\psi}_\infty\}\norm A.
\]
\end{proof}

\begin{proposicion}\label{v4:dirac-tres-espacios}
Sea $(\xi_\alpha)$ una red de Dirac no negativa, de integral uno, en $C_c(G)$; esto significa que su masa fuera de cada vecindad de $e$ tiende a cero. Entonces
\[
 u*\xi_\alpha\to u\text{ en }C_0(G),\qquad
 A\lambda(\xi_\alpha)\to A\text{ en }C_r^*(G),\qquad
 v*\xi_\alpha\to v\text{ en }E,
\]
para $u\in C_0(G)$, $A\in C_r^*(G)$ y $v\in E$, respectivamente. Además, $\widehat\xi_\alpha\to1$ uniformemente sobre cada compacto de $\widehat G$.
\end{proposicion}
\begin{proof}
Para $u\in C_0$, la continuidad de traslaciones en norma uniforme da
\[
 \norm{u*\xi_\alpha-u}_\infty
 \le\int_G\xi_\alpha(y)\norm{L_yu-u}_\infty\,d\mu(y)\longrightarrow0.
\]
Se divide la integral entre una vecindad donde el error es menor que $\eps$ y su complemento, donde es a lo sumo $2\norm u_\infty$ y la masa tiende a cero.

Si $A=\lambda(f)$, la convergencia de identidades aproximadas en $L^1$ da
$\norm{A\lambda(\xi_\alpha)-A}\le\norm{f*\xi_\alpha-f}_1\to0$.
Para $A$ general se aproxima por $\lambda(f)$ y se usa $\norm{\lambda(\xi_\alpha)}\le1$; el error adicional es a lo sumo $2\norm{A-\lambda(f)}$. Para $v\in E$, aproximaciones por $D$ muestran que $v*\xi_\alpha\in E$, con símbolo $\widehat v\widehat\xi_\alpha$. Las dos convergencias ya probadas, una para la coordenada función y otra para el operador de símbolo $\widehat v$, implican convergencia en la norma máxima de $E$.

Finalmente, si $K\subset\widehat G$ es compacto, el lema~\ref{v4:dual-equicontinuidad} da una vecindad $U$ de $e$ tal que $|\chi(y)-1|<\eps$ para $\chi\in K$, $y\in U$. Por tanto
\[
 \sup_{\chi\in K}|\widehat\xi_\alpha(\chi)-1|
 \le\eps+2\int_{G\setminus U}\xi_\alpha\,d\mu\longrightarrow\text{a lo sumo }\eps.
\]
Como $\eps$ es arbitrario, se obtiene la convergencia uniforme en $K$. No se está afirmando convergencia uniforme en todo el dual.
\end{proof}

\section{Inversas exactas de funciones de soporte compacto}
El siguiente paso fortalece el último teorema de la lectura 28. El soporte de una función en $G$ y el de una función en $\widehat G$ viven en espacios distintos; las inclusiones de soporte deben indicar cuál de las dos funciones se considera.

\begin{teorema}\label{v4:inversa-compacta}
Para toda $\eta\in C_c(\widehat G)$ existe un único $u\in E$ con $\widehat u=\eta$. Este elemento pertenece a $L^2(G)$ y admite aproximaciones por elementos de $D$ simultáneamente en $E$ y en $L^2$.
\end{teorema}
\begin{proof}
Si $\eta=0$, tome $u=0$. En otro caso, $K=\supp\eta$ es compacto. La equicontinuidad de $K$ permite elegir una función $h\in C_c(G)$, no negativa, simétrica, de integral uno y apoyada suficientemente cerca de $e$, para que
\[
 |\widehat h(\chi)-1|<\tfrac12\qquad(\chi\in K).
\]
En particular, $k=(\widehat h)^2$ no se anula en $K$. Defina
\[
 a(\chi)=\begin{cases}
 \eta(\chi)/k(\chi),&k(\chi)\ne0,\\
 0,&k(\chi)=0.
 \end{cases}
\]
Esta función es continua: en el abierto donde $k\ne0$ es un cociente continuo; un punto donde $k=0$ queda fuera del cerrado $K$ y tiene una vecindad donde $\eta=0$, por lo que $a$ es idénticamente cero allí. Además $\supp a\subset K$, de modo que $a\in C_c(\widehat G)$.

Por Gelfand existe $A\in C_r^*(G)$ de símbolo $a$. El lema~\ref{v4:regularizacion}, con $\phi=\psi=h$, da
$u=A(h*h)\in E\cap L^2$ y
$\widehat u=a(\widehat h)^2=\eta$, junto con la aproximación simultánea. La unicidad se sigue de la inyectividad de la segunda proyección de $E$.
\end{proof}

\begin{corolario}[El enunciado de aproximación, corregido]\label{v4:aproximacion-lectura28}
Si $\eta\in C_c(\widehat G)$ es real y $\eps>0$, existen $f,g\in E\cap L^2$ tales que
\[
 \widehat f,\widehat g\in C_c(\widehat G),\qquad
 \widehat f\le\eta\le\widehat g,\qquad
 \supp\widehat f,\supp\widehat g\subset\supp\eta,\qquad
 \norm{f-g}_{0,*}<\eps.
\]
Toda $\eta\in C_c(\widehat G)$ es, además, límite uniforme de una sucesión de transformadas de elementos de $E$ cuyos soportes espectrales están contenidos en $\supp\eta$.
\end{corolario}
\begin{proof}
Tome la inversa exacta $u$ del teorema anterior y ponga $f=g=u$. Todas las desigualdades son entonces igualdades, y la norma de la diferencia es cero. Para la última afirmación basta la sucesión constante $\widehat u=\eta$.
\end{proof}

La prueba original intenta multiplicar una desigualdad por $\eta$ sin suponer $\eta\ge0$; si $\eta$ es negativa, el orden se invierte. También escribe soportes de $f,g$ donde deben figurar soportes de $\widehat f,\widehat g$. La división por un factor no nulo en $K$ resuelve ambos problemas y demuestra un resultado más fuerte. No impone soporte compacto a la inversa $u$: esa conclusión sería falsa, por ejemplo, para muchas funciones de soporte compacto en el dual de $\R$.

\section{Construcción de la medida de Haar dual}
\index{medida!de Haar dual}
La existencia abstracta de Haar ya asegura una medida en $\widehat G$, pero todavía no fija su escala respecto de $\mu$. La evaluación de la inversa en $e$ seleccionará exactamente la escala necesaria para Plancherel.

\begin{teorema}\label{v4:haar-dual}
Existe una medida de Haar $\nu$ en $\widehat G$ tal que, para cada $\eta\in C_c(\widehat G)$, su inversa $u$ del teorema~\ref{v4:inversa-compacta} satisface
\[
 u(x)=\int_{\widehat G}\eta(\chi)\overline{\chi(x)}\,d\nu(\chi),
 \qquad
 \norm u_2^2=\int_{\widehat G}|\eta(\chi)|^2\,d\nu(\chi).
\]
Esta medida queda determinada por la primera identidad en $x=e$.
\end{teorema}
\begin{proof}
Para $\eta\in C_c(\widehat G)$, sea $u$ su inversa única y defina $J(\eta)=u(e)$. La unicidad hace a $J$ lineal, y el lema~\ref{v4:positividad-identidad} lo hace positivo: $\eta\ge0$ implica $J(\eta)\ge0$. El teorema de Riesz--Markov proporciona una medida regular de Borel $\nu$, finita en compactos, con
\[
 J(\eta)=\int_{\widehat G}\eta\,d\nu.
\]
Para hacer explícita la acotación local exigida en ese teorema, si $K$ es compacto elija $\rho\in C_c$ no negativa con $\rho\ge1$ en $K$. Para $\supp\eta\subset K$, la positividad implica
$|J(\eta)|\le\norm\eta_\infty J(\rho)$: se multiplica $\eta$ por una fase que vuelva real no negativo su valor bajo $J$ y se compara su parte real con $\norm\eta_\infty\rho$. Así no se está suponiendo que $J$ sea acotado para la norma uniforme en todo $C_c$.

Probemos una identidad que, además, asegura que $\nu$ no es cero. Elija $u_n\in D$ que converjan a $u$ en $E$ y en $L^2$, como permite el teorema~\ref{v4:inversa-compacta}. Las funciones
$q_n=u_n*u_n^*$ pertenecen a $D$: son integrables por Young y pertenecen a $C_0$ por el lema de convolución $L^2$. Ese mismo lema da
\[
 q_n\longrightarrow q=u*u^*\quad\hbox{uniformemente},
 \qquad
 \widehat q_n=|\widehat u_n|^2\longrightarrow|\eta|^2
 \quad\hbox{uniformemente}.
\]
Por tanto $q\in E$ y $\widehat q=|\eta|^2$. Al evaluar en $e$,
\[
 J(|\eta|^2)=q(e)=\int_G|u(y)|^2\,d\mu(y)=\norm u_2^2.
\]
Si $\eta\ne0$, entonces $u\ne0$ por la inyectividad del grafo, y su continuidad y el soporte pleno de Haar implican $\norm u_2>0$. Existen funciones no nulas en $C_c(\widehat G)$; en consecuencia $\nu$ no es la medida cero.

Para $\alpha\in\widehat G$, la función $L_\alpha\eta(\chi)=\eta(\alpha^{-1}\chi)$ es la transformada, en el sentido de $E$, de $\overline\alpha\,u$. Por ello
$J(L_\alpha\eta)=(\overline\alpha\,u)(e)=u(e)=J(\eta)$.
La unicidad en Riesz--Markov implica que $\nu$ es invariante bajo traslaciones, es decir, es una medida de Haar.

Finalmente, $L_{x^{-1}}u$ tiene transformada $\chi\mapsto\overline{\chi(x)}\eta(\chi)$, todavía de soporte compacto. Aplicar la definición de $J$ da
\[
 u(x)=(L_{x^{-1}}u)(e)
 =J\bigl(\overline{\chi(x)}\eta(\chi)\bigr)
 =\int_{\widehat G}\eta(\chi)\overline{\chi(x)}\,d\nu(\chi).
\]
La identidad de normas ya se demostró. La unicidad de $\nu$ se sigue de sus integrales contra todas las funciones de $C_c$.
\end{proof}

\section{Plancherel sin presuponer dualidad}
\index{teorema!de Plancherel, en grupos LCA}
Escribimos inicialmente
\[
 U\eta(x)=\int_{\widehat G}\eta(\chi)\overline{\chi(x)}\,d\nu(\chi)
 \qquad(\eta\in C_c(\widehat G)).
\]
El teorema anterior dice que $U$ es una isometría de $C_c(\widehat G)$, con su norma $L^2$, en $L^2(G)$. Queda comprobar que no falta una parte del espacio de llegada.

\begin{teorema}[Plancherel e inversión]\label{v4:plancherel-lca}
La aplicación $U$ se extiende de modo único a un operador unitario
\[
 U:L^2(\widehat G,\nu)\longrightarrow L^2(G,\mu).
\]
Su inversa $\mathcal F$ coincide con la transformada integral de Fourier en $L^1(G)\cap L^2(G)$, casi en todas partes. En particular,
\[
 \norm{\widehat f}_{L^2(\nu)}=\norm f_{L^2(\mu)}
 \qquad(f\in L^1\cap L^2).
\]
Si además $\widehat f\in L^1(\widehat G,\nu)$, entonces
\[
 f(x)=\int_{\widehat G}\widehat f(\chi)\overline{\chi(x)}\,d\nu(\chi)
\]
para casi todo $x$, y para todo $x$ si se toma el representante continuo del miembro derecho.
\end{teorema}
\begin{proof}
Por densidad de $C_c$ en $L^2(\widehat G)$, la isometría se extiende a todo ese espacio. Su imagen es cerrada. Mostraremos que contiene un subconjunto denso de $L^2(G)$.

Las funciones $z=\phi*\psi$, con $\phi,\psi\in C_c(G)$, tienen envolvente lineal densa en $L^2(G)$. En efecto, $C_c$ es denso en $L^2$, y para $\phi\in C_c$ las convoluciones $\phi*\xi_U$ aproximan $\phi$ en $L^2$. Fijemos una de esas $z$. La identidad aproximada da $\lambda(\xi_U)z\to z$ en $L^2$. Fijado un índice $U$, podemos aproximar $\lambda(\xi_U)$ en norma de operadores por operadores $A$ cuyo símbolo $a$ pertenece a $C_c(\widehat G)$: esto se sigue de $C_r^*(G)\simeq C_0(\widehat G)$ y de la densidad uniforme de $C_c$ en $C_0$. Por regularización,
\[
 Az\in E\cap L^2(G),\qquad
 \widehat{Az}=a\widehat\phi\widehat\psi\in C_c(\widehat G).
\]
La unicidad de la inversa compacta implica que $Az=U(a\widehat\phi\widehat\psi)$. Además,
$\norm{Az-\lambda(\xi_U)z}_2\le\norm{A-\lambda(\xi_U)}\norm z_2$.
Escogiendo primero $U$ y luego $A$, aproximamos $z$ arbitrariamente por elementos de la imagen de $U$. La imagen, al ser cerrada y densa, es todo $L^2(G)$.

Sea ahora $f\in L^1\cap L^2$ y $\eta\in C_c(\widehat G)$. Por la fórmula integral de $U$ y Fubini,
\[
 \ip f{U\eta}_{L^2(G)}
 =\int_{\widehat G}\widehat f(\chi)\overline{\eta(\chi)}\,d\nu(\chi).
\]
La integrabilidad absoluta se justifica por $\norm f_1\norm\eta_1<\infty$. Como $U$ es unitario, el miembro izquierdo es $\ip{\mathcal Ff}\eta$. Así las funciones localmente integrables $\mathcal Ff$ y $\widehat f$ tienen las mismas integrales contra $C_c$. La primera pertenece a $L^2$ y la segunda a $C_0$, por Riemann--Lebesgue; ambas están concentradas, salvo conjuntos nulos, en uniones numerables de compactos. La reducción explicada en la proposición~\ref{v4:grafo-fourier-inyectivo} convierte la igualdad local en igualdad casi en todas partes en todo el dual. Esto demuestra la compatibilidad y Plancherel.

Para justificar la fórmula final, observemos primero que si $k\in L^1(\widehat G)\cap L^2(\widehat G)$, existen $k_n\in C_c(\widehat G)$ con convergencia simultánea en $L^1$ y $L^2$. Se trunca primero la función y su soporte usando integrabilidad; en conjuntos de medida finita, la regularidad de Haar aproxima las funciones simples por funciones continuas de soporte compacto en ambas normas. Entonces $Uk_n\to Uk$ en $L^2$, mientras que las integrales inversas convergen uniformemente porque su diferencia está acotada por $\norm{k_n-k}_1$. Ambos límites coinciden casi en todas partes. Aplique esto a $k=\widehat f$. La integral inversa es continua, pues es límite uniforme de las funciones $Uk_n\in C_0(G)$.
\end{proof}

\begin{observacion}[Normalizaciones]
Si $\mu$ se sustituye por $c\mu$, la medida dual compatible es $c^{-1}\nu$. En efecto, la nueva transformada es $c\widehat f$, y
\[
 \int_{\widehat G}|c\widehat f|^2\,d(c^{-1}\nu)
 =c\int_{\widehat G}|\widehat f|^2\,d\nu
 =\int_G|f|^2\,d(c\mu).
\]
La inversión también conserva exactamente sus factores. En $\R^n$, si $d\mu=dx$ y $\chi_\xi(x)=e^{i\xi\cdot x}$, se tiene $d\nu(\xi)=(2\pi)^{-n}d\xi$, por la fórmula de Fourier euclidiana ya probada. La convención unitaria de los primeros capítulos distribuye el factor como $(2\pi)^{-n/2}$ en cada transformación; aquí la transformada directa no tiene factor y la escala se coloca en la medida dual.
\end{observacion}

\section{Una demostración de la dualidad de Pontryagin}
\index{dualidad!de Pontryagin}
La aplicación canónica al doble dual es
\[
 \iota_G:G\longrightarrow\widehat{\widehat G},\qquad
 \iota_G(x)(\chi)=\chi(x).
\]
Cada $\iota_G(x)$ es un carácter continuo del dual, porque evaluar en el compacto $\{x\}$ es continuo para la topología compacto-abierta. Es un homomorfismo en $x$ y es inyectiva por la separación de puntos. La dificultad consiste en probar que todo carácter del dual es una evaluación y que se recupera exactamente la topología de $G$.

\begin{lema}\label{v4:regular-dual-multiplicadores}
Para $k\in L^1(\widehat G,\nu)$ defina
\[
 \check k(x)=\int_{\widehat G}k(\chi)\overline{\chi(x)}\,d\nu(\chi).
\]
Entonces $\check k\in C_0(G)$ y
\[
 U\lambda_{\widehat G}(k)U^{-1}=M_{\check k},
 \qquad (M_bu)(x)=b(x)u(x).
\]
\end{lema}
\begin{proof}
Para $k\in C_c(\widehat G)$ ya sabemos que $\check k=Uk\in C_0(G)$. Aproximar un $k\in L^1$ por funciones de $C_c$ y usar
$\norm{\check k-\check l}_\infty\le\norm{k-l}_1$
prueba la primera afirmación.

Si $k,\eta\in C_c(\widehat G)$, Fubini y el cambio $\chi=\alpha\beta$ dan
\begin{align*}
 U(k*\eta)(x)
 &=\int_{\widehat G}\int_{\widehat G}
       k(\alpha)\eta(\alpha^{-1}\chi)\overline{\chi(x)}
       \,d\nu(\alpha)\,d\nu(\chi)\\
 &=\left(\int_{\widehat G}k(\alpha)\overline{\alpha(x)}\,d\nu(\alpha)\right)
   \left(\int_{\widehat G}\eta(\beta)\overline{\beta(x)}\,d\nu(\beta)\right)\\
 &=\check k(x)U\eta(x).
\end{align*}
La integrabilidad absoluta está acotada por $\norm k_1\norm\eta_1$. La densidad de $C_c$ en $L^2(\widehat G)$ extiende la identidad a toda $\eta\in L^2$ para $k\in C_c$. Luego la aproximación en $L^1$ de $k$, junto con las cotas de la representación regular y de $\check k$, extiende la identidad a todo $k\in L^1$.
\end{proof}

\begin{lema}\label{v4:algebra-inversas-densa}
La clausura uniforme de $\{\check k:k\in L^1(\widehat G)\}$ es $C_0(G)$. Además, $b\mapsto M_b$ es un isomorfismo isométrico de álgebras $C^*$ de $C_0(G)$ sobre su imagen en $\mathcal B(L^2(G))$.
\end{lema}
\begin{proof}
Para la última afirmación, las identidades de producto y adjunto son puntuales, y $\norm{M_b}\le\norm b_\infty$. Si $0<t<\norm b_\infty$, existe un abierto no vacío donde $|b|>t$. Contiene un conjunto medible $K$ con $0<\mu(K)<\infty$: se toma una vecindad relativamente compacta de un punto dentro de ese abierto, cuya medida es positiva. La función $u=1_K/\sqrt{\mu(K)}$ tiene norma uno y $\norm{M_bu}_2\ge t$. Al hacer $t\uparrow\norm b_\infty$ obtenemos la igualdad de normas. El caso $b=0$ es inmediato.

Sea $B$ la clausura uniforme de las inversas. Es una subálgebra de $C_0(G)$ estable bajo conjugación, porque
$\check{k*l}=\check k\,\check l$ y $\check{k^*}=\overline{\check k}$, identidades que se verifican por Fubini e inversión de Haar. No se anula en ningún punto: para $x\in G$ tome $h\in C_c(\widehat G)$, $h\ge0$, $\int h=1$, y $k(\chi)=\chi(x)h(\chi)$. Entonces $\check k(x)=1$.

Además $B$ separa puntos. Si $x\ne y$, existe $\chi_0$ con $\chi_0(x)\ne\chi_0(y)$ por el corolario~\ref{v4:caracteres-separan}. La función $d(\chi)=\overline{\chi(x)}-\overline{\chi(y)}$ es continua y no nula cerca de $\chi_0$. Elija $h\ge0$ de soporte compacto en esa vecindad, no nula, y ponga $k=\overline d\,h$. Entonces
\[
 \check k(x)-\check k(y)=\int_{\widehat G}|d(\chi)|^2h(\chi)\,d\nu(\chi)>0.
\]
La versión localmente compacta de Stone--Weierstrass da $B=C_0(G)$. Para ver de dónde viene esa versión, si $G$ no es compacto se añaden las constantes a $B$ en la compactificación por un punto $G\cup\{\infty\}$. La condición de no anularse separa cada punto de $G$ del punto $\infty$; las otras hipótesis separan los puntos de $G$. Stone--Weierstrass compacto da densidad en $C(G\cup\{\infty\})$ y, al tomar las funciones que valen cero en $\infty$, densidad en $C_0(G)$. Si $G$ es compacto, la misma prueba puede hacerse añadiendo un punto aislado; alternativamente, la versión compacta sin unidad y con la condición de no anulación se deduce de este argumento.
\end{proof}

\begin{teorema}[Dualidad de Pontryagin]\label{v4:pontryagin}
Para todo grupo LCA $G$, la aplicación canónica $\iota_G:G\to\widehat{\widehat G}$ es un isomorfismo de grupos topológicos.
\end{teorema}
\begin{proof}
Los dos lemas anteriores identifican, mediante conjugación por el unitario $U$, el álgebra $C_r^*(\widehat G)$ con el álgebra de todos los multiplicadores $M_b$, $b\in C_0(G)$. En efecto, el álgebra regular es la clausura en norma de operadores de los $\lambda_{\widehat G}(k)$; después de conjugar, esa clausura es la clausura de los $M_{\check k}$. Como la norma de multiplicación coincide con la norma uniforme, ésta es exactamente $\{M_b:b\in C_0(G)\}$. Por tanto
\[
 C_r^*(\widehat G)\simeq C_0(G)
\]
es un isomorfismo isométrico de álgebras $C^*$.

Recordemos con detalle por qué $\Delta(C_0(G))$ es homeomórfico a $G$. En un espacio compacto $X$, todo carácter $\rho$ de $C(X)$ es una evaluación. Si no existiera un cero común de $\ker\rho$, para cada punto habría una función del núcleo no nula allí. Un número finito $f_1,\ldots,f_m$ bastaría para cubrir $X$ por esas regiones de no anulación. Entonces $s=\sum|f_j|^2$ pertenecería al ideal $\ker\rho$ y sería estrictamente positiva. Su inversa $1/s$ sería continua; así $1=s(1/s)$ pertenecería al núcleo, contradicción. Hay por tanto un punto $x$ donde se anulan todas las funciones del núcleo. Como $f-\rho(f)1$ pertenece al núcleo, $\rho(f)=f(x)$ para todo $f$. La separación de puntos da unicidad. La aplicación $x\mapsto\operatorname{ev}_x$ es una biyección continua de compacto a Hausdorff, luego es un homeomorfismo.

Para $G$ no compacto, la unitalización de $C_0(G)$ es $C(G\cup\{\infty\})$; el carácter de evaluación en $\infty$ es el único cuya restricción a $C_0(G)$ es cero. Al quitarlo se obtiene el homeomorfismo $G\simeq\Delta(C_0(G))$. Si $G$ es compacto se aplica directamente el párrafo anterior.

Por otro lado, el teorema~\ref{v4:cstar-fourier}, aplicado al grupo LCA $\widehat G$, identifica $\Delta(C_r^*(\widehat G))$ con $\widehat{\widehat G}$. Bajo estas identificaciones, evaluar el multiplicador $M_{\check k}$ en $x$ produce
\[
 \check k(x)=\int_{\widehat G}k(\chi)\overline{\chi(x)}\,d\nu(\chi)
 =\widehat k\bigl(\iota_G(x)^{-1}\bigr).
\]
La densidad de los operadores $\lambda_{\widehat G}(k)$ determina el carácter completo. En consecuencia, el homeomorfismo entre $G$ y $\widehat{\widehat G}$ que acabamos de construir es $x\mapsto\iota_G(x)^{-1}$. La inversión es un homeomorfismo del grupo doble dual, de modo que $\iota_G$ también es un homeomorfismo. Ya sabíamos que es un homomorfismo; queda demostrada la dualidad.
\end{proof}

El orden de las dependencias importa: primero probamos que los caracteres separan puntos usando el álgebra regular; después construimos Haar dual y Plancherel; sólo entonces usamos el álgebra regular del dual para obtener la sobreyectividad de $\iota_G$. En ningún paso se asumió previamente que un carácter del dual fuera una evaluación.

\section{Aniquiladores, cocientes y extensión de caracteres}
\index{aniquilador!de un subgrupo}\index{caracter@carácter!extension@extensión}
Sea $H$ un subgrupo cerrado de $G$. El grupo $H$ y el cociente $G/H$, con sus topologías naturales, son LCA. Definimos su \emph{aniquilador} por
\[
 H^\perp=\{\chi\in\widehat G:\chi(h)=1\text{ para todo }h\in H\}.
\]
Es un subgrupo cerrado del dual, porque es intersección de los conjuntos cerrados definidos por las evaluaciones $\chi(h)=1$.

\begin{lema}[El dual de un cociente]\label{v4:dual-cociente-encaje}
Si $q:G\to G/H$ es el cociente, la aplicación
\[
 \widehat q:\widehat{G/H}\longrightarrow\widehat G,
 \qquad\gamma\longmapsto\gamma\circ q,
\]
es un isomorfismo de grupos topológicos sobre $H^\perp$.
\end{lema}
\begin{proof}
El carácter $\gamma\circ q$ es trivial en $H$. Recíprocamente, si $\chi\in H^\perp$, la fórmula $\gamma(xH)=\chi(x)$ está bien definida y es un homomorfismo. Es continua por la definición de topología cociente, porque $\gamma\circ q=\chi$ es continua. Esto prueba la biyección algebraica. Si $K\subset G$ es compacto, $q(K)$ es compacto; por tanto la convergencia uniforme en compactos de $G/H$ implica la convergencia uniforme de los caracteres compuestos en $K$, lo cual prueba continuidad de $\widehat q$.

Para la inversa necesitamos levantar compactos. Sea $C\subset G/H$ compacto y elija una vecindad relativamente compacta $V$ de $e$ en $G$. Como $q$ es abierta, los conjuntos $q(xV)$ forman una cubierta abierta de $C$. Un número finito $q(x_1V),\ldots,q(x_mV)$ basta. El conjunto
\[
 L=\bigl(x_1\overline V\cup\cdots\cup x_m\overline V\bigr)\cap q^{-1}(C)
\]
es compacto y satisface $q(L)=C$: todo punto de $C$ tiene un representante en alguno de los $x_jV$, y luego ese representante está en $L$. En consecuencia,
\[
 \sup_{z\in C}|\gamma_1(z)-\gamma_2(z)|
 =\sup_{x\in L}|(\gamma_1\circ q)(x)-(\gamma_2\circ q)(x)|.
\]
La convergencia uniforme en compactos dentro de $H^\perp$ implica la correspondiente convergencia en $\widehat{G/H}$. Esto prueba continuidad de la inversa, sin suponer compacidad de la preimagen completa $q^{-1}(C)$.
\end{proof}

\begin{proposicion}[Doble aniquilador]\label{v4:doble-aniquilador}
Bajo la identificación canónica de $G$ con su doble dual,
\[
 (H^\perp)^\perp=H.
\]
\end{proposicion}
\begin{proof}
Si $h\in H$, cada $\chi\in H^\perp$ satisface $\chi(h)=1$, de modo que $h$ pertenece al doble aniquilador. Si $x\notin H$, entonces $q(x)\ne H$ en $G/H$. Los caracteres de este grupo separan puntos, por el corolario~\ref{v4:caracteres-separan}. Existe $\gamma\in\widehat{G/H}$ con $\gamma(q(x))\ne1$. Su composición $\chi=\gamma\circ q$ pertenece a $H^\perp$ pero $\chi(x)\ne1$. Así $x$ no pertenece al doble aniquilador.
\end{proof}

\begin{teorema}[Extensión y exactitud topológica]\label{v4:dual-exactitud}
La restricción $r:\widehat G\to\widehat H$, $r(\chi)=\chi|_H$, es sobreyectiva y abierta, y su núcleo es $H^\perp$. Por tanto
\[
 \widehat G/H^\perp\simeq\widehat H
\]
como grupos topológicos. En particular, todo carácter continuo de un subgrupo cerrado se extiende a un carácter continuo de $G$.
\end{teorema}
\begin{proof}
La restricción es continua, pues un compacto de $H$ también es compacto en $G$, y su núcleo es $H^\perp$ por definición. Sea $Q=\widehat G/H^\perp$ y sea $q_2:\widehat G\to Q$ la aplicación cociente. La restricción induce un homomorfismo continuo e inyectivo
\[
 j:Q\longrightarrow\widehat H,
 \qquad j(\chi H^\perp)=\chi|_H.
\]
Todavía no suponemos que $j$ sea sobreyectivo. Consideremos su dual
$\widehat j:\widehat{\widehat H}\to\widehat Q$, dado por composición. Es continuo porque la imagen de un compacto por $j$ es compacta.

Por el lema~\ref{v4:dual-cociente-encaje}, la aplicación $\widehat{q_2}:\widehat Q\to\widehat{\widehat G}$ identifica $\widehat Q$ homeomórficamente con $(H^\perp)^\perp$. Por la proposición anterior y Pontryagin, ese subgrupo es $\iota_G(H)$ con la topología heredada. Para $h\in H$ y $\chi\in\widehat G$, un cálculo de evaluaciones da
\[
 \bigl(\widehat{q_2}\,\widehat j(\iota_H(h))\bigr)(\chi)
 =\iota_H(h)\bigl(j(q_2(\chi))\bigr)
 =\chi(h)=\iota_G(h)(\chi).
\]
Pontryagin identifica $H$ homeomórficamente con $\widehat{\widehat H}$ y también con $\iota_G(H)$ como subespacio del doble dual de $G$. La identidad anterior, junto con el encaje de $\widehat{q_2}$, demuestra que $\widehat j$ es un isomorfismo de grupos topológicos.

Falta pasar de esa afirmación a $j$, sin inferirla sólo de una biyección de conjuntos. Para cualquier homomorfismo continuo $t:X\to Y$ de grupos LCA se cumple la identidad de naturalidad
\[
 \iota_Y\circ t=\widehat{\widehat t}\circ\iota_X.
\]
En efecto, ambos miembros evaluados en $\gamma\in\widehat Y$ valen $\gamma(t(x))$. Si $\widehat t$ es un isomorfismo topológico, su dual $\widehat{\widehat t}$ también lo es: el dual de la inversa es la inversa, y ambas aplicaciones son continuas por imágenes compactas. Pontryagin y la identidad de naturalidad implican entonces que $t$ es un isomorfismo topológico. Aplique este argumento a $t=j$. Resulta que $j$ es sobreyectivo y abierto. Finalmente $r=j\circ q_2$ es sobreyectivo y abierto porque lo son sus dos factores. Esto completa la demostración y la dependencia pendiente del ejercicio~\ref{ej:v4-l25-2}.
\end{proof}

\section{Ejercicios de inversión, soporte y aniquiladores}
Las lecturas 27 y 28 proponen varias demostraciones dentro de la teoría; todas quedaron desarrolladas en las secciones anteriores. Los cinco ejercicios siguientes son ampliaciones de esta edición y no se atribuyen a los apuntes originales.

\begin{ejercicio}[Ampliación: Fourier en un grupo finito]\label{ej:v4-ampl-16-1}
Sea $G$ un grupo abeliano finito de orden $N$, con medida de contar. Determine la medida dual de Plancherel y demuestre explícitamente las fórmulas de inversión y Parseval.
\end{ejercicio}


\begin{ejercicio}[Ampliación: el aniquilador de una retícula]\label{ej:v4-ampl-16-2}
Sea $a>0$ y $H=a\Z\subset\R$. Identifique $H^\perp$ usando los caracteres $\chi_\xi(x)=e^{i\xi x}$. Deduzca el dual de $\R/a\Z$ y describa la aplicación de restricción $\widehat\R\to\widehat{a\Z}$, incluyendo su núcleo y su topología.
\end{ejercicio}


\begin{ejercicio}[Ampliación: círculo, enteros y signos]\label{ej:v4-ampl-16-3}
En $G=\T$ use Haar normalizada $dt/(2\pi)$, con $z=e^{it}$, y caracteres $\chi_n(z)=z^n$. Determine la medida dual en $\Z$ y escriba la inversión de un polinomio trigonométrico. Explique por qué los exponentes en la inversión llevan signo negativo con nuestra convención.
\end{ejercicio}


\begin{ejercicio}[Ampliación: positividad espectral y valores negativos]\label{ej:v4-ampl-16-4}
Muestre con un ejemplo que $\widehat f\ge0$ no implica $f\ge0$ punto por punto. Verifique que el ejemplo sí satisface $f(e)\ge0$ y explique por qué no contradice el lema~\ref{v4:positividad-identidad}.
\end{ejercicio}


\begin{ejercicio}[Ampliación: dos soportes que no se confunden]\label{ej:v4-ampl-16-5}
En $G=\R$ con medida de Lebesgue, sea $\eta(\xi)=\max\{1-|\xi|,0\}$. Calcule su inversa para la medida dual $(2\pi)^{-1}d\xi$. Pruebe que esa inversa pertenece a $L^1\cap L^2\cap C_0$ y que su soporte no es compacto, aunque $\supp\eta=[-1,1]$.
\end{ejercicio}


La construcción completa enlaza cuatro niveles: una función integrable, su operador de convolución, su símbolo continuo en el dual y su transformación unitaria en $L^2$. La dualidad y la extensión de caracteres son consecuencias de ese enlace, con la topología de cada paso controlada. El ejercicio de exactitud del capítulo~\ref{cap:v4-14} queda así integrado en la teoría demostrada, y los espacios $C_r^*(G)$ y $C_0^*(G)$ conservan sus papeles distintos a lo largo del argumento.

